
\documentclass[11pt]{amsart}

% Encoding and fonts
\usepackage[T1]{fontenc}
\usepackage[utf8]{inputenc}
\usepackage{lmodern}

% Math packages
\usepackage{amsmath}
\usepackage{amssymb}
\usepackage{amsfonts}
\usepackage{amsthm}
\usepackage{mathtools}

% Layout and utilities
\usepackage{geometry}
\usepackage{enumitem}
\usepackage{booktabs}
\usepackage{longtable}
\usepackage{array}
\usepackage{microtype}
\usepackage{listings}
\usepackage[hidelinks]{hyperref}

\geometry{margin=1in}

\lstdefinestyle{matlabclean}{
  basicstyle=\ttfamily\footnotesize,
  breaklines=true,
  breakatwhitespace=false,
  columns=flexible,
  keepspaces=true,
  showstringspaces=false,
  frame=none,
  numbers=none,
  xleftmargin=0pt,
  tabsize=4,
  aboveskip=0.8em,
  belowskip=0.8em
}

\lstdefinestyle{consoleclean}{
  basicstyle=\ttfamily\scriptsize,
  breaklines=true,
  breakatwhitespace=false,
  columns=flexible,
  keepspaces=true,
  showstringspaces=false,
  frame=none,
  numbers=none,
  xleftmargin=0pt,
  tabsize=2,
  aboveskip=0.8em,
  belowskip=0.8em
}

\numberwithin{equation}{section}
\setcounter{tocdepth}{3}
\setcounter{secnumdepth}{3}

% Theorem-like environments
\theoremstyle{plain}
\newtheorem{theorem}{Theorem}[section]
\newtheorem{proposition}[theorem]{Proposition}
\newtheorem{lemma}[theorem]{Lemma}
\newtheorem{corollary}[theorem]{Corollary}
\newtheorem{conjecture}[theorem]{Conjecture}

\theoremstyle{definition}
\newtheorem{definition}[theorem]{Definition}
\newtheorem{experiment}[theorem]{Experiment}
\newtheorem{heuristic}[theorem]{Heuristic}
\newtheorem{question}[theorem]{Question}
\newtheorem{example}[theorem]{Example}

\theoremstyle{remark}
\newtheorem{remark}[theorem]{Remark}
\newtheorem{warning}[theorem]{Warning}

% Operators and notation
\DeclareMathOperator{\FP}{FP}
\DeclareMathOperator{\Var}{Var}
\DeclareMathOperator{\supp}{supp}
\DeclareMathOperator{\sgn}{sgn}
\DeclareMathOperator{\Li}{Li}
\DeclareMathOperator{\Ei}{Ei}
\newcommand{\RH}{\mathrm{RH}}
\newcommand{\PP}{\mathcal{P\!P}}
\newcommand{\1}{\mathbf{1}}
\newcommand{\dd}{\,\mathrm{d}}
\newcommand{\cH}{\mathcal H}
\newcommand{\cJ}{\mathcal J}
\newcommand{\cD}{\mathcal D}
\newcommand{\cM}{\mathcal M}
\newcommand{\cA}{\mathcal A}
\newcommand{\cP}{\mathcal P}
\newcommand{\cR}{\mathcal R}
\newcommand{\cV}{\mathcal V}

\title[Exponential Zeta Walks and Prime-Power Identities]{
Exponential Sums over the Riemann Zeros versus the
Riemann--von Mangoldt Clock, with Finite Möbius Interpolation
and Prime-Power Spectral Identities
}

\author{Víctor Casaña-Giner}
\date{}

\allowdisplaybreaks
\begin{document}

\maketitle
\tableofcontents

\bigskip

\begin{remark}
\label{rem:status_policy}
This manuscript develops the analytic identities in full theorem--proof
form.  Standard and literature-based material is cited at the point of use.
Original finite identities are proved directly.  Heuristics and conjectures
are explicitly labelled.  Numerical evidence is confined to the final part,
immediately before the conclusions, and is never used as a substitute for an
analytic proof.
\end{remark}


\section{Introduction: the exponential zero sum versus the Riemann--von Mangoldt clock}
\label{sec:introduction}

The point of departure is not the finite interpolation introduced later.
It is the direct comparison between the exponential walk formed by the
actual ordinates of the nontrivial zeros and the walk formed by the smooth
Riemann--von Mangoldt clock.  Let
\begin{equation}
0<\gamma_1\leq\gamma_2\leq\cdots
\label{eq:intro_ordered_ordinates}
\end{equation}
denote the positive ordinates, repeated according to multiplicity, and fix a
frequency \(\omega>0\).  The primary object is
\begin{equation}
S_N(\omega)
=
\sum_{k=1}^{N}e^{i\omega\gamma_k}.
\label{eq:intro_SN_omega}
\end{equation}
The smooth comparison sequence is obtained by inverting the principal
Riemann--von Mangoldt counting function.  Writing
\begin{equation}
N_0(T)
=
\frac{T}{2\pi}\log\!\left(\frac{T}{2\pi e}\right)+\frac78,
\label{eq:intro_N0}
\end{equation}
we define \(\widetilde\gamma_k\) by
\begin{equation}
N_0(\widetilde\gamma_k)=k-\frac12
\label{eq:intro_rvm_inverse_definition}
\end{equation}
and form
\begin{equation}
S_N^{\mathrm{RVM}}(\omega)
=
\sum_{k=1}^{N}e^{i\omega\widetilde\gamma_k}.
\label{eq:intro_SN_RVM}
\end{equation}
The discrepancy
\begin{equation}
\mathcal E_N(\omega)
=
S_N(\omega)-S_N^{\mathrm{RVM}}(\omega)
\label{eq:intro_discrepancy}
\end{equation}
is the object from which the subsequent prime-power investigation arose.
The smooth clock reproduces the main zero density.  For admissible
regularized test functions, the arithmetic side of the explicit formula is
supported on prime powers; this does not by itself assert that every component
of the sharp residual is a prime-power mode.  The two branches with frequencies
\begin{equation}
\omega-\log n
\qquad\text{and}\qquad
\omega+\log n
\label{eq:intro_two_frequencies}
\end{equation}
are therefore not an auxiliary model imposed on the data.  They are the
Fourier displacements of the two prime terms in the Guinand--Weil explicit
formula.

A second construction was found while studying how a finite Möbius section
can be completed by the reciprocal zeta function.  It begins with an affine
interpolation of the terminal prime coefficients.  The smallest numerical
example already contains the complete mechanism.  For \(N=10\),
\begin{align}
A_{10}(s)
&=
1-2^{-s}-3^{-s}-5^{-s}+6^{-s}+10^{-s},
\label{eq:intro_A10}
\\
P_{10}(s)
&=
7^{-s},
\label{eq:intro_P10}
\\
B_{10}
&=
A_{10}(0)=0.
\label{eq:intro_B10}
\end{align}
If the unknown terminal value is denoted by the letter \(m\), the first
interpolation is
\begin{equation}
\mathcal I_{10}(s;m)
=
A_{10}(s)+m7^{-s}.
\label{eq:intro_I10}
\end{equation}
Solving
\begin{equation}
\mathcal I_{10}(s;m)=\zeta(s)^{-1}
\label{eq:intro_I10_completion}
\end{equation}
for the number \(m\) produces
\begin{equation}
\Phi_{10}(s)
=
7^s\left(\zeta(s)^{-1}-A_{10}(s)\right).
\label{eq:intro_Phi10}
\end{equation}
The general symbolic form is
\begin{equation}
\mathcal I_N(s;m)
=
A_N(s)+\frac{m-B_N}{h_N}P_N(s),
\label{eq:intro_general_interpolation}
\end{equation}
and solving \(\mathcal I_N(s;m)=\zeta(s)^{-1}\) yields
\begin{equation}
\Phi_N(s)
=
B_N+h_N\frac{\zeta(s)^{-1}-A_N(s)}{P_N(s)}.
\label{eq:intro_general_Phi}
\end{equation}
This finite interpolation is analytically independent of the phase-walk
experiment, but it uses the same reciprocal-zeta and prime-support
architecture.  It is included because it gives exact finite identities,
not because it is the historical origin of the project.

The principal rigorous results established below are the following.
First, the Riemann--von Mangoldt comparison is written exactly as a
Stieltjes integral and its smooth antiderivative is derived.  Second, the
finite interpolation is solved completely: its value at the origin, its
limit on the positive real axis, its exponential rate, its rounding
criterion, its poles, and its finite parts are proved.  Third, the two
prime-power branches are derived from a regularized explicit formula and
their signed invariant
\begin{equation}
\Delta_\omega(Q)
=
\sum_{n=p^r>Q}
\left(|c_n^-|^2-|c_n^+|^2\right)
\label{eq:intro_Delta}
\end{equation}
is shown to satisfy
\begin{equation}
\Delta_\omega(Q)
\sim
\frac{\omega}{\pi^2\log Q}.
\label{eq:intro_Delta_asymptotic}
\end{equation}
Fourth, the complete-zero walk and the ordinate-only walk are separated by
an atomic positive measure.  This gives several exact routes to the Riemann
Hypothesis.  None is claimed to be completed: each route is stated with the
precise missing uniform or local estimate.

Numerical material is placed only in the final part.  It supports the
analytic structure but is not used in any proof.

\part{The exponential walk and the Riemann--von Mangoldt benchmark}

\section{The zero-counting function and its smooth inverse}
\label{sec:rvm_counting_inverse}

\begin{definition}[Midpoint zero count]
\label{def:midpoint_zero_count}
Let \(N(T)\) be the number of nontrivial zeros \(\rho=\beta+i\gamma\)
with \(0<\gamma\leq T\), counted with multiplicity.  At a zero ordinate
\(T=\tau\), let
\begin{equation}
N^{\#}(\tau)
=
\frac12\left(N(\tau^-)+N(\tau^+)\right).
\label{eq:midpoint_zero_count}
\end{equation}
For a simple ordered ordinate \(\gamma_k\), this equals \(k-1/2\).
\end{definition}

\begin{definition}[Midpoint argument function and exact smooth remainder]
\label{def:rvm_midpoint_argument_remainder}
For \(T\) that is not a zero ordinate, let
\begin{equation}
S(T)
=
\frac1\pi\arg\zeta\!\left(\frac12+iT\right),
\label{eq:rvm_S_function}
\end{equation}
with the standard path from \(2\) to \(2+iT\) and then to
\(1/2+iT\).  At a zero ordinate \(\tau\), define \(S^{\#}(\tau)\) as the
average of the two one-sided limits.  Define
\begin{equation}
R_\vartheta(T)
=
1+\frac{\vartheta(T)}{\pi}-N_0(T),
\label{eq:rvm_exact_smooth_remainder}
\end{equation}
where
\begin{equation}
\vartheta(T)
=
\Im\log\Gamma\!\left(\frac14+\frac{iT}{2}\right)
-
\frac T2\log\pi.
\label{eq:riemann_siegel_theta}
\end{equation}
\end{definition}

\begin{theorem}[Riemann--von Mangoldt formula]
\label{thm:rvm_formula}
For every \(T>0\), with midpoint interpretation at zero ordinates,
\begin{equation}
\boxed{
N^{\#}(T)
=
N_0(T)+S^{\#}(T)+R_\vartheta(T).
}
\label{eq:rvm_formula}
\end{equation}
Moreover,
\begin{equation}
R_\vartheta(T)=O(T^{-1})
\qquad(T\to\infty).
\label{eq:rvm_remainder}
\end{equation}
The exact counting identity and the asymptotic expansion are standard; see
\cite[Chapter~9]{Titchmarsh1986}.
\end{theorem}

\begin{proof}
For \(T\) that is not a zero ordinate, the argument principle applied to
the completed zeta function
\begin{equation}
\xi(s)
=
\frac12s(s-1)\pi^{-s/2}\Gamma\!\left(\frac s2\right)\zeta(s)
\label{eq:rvm_xi_definition}
\end{equation}
gives the exact identity
\begin{equation}
N(T)
=
1+\frac{\vartheta(T)}{\pi}+S(T).
\label{eq:rvm_theta_exact}
\end{equation}
At a zero ordinate, the average of the two one-sided identities gives the
same formula with \(N^{\#}\) and \(S^{\#}\).  Subtracting \(N_0(T)\) and
using the definition \eqref{eq:rvm_exact_smooth_remainder} proves
\eqref{eq:rvm_formula} exactly.

It remains to estimate \(R_\vartheta\).  Stirling's formula in a fixed
sector gives
\begin{equation}
\log\Gamma(z)
=
\left(z-\frac12\right)\log z-z+\frac12\log(2\pi)
+
O(|z|^{-1}).
\label{eq:rvm_stirling}
\end{equation}
With \(z=1/4+iT/2\), taking imaginary parts yields
\begin{equation}
\vartheta(T)
=
\frac T2\log\!\left(\frac{T}{2\pi}\right)
-
\frac T2
-
\frac\pi8
+
O(T^{-1}).
\label{eq:rvm_theta_asymptotic}
\end{equation}
Consequently,
\begin{align}
1+\frac{\vartheta(T)}{\pi}
&=
\frac{T}{2\pi}\log\!\left(\frac{T}{2\pi e}\right)
+\frac78+O(T^{-1})
\label{eq:rvm_substitution_one}
\\
&=
N_0(T)+O(T^{-1}),
\label{eq:rvm_substitution_two}
\end{align}
which is exactly \eqref{eq:rvm_remainder}.
\end{proof}

\begin{theorem}[Principal Lambert--\(W\) inverse of the smooth clock]
\label{thm:rvm_lambert_inverse}
Let \(x>-1/8\).  The equation \(N_0(T)=x\) has a unique solution
\(T>2\pi\), namely
\begin{equation}
\boxed{
T
=
\frac{2\pi(x-7/8)}
{W_0\!\left((x-7/8)/e\right)}.
}
\label{eq:rvm_lambert_inverse}
\end{equation}
Here \(W_0\) is the principal real branch, and the value at \(x=7/8\) is
understood by continuity as \(T=2\pi e\).  In particular, for every
\(k\geq1\),
\begin{equation}
\widetilde\gamma_k
=
\frac{2\pi(k-11/8)}
{W_0\!\left((k-11/8)/e\right)}.
\label{eq:rvm_gamma_tilde}
\end{equation}
\end{theorem}

\begin{proof}
Put
\begin{equation}
y=x-\frac78,
\qquad
u=\frac{T}{2\pi e}.
\label{eq:rvm_inverse_variables}
\end{equation}
The equation \(N_0(T)=x\) is
\begin{equation}
y=e\,u\log u.
\label{eq:rvm_inverse_step_one}
\end{equation}
Writing \(v=\log u\) gives
\begin{equation}
ve^v=\frac ye.
\label{eq:rvm_inverse_step_two}
\end{equation}
Since \(x>-1/8\), one has \(y/e>-1/e\), so the principal branch
\(W_0(y/e)\) is real and greater than \(-1\).  Thus
\begin{equation}
v=W_0(y/e),
\qquad
u=e^v=\frac{y/e}{W_0(y/e)},
\label{eq:rvm_inverse_step_three}
\end{equation}
and hence \eqref{eq:rvm_lambert_inverse}.  The inequality
\(W_0(y/e)>-1\) implies \(u>e^{-1}\), hence \(T>2\pi\).
Conversely, \(N_0'(T)=(2\pi)^{-1}\log(T/2\pi)>0\) on \((2\pi,\infty)\),
so this solution is unique in that interval.  At \(y=0\),
\(W_0(y/e)\sim y/e\), giving the continuous value \(T=2\pi e\).
Finally, \(x=k-1/2>-1/8\) for every \(k\geq1\), and substitution gives
\eqref{eq:rvm_gamma_tilde}.
\end{proof}

\section{Stieltjes representation and the smooth RVM walk}
\label{sec:stieltjes_rvm_walk}

\begin{definition}[Symmetric phase sum]
\label{def:symmetric_phase_sum}
Define
\begin{equation}
S_\omega^{\#}(T)
=
\sum_{0<\gamma<T}m(\gamma)e^{i\omega\gamma}
+
\frac12
\sum_{\gamma=T}m(\gamma)e^{i\omega\gamma}.
\label{eq:symmetric_phase_sum}
\end{equation}
\end{definition}

\begin{proposition}[Atomic representation and integration by parts]
\label{prop:stieltjes_phase_sum}
Let
\begin{equation}
\dd N
=
\sum_{\gamma>0}m(\gamma)\delta_\gamma.
\label{eq:dN_atomic}
\end{equation}
Then
\begin{equation}
\boxed{
S_\omega^{\#}(T)
=
\int_{(0,T)}e^{i\omega t}\,\dd N(t)
+
\frac12e^{i\omega T}\Delta N(T).
}
\label{eq:stieltjes_phase_sum}
\end{equation}
If \(0<T_0<T\) and neither endpoint is a zero ordinate, then
\begin{align}
\int_{(T_0,T]}e^{i\omega t}\,\dd N(t)
={}&
e^{i\omega T}N(T)-e^{i\omega T_0}N(T_0)
\label{eq:stieltjes_integration_by_parts_one}
\\
&-
i\omega\int_{T_0}^{T}e^{i\omega t}N(t)\,\dd t.
\label{eq:stieltjes_integration_by_parts_two}
\end{align}
The midpoint formula at a zero endpoint is obtained by averaging the two
one-sided identities.
\end{proposition}

\begin{proof}
The first identity is the direct evaluation of the atomic measure
\(\dd N\): atoms strictly below \(T\) receive full weight and the atom at
\(T\) receives half weight.  Merely changing the point value of the
counting function would not change its Stieltjes measure, which is why the
explicit endpoint term is required.

For endpoints away from atoms, \(N\) is a right-continuous
bounded-variation function and \(f(t)=e^{i\omega t}\) is continuously
differentiable.  The Riemann--Stieltjes integration-by-parts formula gives
\begin{equation}
\int_{(T_0,T]} f\,\dd N
=
f(T)N(T)-f(T_0)N(T_0)-\int_{T_0}^{T}N(t)f'(t)\,\dd t.
\label{eq:stieltjes_general_ibp}
\end{equation}
Substitution of \(f'(t)=i\omega e^{i\omega t}\) proves the displayed
formula.  Averaging the limits from the two sides supplies half of any
endpoint atom.
\end{proof}

\begin{theorem}[Smooth Riemann--von Mangoldt contribution]
\label{thm:smooth_rvm_contribution}
The Stieltjes integral of the smooth density
\begin{equation}
N_0'(t)
=
\frac1{2\pi}\log\!\left(\frac{t}{2\pi}\right)
\label{eq:N0_derivative}
\end{equation}
has the antiderivative
\begin{equation}
\boxed{
A_{\mathrm{smooth},\omega}(T)
=
\frac{
 e^{i\omega T}\log(T/2\pi)-\Ei(i\omega T)
}{2\pi i\omega}
+C_\omega.
}
\label{eq:smooth_rvm_antiderivative}
\end{equation}
Consequently, for any \(T_0>0\),
\begin{equation}
\int_{T_0}^{T}e^{i\omega t}N_0'(t)\,\dd t
=
A_{\mathrm{smooth},\omega}(T)
-
A_{\mathrm{smooth},\omega}(T_0).
\label{eq:smooth_rvm_definite}
\end{equation}
\end{theorem}

\begin{proof}
Differentiate the numerator in \eqref{eq:smooth_rvm_antiderivative}.
The derivative of the first term is
\begin{equation}
\frac{\dd}{\dd T}
\left[e^{i\omega T}\log(T/2\pi)\right]
=
i\omega e^{i\omega T}\log(T/2\pi)
+
\frac{e^{i\omega T}}{T}.
\label{eq:smooth_derivative_first}
\end{equation}
Since \(\frac{\dd}{\dd z}\Ei(z)=e^z/z\), the chain rule gives
\begin{equation}
\frac{\dd}{\dd T}\Ei(i\omega T)
=
\frac{e^{i\omega T}}{T}.
\label{eq:smooth_derivative_Ei}
\end{equation}
The two \(T^{-1}\) terms cancel.  Division by \(2\pi i\omega\)
leaves exactly \(e^{i\omega T}N_0'(T)\).  The definite identity follows
from the fundamental theorem of calculus on any interval avoiding the
chosen branch cut; a change of branch alters only the constant
\(C_\omega\).
\end{proof}

\begin{corollary}[Exact RVM remainder decomposition]
\label{cor:rvm_remainder_decomposition}
Let \(0<T_0<T\) and suppose first that neither endpoint is a zero
ordinate.  Then
\begin{align}
\int_{(T_0,T]}e^{i\omega t}\,\dd N(t)
={}&
A_{\mathrm{smooth},\omega}(T)
-
A_{\mathrm{smooth},\omega}(T_0)
\label{eq:rvm_remainder_decomposition_one}
\\
&+
\int_{(T_0,T]}e^{i\omega t}\,\dd S^{\#}(t)
+
\int_{T_0}^{T}e^{i\omega t}R_\vartheta'(t)\,\dd t.
\label{eq:rvm_remainder_decomposition_two}
\end{align}
At a zero endpoint, the corresponding symmetric identity is obtained by
averaging the one-sided limits.
\end{corollary}

\begin{proof}
The exact identity \eqref{eq:rvm_formula} is an equality of locally
bounded-variation functions:
\begin{equation}
N^{\#}(t)=N_0(t)+S^{\#}(t)+R_\vartheta(t).
\label{eq:rvm_linear_decomposition_proof}
\end{equation}
Away from endpoint atoms, taking Stieltjes differentials gives
\begin{equation}
\dd N
=
N_0'(t)\,\dd t+\dd S^{\#}(t)+R_\vartheta'(t)\,\dd t.
\label{eq:rvm_differential_decomposition}
\end{equation}
Integrating against \(e^{i\omega t}\) over \((T_0,T]\) and invoking
Theorem~\ref{thm:smooth_rvm_contribution} for the \(N_0'\)-term proves
the formula.  Midpoint endpoint weights follow from the same one-sided
averaging used in Proposition~\ref{prop:stieltjes_phase_sum}.
\end{proof}

\section{Exact comparison of the indexed zero walk with the smooth clock}
\label{sec:indexed_rvm_comparison}

\begin{definition}
\label{def:index_displacement}
For each \(k\), define the RVM displacement
\begin{equation}
\delta_k
=
\gamma_k-\widetilde\gamma_k.
\label{eq:index_displacement}
\end{equation}
\end{definition}

\begin{lemma}[Exact phase perturbation identity]
\label{lem:phase_perturbation}
For real \(x\) and \(h\),
\begin{equation}
e^{i(x+h)}-e^{ix}
=
ih\int_0^1e^{i(x+th)}\,\dd t.
\label{eq:phase_perturbation_integral}
\end{equation}
In particular,
\begin{align}
|e^{i(x+h)}-e^{ix}|
&\leq |h|,
\label{eq:phase_lipschitz}
\\
\left|e^{i(x+h)}-e^{ix}-ihe^{ix}\right|
&\leq\frac{h^2}{2}.
\label{eq:phase_quadratic_remainder}
\end{align}
\end{lemma}

\begin{proof}
Define \(F(t)=e^{i(x+th)}\) on \([0,1]\).  Since
\(F'(t)=ihe^{i(x+th)}\), the fundamental theorem of calculus gives
\begin{equation}
e^{i(x+h)}-e^{ix}
=
F(1)-F(0)
=
ih\int_0^1e^{i(x+th)}\,\dd t.
\label{eq:phase_perturbation_ftc_detailed}
\end{equation}
The integrand has modulus one, so taking absolute values yields the
Lipschitz estimate.  For the quadratic estimate, subtract \(ihe^{ix}\):
\begin{equation}
e^{i(x+h)}-e^{ix}-ihe^{ix}
=
ihe^{ix}\int_0^1(e^{ith}-1)\,\dd t.
\label{eq:phase_quadratic_identity_detailed}
\end{equation}
Apply the already proved Lipschitz bound to \(e^{ith}-1\), obtaining
\(|e^{ith}-1|\leq t|h|\).  Integration of \(t\) from zero to one gives
\(h^2/2\).
\end{proof}

\begin{theorem}[Exact and first-order decompositions of the discrepancy]
\label{thm:indexed_discrepancy_decomposition}
For every \(N\),
\begin{equation}
\mathcal E_N(\omega)
=
i\omega
\sum_{k=1}^{N}
\delta_k
\int_0^1
 e^{i\omega(\widetilde\gamma_k+t\delta_k)}\,\dd t.
\label{eq:exact_indexed_discrepancy}
\end{equation}
Moreover,
\begin{equation}
\mathcal E_N(\omega)
=
i\omega\sum_{k=1}^{N}
\delta_ke^{i\omega\widetilde\gamma_k}
+
\mathcal R_N(\omega),
\label{eq:first_order_indexed_discrepancy}
\end{equation}
where
\begin{equation}
|\mathcal R_N(\omega)|
\leq
\frac{\omega^2}{2}
\sum_{k=1}^{N}\delta_k^2.
\label{eq:indexed_discrepancy_remainder}
\end{equation}
\end{theorem}

\begin{proof}
For each \(k\), apply Lemma~\ref{lem:phase_perturbation} to
\begin{equation}
x=\omega\widetilde\gamma_k,
\qquad
h=\omega(\gamma_k-\widetilde\gamma_k)=\omega\delta_k.
\label{eq:indexed_discrepancy_substitution}
\end{equation}
The resulting identity is
\begin{equation}
e^{i\omega\gamma_k}-e^{i\omega\widetilde\gamma_k}
=
i\omega\delta_k
\int_0^1e^{i\omega(\widetilde\gamma_k+t\delta_k)}\,\dd t.
\label{eq:indexed_discrepancy_termwise}
\end{equation}
Summing from \(1\) to \(N\) gives the exact formula.  For the first-order
form, the quadratic estimate in the lemma gives a remainder \(r_k\) with
\begin{equation}
|r_k|\leq\frac{\omega^2\delta_k^2}{2}.
\label{eq:indexed_term_remainder_bound}
\end{equation}
The triangle inequality then gives the stated bound for the sum of all
remainders.  No restriction on the signs or dependence of the displacements
is needed.
\end{proof}

\begin{proposition}[Displacement in terms of the Riemann--von Mangoldt remainder]
\label{prop:displacement_rvm_remainder}
Assume \(\gamma_k\) is interpreted with midpoint counting, and let
\(\xi_k\) lie between \(\gamma_k\) and \(\widetilde\gamma_k\).  Then
\begin{equation}
\boxed{
\delta_k
=
-
\frac{S^{\#}(\gamma_k)+R(\gamma_k)}
{N_0'(\xi_k)}.
}
\label{eq:displacement_rvm_formula}
\end{equation}
Equivalently,
\begin{equation}
\delta_k
=
-
\frac{2\pi\,[S^{\#}(\gamma_k)+R(\gamma_k)]}
{\log(\xi_k/2\pi)}.
\label{eq:displacement_rvm_explicit}
\end{equation}
\end{proposition}

\begin{proof}
By definition of the smooth inverse and midpoint count,
\begin{equation}
N_0(\widetilde\gamma_k)
=
k-\frac12
=
N^{\#}(\gamma_k).
\label{eq:displacement_count_equalities}
\end{equation}
The Riemann--von Mangoldt formula gives
\begin{equation}
N_0(\widetilde\gamma_k)-N_0(\gamma_k)
=
S^{\#}(\gamma_k)+R(\gamma_k).
\label{eq:displacement_count_difference}
\end{equation}
The mean-value theorem supplies a point \(\xi_k\) between the two
ordinates such that the left-hand side equals
\begin{equation}
N_0'(\xi_k)(\widetilde\gamma_k-\gamma_k)
=
-N_0'(\xi_k)\delta_k.
\label{eq:displacement_mvt}
\end{equation}
Solving for \(\delta_k\) proves the result.
\end{proof}

\begin{remark}
\label{rem:rvm_origin_prime_formula}
Equations \eqref{eq:first_order_indexed_discrepancy} and
\eqref{eq:displacement_rvm_explicit} explain why a comparison with the RVM
clock is sensitive to the argument of zeta.  The explicit formula then
rewrites that argument contribution in terms of prime powers.  This is the
analytic route by which the prime-power branches entered the project.
\end{remark}

\part{Finite Möbius interpolation}


\section{The first affine interpolation: numerical and symbolic forms}
\label{sec:first_affine_interpolation}

\begin{definition}[Terminal affine interpolation]
\label{def:terminal_affine_interpolation}
Let \(Q_N=\{q_1,\ldots,q_{h_N}\}\) be the primes in \((N/2,N]\), and
let \(A_N,P_N,B_N\) be defined in Section~\ref{sec:mobius_definitions}.
For a prescribed terminal sum \(m\in\mathbb C\), define
\begin{equation}
\boxed{
\mathcal I_N(s;m)
=
A_N(s)+\frac{m-B_N}{h_N}P_N(s).
}
\label{eq:terminal_affine_interpolation}
\end{equation}
This is the symmetric affine interpolation in which all terminal-prime
coefficients are assigned their common mean
\begin{equation}
\lambda_N(m)=\frac{m-B_N}{h_N}.
\label{eq:terminal_mean_coefficient}
\end{equation}
\end{definition}

\begin{proposition}[Recovery of the true finite Möbius polynomial]
\label{prop:true_terminal_value}
For the true Mertens value \(m=M(N)\),
\begin{equation}
\mathcal I_N(s;M(N))
=
M_s(N).
\label{eq:true_terminal_interpolation}
\end{equation}
\end{proposition}

\begin{proof}
The set \(Q_N\) contains only primes, so every terminal Möbius coefficient
is exactly \(-1\).  Therefore its total contribution at \(s=0\) is
\(-h_N\), which proves \(M(N)=B_N-h_N\).  Hence the affine mean coefficient
for the true terminal value is
\begin{equation}
\frac{M(N)-B_N}{h_N}
=
\frac{-h_N}{h_N}
=-1.
\label{eq:true_terminal_mean_proof_detailed}
\end{equation}
Substitution assigns coefficient \(-1\) to every term in \(P_N(s)\), so
\begin{equation}
\mathcal I_N(s;M(N))
=
A_N(s)-P_N(s).
\label{eq:true_terminal_interpolation_proof_detailed}
\end{equation}
The right side is exactly the decomposition of \(M_s(N)\).
\end{proof}

\begin{example}[The numerical model \(N=10\)]
\label{ex:N10_affine_interpolation}
The Möbius values up to ten are
\begin{equation}
(1,-1,-1,0,-1,1,-1,0,0,1).
\label{eq:N10_mobius_vector}
\end{equation}
The only terminal prime is \(7\), so
\begin{align}
Q_{10}&=\{7\},
&h_{10}&=1,
&P_{10}(s)&=7^{-s},
\label{eq:N10_terminal_data_detailed}
\\
A_{10}(s)
&=
1-2^{-s}-3^{-s}-5^{-s}+6^{-s}+10^{-s},
&B_{10}&=0.
\label{eq:N10_A_B_detailed}
\end{align}
Thus
\begin{equation}
\mathcal I_{10}(s;m)
=
A_{10}(s)+m7^{-s}.
\label{eq:N10_interpolation_detailed}
\end{equation}
The actual terminal value is \(M(10)=-1\), which restores the coefficient
\(-1\) of \(7^{-s}\).
\end{example}

\begin{theorem}[Completion equation]
\label{thm:completion_equation}
At every point where \(P_N(s)\neq0\), the unique value of \(m\) satisfying
\begin{equation}
\mathcal I_N(s;m)=\zeta(s)^{-1}
\label{eq:completion_equation}
\end{equation}
is
\begin{equation}
\boxed{
 m=\Phi_N(s)
 =
 B_N+h_N\frac{\zeta(s)^{-1}-A_N(s)}{P_N(s)}.
}
\label{eq:completion_solution}
\end{equation}
\end{theorem}

\begin{proof}
Starting from the affine interpolation, the completion equation is
\begin{equation}
A_N(s)+\frac{m-B_N}{h_N}P_N(s)
=
\zeta(s)^{-1}.
\label{eq:completion_equation_expanded}
\end{equation}
Because \(h_N\neq0\) and \(P_N(s)\neq0\), this is a nondegenerate linear
equation in \(m\).  Subtract \(A_N(s)\), multiply by \(h_N/P_N(s)\), and
add \(B_N\) to obtain the displayed solution.  If two values \(m_1,m_2\)
solved the equation, subtracting their equations would give
\begin{equation}
\frac{m_1-m_2}{h_N}P_N(s)=0,
\label{eq:completion_uniqueness}
\end{equation}
so \(m_1=m_2\).  This proves both existence and uniqueness.
\end{proof}

\begin{remark}
\label{rem:Phi_meaning}
The function \(\Phi_N(s)\) is therefore not introduced by an arbitrary
quotient.  It is the terminal value selected by requiring the affine finite
interpolation to coincide with the completed reciprocal-zeta Dirichlet
series at the spectral parameter \(s\).
\end{remark}

\section{Terminal half-interval decomposition}
\label{sec:mobius_definitions}

\begin{definition}
\label{def:terminal_prime_set}
For an integer \(N\geq 2\), define
\begin{equation}
Q_N
=
\left\{
q\ \text{prime}:\frac{N}{2}<q\leq N
\right\},
\label{eq:terminal_prime_set}
\end{equation}
and put
\begin{equation}
h_N=|Q_N|,
\qquad
q_N=\min Q_N.
\label{eq:hN_qN}
\end{equation}
Bertrand's postulate guarantees \(h_N\geq 1\).
\end{definition}

\begin{definition}
\label{def:mobius_polynomials}
Let
\begin{align}
M_s(N)
&=
\sum_{n\leq N}\frac{\mu(n)}{n^s},
\label{eq:MsN_definition}
\\
P_N(s)
&=
\sum_{q\in Q_N}q^{-s},
\label{eq:PN_definition}
\\
A_N(s)
&=
\sum_{\substack{n\leq N\\n\notin Q_N}}
\frac{\mu(n)}{n^s},
\label{eq:AN_definition}
\\
B_N
&=
A_N(0)
=
\sum_{\substack{n\leq N\\n\notin Q_N}}\mu(n).
\label{eq:BN_definition}
\end{align}
Since every \(q\in Q_N\) is prime and \(\mu(q)=-1\),
\begin{equation}
M_s(N)=A_N(s)-P_N(s)
\label{eq:Ms_A_minus_P}
\end{equation}
and
\begin{equation}
M(N)=B_N-h_N.
\label{eq:M_B_minus_h}
\end{equation}
\end{definition}

\begin{definition}
\label{def:tail_RN}
For \(\Re s>1\), define the exact tail
\begin{equation}
R_N(s)
=
\frac{1}{\zeta(s)}-M_s(N)
=
\sum_{n>N}\frac{\mu(n)}{n^s}.
\label{eq:RN_definition}
\end{equation}
The first equality supplies a meromorphic continuation of \(R_N\).
\end{definition}

\begin{definition}
\label{def:PhiN}
Define
\begin{equation}
\Phi_N(s)
=
B_N+
h_N\frac{\zeta(s)^{-1}-A_N(s)}{P_N(s)}
\label{eq:PhiN_definition}
\end{equation}
at points at which the quotient is defined, and meromorphically elsewhere.
\end{definition}

\section{Canonical identity and endpoint values}
\label{sec:canonical_identity}

\begin{theorem}[Canonical orientation]
\label{thm:canonical_orientation}
For every \(N\geq2\),
\begin{equation}
\boxed{
\Phi_N(s)
=
B_N+
h_N\frac{\zeta(s)^{-1}-A_N(s)}{P_N(s)}
=
M(N)+h_N\frac{R_N(s)}{P_N(s)}.
}
\label{eq:canonical_orientation}
\end{equation}
\end{theorem}

\begin{proof}
The functions \(A_N(s)\) and \(P_N(s)\) are finite Dirichlet polynomials
and hence entire.  In the half-plane \(\Re s>1\), absolute convergence of
the reciprocal-zeta Dirichlet series gives
\begin{equation}
\zeta(s)^{-1}
=
M_s(N)+R_N(s).
\label{eq:canonical_proof_reciprocal_split}
\end{equation}
The terminal decomposition \(M_s(N)=A_N(s)-P_N(s)\) therefore yields
\begin{align}
\zeta(s)^{-1}-A_N(s)
&=
M_s(N)+R_N(s)-A_N(s)
\label{eq:canonical_proof_step_one}
\\
&=
R_N(s)-P_N(s).
\label{eq:canonical_proof_step_two}
\end{align}
Divide by \(P_N(s)\) at a point where it is nonzero, multiply by \(h_N\),
and add \(B_N\):
\begin{align}
\Phi_N(s)
&=
B_N+h_N\left(\frac{R_N(s)}{P_N(s)}-1\right)
\label{eq:canonical_proof_step_three}
\\
&=
B_N-h_N+h_N\frac{R_N(s)}{P_N(s)}.
\label{eq:canonical_proof_step_four}
\end{align}
Because the terminal primes all have Möbius value \(-1\),
\(M(N)=B_N-h_N\).  This proves the identity on \(\Re s>1\).  Both sides
are meromorphic continuations of the same function, so the identity extends
to every connected component of their common domain and hence
meromorphically to the plane.
\end{proof}

\begin{corollary}[Value at the origin]
\label{cor:Phi_origin}
The meromorphic interpolation satisfies
\begin{equation}
\boxed{\Phi_N(0)=-2.}
\label{eq:Phi_origin}
\end{equation}
\end{corollary}

\begin{proof}
The point \(s=0\) is regular for every finite Dirichlet polynomial in the
formula.  The classical special value \(\zeta(0)=-1/2\) gives
\begin{equation}
\zeta(0)^{-1}=-2.
\label{eq:origin_inverse_zeta}
\end{equation}
Moreover, by definition,
\begin{equation}
A_N(0)=B_N,
\qquad
P_N(0)=\sum_{q\in Q_N}1=h_N.
\label{eq:origin_finite_values}
\end{equation}
Since \(h_N\geq1\), the denominator does not vanish.  Direct substitution
in the completion solution gives
\begin{align}
\Phi_N(0)
&=
B_N+h_N\frac{-2-B_N}{h_N}
\label{eq:origin_substitution_detailed}
\\
&=
B_N-2-B_N
=-2.
\label{eq:origin_cancellation_detailed}
\end{align}
No limiting argument is involved.
\end{proof}

\begin{corollary}[Large positive real limit]
\label{cor:Phi_infinity}
Along the positive real axis,
\begin{equation}
\boxed{
\lim_{\sigma\to+\infty}\Phi_N(\sigma)=M(N).
}
\label{eq:Phi_infinity}
\end{equation}
\end{corollary}

\begin{proof}
For \(\sigma>1\), the tail satisfies
\begin{equation}
|R_N(\sigma)|
\leq
\sum_{n>N}n^{-\sigma}
\leq
\int_N^\infty x^{-\sigma}\,\dd x
=
\frac{N^{1-\sigma}}{\sigma-1}.
\label{eq:infinity_tail_bound}
\end{equation}
The term \(q_N^{-\sigma}\) occurs in \(P_N(\sigma)\), and all terms are
positive, hence
\begin{equation}
P_N(\sigma)\geq q_N^{-\sigma}.
\label{eq:infinity_denominator_bound}
\end{equation}
Consequently,
\begin{equation}
\left|h_N\frac{R_N(\sigma)}{P_N(\sigma)}\right|
\leq
\frac{h_NN}{\sigma-1}
\left(\frac{q_N}{N}\right)^\sigma.
\label{eq:infinity_ratio_bound}
\end{equation}
For \(N\geq3\), Bertrand's theorem gives a prime in \((N/2,N)\), so
\(q_N<N\), and the exponential factor tends to zero.  For \(N=2\),
\(q_N/N=1\), but the same bound is \(2/(\sigma-1)\), which also tends to
zero.  Hence the second expression in
Theorem~\ref{thm:canonical_orientation} tends to \(M(N)\) for every
\(N\geq2\).
\end{proof}

\begin{example}[The \(N=10\) orientation check]
\label{exp:N10_orientation}
For \(N=10\),
\begin{equation}
Q_{10}=\{7\},
\qquad
h_{10}=1,
\qquad
M(10)=-1,
\qquad
B_{10}=0.
\label{eq:N10_data}
\end{equation}
The inverted quotient \(P_N/R_N\) would give \(-1/2\) at \(s=0\).
The quotient in \eqref{eq:canonical_orientation} gives
\begin{equation}
0+\frac{-2}{1}=-2.
\label{eq:N10_correct_value}
\end{equation}
This finite check confirms the orientation of the quotient in Theorem~\ref{thm:canonical_orientation}.
\end{example}

\section{Exponential approach, rounding, and recovery of the first tail term}
\label{sec:asymptotic_recovery}

\begin{definition}
\label{def:aN}
Let
\begin{equation}
a_N
=
\min\{n>N:\mu(n)\neq0\}.
\label{eq:aN_definition}
\end{equation}
\end{definition}

\begin{theorem}[Leading exponential term]
\label{thm:leading_exponential}
As \(\sigma\to+\infty\),
\begin{equation}
\boxed{
\Phi_N(\sigma)-M(N)
=
h_N\mu(a_N)
\left(\frac{q_N}{a_N}\right)^\sigma
\left(1+o(1)\right).
}
\label{eq:leading_exponential}
\end{equation}
\end{theorem}

\begin{proof}
Factor the least terminal prime from the denominator:
\begin{equation}
P_N(\sigma)
=
q_N^{-\sigma}
\left[
1+
\sum_{\substack{q\in Q_N\\q>q_N}}
\left(\frac{q_N}{q}\right)^\sigma
\right].
\label{eq:leading_P_factorization}
\end{equation}
Every ratio in the bracket is strictly smaller than one, and there are only
finitely many of them.  Thus the bracket is \(1+o(1)\).

For the numerator, factor the first nonzero Möbius term after \(N\):
\begin{equation}
R_N(\sigma)
=
\mu(a_N)a_N^{-\sigma}
\left[
1+
\sum_{n>a_N}
\frac{\mu(n)}{\mu(a_N)}
\left(\frac{a_N}{n}\right)^\sigma
\right].
\label{eq:leading_R_factorization}
\end{equation}
For \(\sigma\geq2\), the absolute value of the summand is bounded by
\((a_N/n)^2\), whose sum over \(n>a_N\) is finite.  Each fixed summand
tends to zero as \(\sigma\to\infty\).  Dominated convergence therefore
shows that the bracket is \(1+o(1)\).  Division gives
\begin{equation}
\frac{R_N(\sigma)}{P_N(\sigma)}
=
\mu(a_N)
\left(\frac{q_N}{a_N}\right)^\sigma
(1+o(1)).
\label{eq:leading_ratio_detailed}
\end{equation}
Multiply by \(h_N\) and use the canonical identity.
\end{proof}

\begin{remark}
\label{rem:orientation_growth}
The factor is \((q_N/a_N)^\sigma\), not
\((a_N/q_N)^\sigma\).  Since \(q_N\leq N<a_N\), the correct factor
decays.  The inverted factor would contradict
Corollary~\ref{cor:Phi_infinity}.
\end{remark}

\begin{theorem}[Uniform real-axis error bound]
\label{thm:rounding_bound}
For every real \(\sigma>1\),
\begin{equation}
\boxed{
|\Phi_N(\sigma)-M(N)|
\leq
\frac{h_NN}{\sigma-1}
\left(\frac{q_N}{N}\right)^\sigma.
}
\label{eq:rounding_bound}
\end{equation}
\end{theorem}

\begin{proof}
The reciprocal-zeta tail is absolutely convergent for \(\sigma>1\), so
\begin{align}
|R_N(\sigma)|
&\leq
\sum_{n>N}|\mu(n)|n^{-\sigma}
\label{eq:rounding_abs_mu}
\\
&\leq
\sum_{n>N}n^{-\sigma}
\label{eq:rounding_drop_mu}
\\
&\leq
\int_N^\infty x^{-\sigma}\,\dd x
=
\frac{N^{1-\sigma}}{\sigma-1}.
\label{eq:rounding_integral_detailed}
\end{align}
The denominator is a sum of positive real terms.  Retaining only its
largest term gives
\begin{equation}
P_N(\sigma)
=
\sum_{q\in Q_N}q^{-\sigma}
\geq
q_N^{-\sigma}.
\label{eq:rounding_P_positive}
\end{equation}
The canonical identity now gives
\begin{align}
|\Phi_N(\sigma)-M(N)|
&=
h_N\frac{|R_N(\sigma)|}{P_N(\sigma)}
\label{eq:rounding_identity_use}
\\
&\leq
\frac{h_NN^{1-\sigma}q_N^\sigma}{\sigma-1}
\label{eq:rounding_combine}
\\
&=
\frac{h_NN}{\sigma-1}
\left(\frac{q_N}{N}\right)^\sigma.
\label{eq:rounding_final_detailed}
\end{align}
\end{proof}

\begin{corollary}[Finite rounding criterion]
\label{cor:rounding_criterion}
If
\begin{equation}
\boxed{
\frac{h_NN}{\sigma-1}
\left(\frac{q_N}{N}\right)^\sigma
<
\frac12,
}
\label{eq:rounding_criterion}
\end{equation}
then \(M(N)\) is the unique nearest integer to the real number
\(\Phi_N(\sigma)\).
\end{corollary}

\begin{proof}
Theorem~\ref{thm:rounding_bound} and
\eqref{eq:rounding_criterion} give
\begin{equation}
|\Phi_N(\sigma)-M(N)|<\frac12.
\label{eq:rounding_criterion_proof_ball}
\end{equation}
For any integer \(k\neq M(N)\), the reverse triangle inequality yields
\begin{align}
|\Phi_N(\sigma)-k|
&\geq
|k-M(N)|-|\Phi_N(\sigma)-M(N)|
\label{eq:rounding_reverse_triangle_one}
\\
&>
1-\frac12
=
\frac12.
\label{eq:rounding_reverse_triangle_two}
\end{align}
Thus \(M(N)\) is strictly closer to \(\Phi_N(\sigma)\) than every other
integer.
\end{proof}

\begin{definition}
\label{def:EN}
Define
\begin{equation}
E_N(s)
=
\frac{\Phi_N(s)-M(N)}{h_N}
=
\frac{R_N(s)}{P_N(s)}.
\label{eq:EN_definition}
\end{equation}
\end{definition}

\begin{corollary}[Recovery of \(a_N\)]
\label{cor:aN_recovery}
The first nonzero Möbius term after \(N\) is recovered by
\begin{equation}
\boxed{
a_N
=
q_N
\lim_{\sigma\to+\infty}
|E_N(\sigma)|^{-1/\sigma}.
}
\label{eq:aN_recovery}
\end{equation}
\end{corollary}

\begin{proof}
By Theorem~\ref{thm:leading_exponential} and
Definition~\ref{def:EN},
\begin{equation}
E_N(\sigma)
=
\mu(a_N)
\left(\frac{q_N}{a_N}\right)^\sigma
(1+o(1)).
\label{eq:aN_recovery_asymptotic}
\end{equation}
Since \(|\mu(a_N)|=1\), taking absolute values and then
\(-1/\sigma\)-th powers gives
\begin{equation}
|E_N(\sigma)|^{-1/\sigma}
=
\frac{a_N}{q_N}
\,|1+o(1)|^{-1/\sigma}
\longrightarrow
\frac{a_N}{q_N}.
\label{eq:aN_recovery_root_limit}
\end{equation}
Multiplication by \(q_N\) proves \eqref{eq:aN_recovery}.
\end{proof}

\section{Feasible terminal values and non-distinction in \(\Re s>1\)}
\label{sec:feasible_values}

\begin{definition}
\label{def:feasible_interpolation}
For an integer parameter \(m\), define
\begin{equation}
F_N(s;m)
=
M_s(N)
+
P_N(s)
\left(
1+\frac{m-B_N}{h_N}
\right).
\label{eq:feasible_interpolation}
\end{equation}
A value \(m\) is called terminally feasible when
\begin{equation}
-1\leq\frac{m-B_N}{h_N}\leq1.
\label{eq:feasible_interval}
\end{equation}
\end{definition}

\begin{proposition}[Exact substitution identity]
\label{prop:feasible_exact_identity}
The expression \eqref{eq:feasible_interpolation} is the exact result of
substituting
\begin{equation}
A_N(s)=M_s(N)+P_N(s)
\label{eq:A_M_plus_P}
\end{equation}
into the terminal-value interpolation.
\end{proposition}

\begin{proof}
From \(M_s(N)=A_N(s)-P_N(s)\) one has
\(A_N(s)=M_s(N)+P_N(s)\).  Substitution in
\eqref{eq:terminal_affine_interpolation} gives
\begin{align}
\mathcal I_N(s;m)
&=
M_s(N)+P_N(s)
+\frac{m-B_N}{h_N}P_N(s)
\label{eq:feasible_substitution_step}
\\
&=
F_N(s;m),
\label{eq:feasible_substitution_result}
\end{align}
which is the claimed identity.
\end{proof}

\begin{theorem}[The half-plane limit does not select \(M(N)\)]
\label{thm:halfplane_nondistinction}
Fix \(s\) with \(\Re s>1\).  For every sequence of feasible integers
\(m(N)\),
\begin{equation}
\boxed{
F_N(s;m(N))
\longrightarrow
\frac1{\zeta(s)}
\qquad
(N\to\infty).
}
\label{eq:halfplane_nondistinction}
\end{equation}
Thus convergence in the absolutely convergent half-plane does not
distinguish \(m(N)=M(N)\) from other feasible terminal values.
\end{theorem}

\begin{proof}
Fix \(s\) with \(\sigma=\Re s>1\).  Absolute convergence of
\(\sum\mu(n)n^{-s}\) gives
\begin{equation}
M_s(N)\longrightarrow\zeta(s)^{-1}.
\label{eq:halfplane_Ms_limit_detailed}
\end{equation}
For the terminal prime polynomial,
\begin{align}
|P_N(s)|
&\leq
\sum_{N/2<q\leq N}q^{-\sigma}
\label{eq:halfplane_P_bound_one}
\\
&\leq
\sum_{N/2<n\leq N}n^{-\sigma}
\label{eq:halfplane_P_bound_two}
\\
&\leq
\int_{N/2-1}^{N}x^{-\sigma}\,\dd x
=
O_\sigma(N^{1-\sigma}).
\label{eq:halfplane_P_bound_three}
\end{align}
Feasibility means \(|(m(N)-B_N)/h_N|\leq1\), so the multiplier
\(1+(m(N)-B_N)/h_N\) has absolute value at most two.  The entire terminal
correction is therefore \(O_\sigma(N^{1-\sigma})\), uniformly over all
feasible choices.  Adding it to \(M_s(N)\) proves the common limit.
\end{proof}

\begin{remark}
\label{rem:arithmetic_not_complex_selection}
The selection of the true value \(M(N)\) is arithmetical and
combinatorial.  It is not produced by convergence to \(1/\zeta(s)\) in
\(\Re s>1\), because the entire feasible family has the same limit.
\end{remark}

\section{Poles and finite parts at zeta zeros}
\label{sec:finite_parts}

\begin{definition}
\label{def:finite_part}
If a meromorphic function has a Laurent expansion
\begin{equation}
f(s)
=
\frac{a_{-1}}{s-\rho}+a_0+O(s-\rho),
\label{eq:finite_part_laurent}
\end{equation}
its finite part at \(\rho\) is
\begin{equation}
\FP_{s=\rho}f(s)=a_0.
\label{eq:finite_part_definition}
\end{equation}
\end{definition}

\begin{proposition}[Finite part of the reciprocal zeta function]
\label{prop:finite_part_inverse_zeta}
Let \(\rho\) be a simple zero of \(\zeta\).  Then
\begin{equation}
\frac1{\zeta(s)}
=
\frac{1}{\zeta'(\rho)(s-\rho)}
-
\frac{\zeta''(\rho)}{2\zeta'(\rho)^2}
+
O(s-\rho),
\label{eq:inverse_zeta_laurent}
\end{equation}
and hence
\begin{equation}
\FP_{s=\rho}\frac1{\zeta(s)}
=
-\frac{\zeta''(\rho)}{2\zeta'(\rho)^2}.
\label{eq:inverse_zeta_finite_part}
\end{equation}
\end{proposition}

\begin{proof}
Write \(z=s-\rho\).  Since the zero is simple,
\begin{equation}
\zeta(\rho+z)
=
\zeta'(\rho)z
+
\frac12\zeta''(\rho)z^2
+
O(z^3).
\label{eq:inverse_zeta_taylor_detailed}
\end{equation}
Factor the linear term:
\begin{equation}
\zeta(\rho+z)
=
\zeta'(\rho)z
\left[
1+
\frac{\zeta''(\rho)}{2\zeta'(\rho)}z
+O(z^2)
\right].
\label{eq:inverse_zeta_factor_detailed}
\end{equation}
Using \((1+w)^{-1}=1-w+O(w^2)\),
\begin{align}
\frac1{\zeta(\rho+z)}
&=
\frac1{\zeta'(\rho)z}
\left[
1-
\frac{\zeta''(\rho)}{2\zeta'(\rho)}z
+O(z^2)
\right]
\label{eq:inverse_zeta_invert_detailed}
\\
&=
\frac1{\zeta'(\rho)z}
-
\frac{\zeta''(\rho)}{2\zeta'(\rho)^2}
+O(z).
\label{eq:inverse_zeta_laurent_detailed}
\end{align}
The constant Laurent coefficient is the finite part.
\end{proof}

\begin{proposition}[Finite part of \(\Phi_N\)]
\label{prop:finite_part_Phi}
Assume that \(\rho\) is a simple zero of \(\zeta\) and
\(P_N(\rho)\neq0\).  Then
\begin{align}
\FP_{s=\rho}\Phi_N(s)
={}&
B_N
+
h_N
\frac{
-\dfrac{\zeta''(\rho)}{2\zeta'(\rho)^2}
-A_N(\rho)
}{
P_N(\rho)
}
\label{eq:Phi_finite_part_first}
\\
&-
h_N
\frac{P_N'(\rho)}
{\zeta'(\rho)P_N(\rho)^2}.
\label{eq:Phi_finite_part_second}
\end{align}
\end{proposition}

\begin{proof}
Set \(z=s-\rho\).  Proposition~\ref{prop:finite_part_inverse_zeta} and
the Taylor expansion of the entire function \(A_N\) give
\begin{equation}
\zeta(s)^{-1}-A_N(s)
=
\frac{a_{-1}}{z}+a_0+O(z),
\label{eq:Phi_fp_numerator_detailed}
\end{equation}
where
\begin{equation}
a_{-1}=\frac1{\zeta'(\rho)},
\qquad
a_0=-\frac{\zeta''(\rho)}{2\zeta'(\rho)^2}-A_N(\rho).
\label{eq:Phi_fp_a_coefficients}
\end{equation}
Likewise,
\begin{equation}
P_N(s)
=
p_0+p_1z+O(z^2),
\qquad
p_0=P_N(\rho),\quad p_1=P_N'(\rho).
\label{eq:Phi_fp_denominator_detailed}
\end{equation}
Because \(p_0\neq0\),
\begin{equation}
\frac1{P_N(s)}
=
\frac1{p_0}-\frac{p_1}{p_0^2}z+O(z^2).
\label{eq:Phi_fp_inverse_denominator}
\end{equation}
Multiplication yields
\begin{align}
\frac{\zeta(s)^{-1}-A_N(s)}{P_N(s)}
={}&
\frac{a_{-1}}{p_0z}
+
\frac{a_0}{p_0}
-
\frac{a_{-1}p_1}{p_0^2}
+O(z).
\label{eq:Phi_fp_quotient_detailed}
\end{align}
Multiplying by \(h_N\) and adding \(B_N\) gives the asserted finite part.
\end{proof}

\begin{definition}
\label{def:finite_part_scale}
For a simple zero \(\rho\), define the normalized finite-part observable
\begin{equation}
\mathcal F_N(\rho)
=
N^{-\rho}\FP_{s=\rho}\Phi_N(s).
\label{eq:finite_part_scale}
\end{equation}
\end{definition}





\part{Prime-power phase walks over zeta ordinates}

\section{Phase walks, endpoint conventions, and a universal half}
\label{sec:phase_walks}

\begin{definition}
\label{def:phase_walk_q}
Fix \(q>1\) and write
\begin{equation}
\omega=\log q.
\label{eq:omega_definition}
\end{equation}
For the positive ordinates, counted with multiplicity, define
\begin{equation}
S_q(T)
=
\sum_{0<\gamma\leq T}m(\gamma)e^{i\omega\gamma}.
\label{eq:Sq_definition}
\end{equation}
At an ordinate \(\gamma_k\), the right, left, and midpoint values are
\begin{align}
S_k^{\mathrm{right}}
&=
\sum_{j\leq k}b_j,
\label{eq:right_sum}
\\
S_k^{\mathrm{left}}
&=
\sum_{j<k}b_j,
\label{eq:left_sum}
\\
S_k^{\mathrm{mid}}
&=
\sum_{j<k}b_j+\frac12b_k,
\qquad
b_k=e^{i\omega\gamma_k}.
\label{eq:mid_sum}
\end{align}
\end{definition}

\begin{remark}
\label{rem:indexwise_vs_atomic_midpoint}
The quantities \(S_k^{\mathrm{right}},S_k^{\mathrm{left}}\), and
\(S_k^{\mathrm{mid}}\) are indexwise: an ordinate of multiplicity \(m\)
appears \(m\) times in the ordered list.  The symmetric midpoint of the
whole atom at a distinct ordinate \(\tau\) instead contains
\(\frac12m(\tau)e^{i\omega\tau}\).
\end{remark}

\begin{proposition}[Endpoint conversion]
\label{prop:endpoint_conversion}
At every index \(k\),
\begin{align}
\overline{b_k}S_k^{\mathrm{right}}
&=
\overline{b_k}S_k^{\mathrm{mid}}+\frac12,
\label{eq:right_mid_conversion}
\\
\overline{b_k}S_k^{\mathrm{left}}
&=
\overline{b_k}S_k^{\mathrm{mid}}-\frac12.
\label{eq:left_mid_conversion}
\end{align}
Thus changing the endpoint convention translates the carrier coefficient
by \(\pm1/2\) and does not change a centered variance.
\end{proposition}

\begin{proof}
The right sum differs from the indexwise midpoint sum by
\(\frac12b_k\), while the left sum differs by \(-\frac12b_k\).
Multiplication by \(\overline b_k\) and \(|b_k|=1\) gives the two
identities.  Translating all sampled values by a fixed constant does not
alter their centered variance.
\end{proof}

\begin{theorem}[Exact triangular half identity]
\label{thm:triangular_half}
Let \(|b_k|=1\) and
\begin{equation}
S_k=\sum_{j\leq k}b_j.
\label{eq:triangular_partial_sums}
\end{equation}
Then
\begin{equation}
\boxed{
\Re\left(
\frac1N\sum_{k=1}^N\overline{b_k}S_k
\right)
=
\frac12+\frac{|S_N|^2}{2N}.
}
\label{eq:triangular_half_identity}
\end{equation}
\end{theorem}

\begin{proof}
Let \(b=(b_1,\ldots,b_N)^{\mathsf T}\), and let \(L\) be the lower
triangular matrix with \(L_{kj}=1\) for \(j\leq k\).  Then the vector of
partial sums is \(Lb\), and
\begin{equation}
\sum_{k=1}^{N}\overline{b_k}S_k
=
b^\ast Lb.
\label{eq:triangular_form_detailed}
\end{equation}
Let \(J\) denote the all-ones matrix.  Entrywise, \(L+L^\ast\) equals
\(J+I\): off the diagonal exactly one of \(L_{jk},L_{kj}\) equals one,
while on the diagonal both do.  Hence
\begin{align}
2\Re(b^\ast Lb)
&=
b^\ast(L+L^\ast)b
\label{eq:triangular_real_part_one}
\\
&=
b^\ast Jb+b^\ast b
\label{eq:triangular_real_part_two}
\\
&=
\left|\sum_{k=1}^{N}b_k\right|^2+N,
\label{eq:triangular_real_part_three}
\end{align}
because \(|b_k|=1\).  Division by \(2N\) proves the identity.
\end{proof}

\begin{corollary}
\label{cor:midpoint_triangular}
For midpoint evaluation,
\begin{equation}
\Re\left(
\frac1N\sum_{k=1}^N
\overline{b_k}S_k^{\mathrm{mid}}
\right)
=
\frac{|S_N|^2}{2N}.
\label{eq:midpoint_triangular}
\end{equation}
\end{corollary}

\begin{proof}
By Proposition~\ref{prop:endpoint_conversion},
\begin{equation}
\overline b_kS_k^{\mathrm{mid}}
=
\overline b_kS_k^{\mathrm{right}}-\frac12.
\label{eq:midpoint_corollary_step}
\end{equation}
Average over \(k\), take real parts, and apply
Theorem~\ref{thm:triangular_half}.
\end{proof}

\begin{warning}
\label{warn:triangular_scope}
Theorem~\ref{thm:triangular_half} is universal for unit phases.  It does
not identify prime-power structure and does not imply the Riemann
Hypothesis.  It applies exactly to the raw full-sample cumulative walk,
not automatically to a residual after smooth and prime-power subtraction.
\end{warning}

\section{Complete-zero weights and the positive horizontal defect}
\label{sec:positive_defect}

\begin{definition}
\label{def:complete_zero_weight}
For a positive ordinate \(\tau\), define
\begin{equation}
W_q(\tau)
=
\sum_{\rho:\Im\rho=\tau}
q^{\beta-\frac12},
\label{eq:Wq_definition}
\end{equation}
where zeros are counted with multiplicity.
\end{definition}

\begin{definition}
\label{def:Dq}
Define
\begin{equation}
D_q(\tau)
=
\sum_{\substack{\rho=\beta+i\tau\\\beta>1/2}}
\left[
\cosh\left((\beta-\tfrac12)\log q\right)-1
\right].
\label{eq:Dq_definition}
\end{equation}
\end{definition}

\begin{theorem}[Positive defect at each ordinate]
\label{thm:positive_defect}
For every \(q>1\) and every positive ordinate \(\tau\),
\begin{equation}
\boxed{
W_q(\tau)=m(\tau)+2D_q(\tau),
\qquad
D_q(\tau)\geq0.
}
\label{eq:W_m_D}
\end{equation}
Moreover,
\begin{equation}
D_q(\tau)=0
\label{eq:Dq_zero_condition}
\end{equation}
if and only if every zero at ordinate \(\tau\) lies on the critical line.
\end{theorem}

\begin{proof}
The functional equation and conjugation symmetry imply that, at a fixed
positive ordinate \(\tau\), every zero \(\beta+i\tau\) off the critical
line is accompanied with equal multiplicity by
\(1-\beta+i\tau\); see \cite[Chapter~2]{Titchmarsh1986}.  Put
\(\delta=\beta-1/2>0\).  The pair contributes to \(W_q(\tau)\)
\begin{align}
q^{\beta-1/2}+q^{1/2-\beta}
&=
q^\delta+q^{-\delta}
\label{eq:defect_pair_one}
\\
&=
2\cosh(\delta\log q)
\label{eq:defect_pair_two}
\\
&=
2+2\left[\cosh(\delta\log q)-1\right].
\label{eq:defect_pair_three}
\end{align}
The first term \(2\) is exactly the contribution of two unit weights to
the multiplicity \(m(\tau)\); the second is the excess.  A zero on the
critical line has \(\delta=0\) and contributes one to both quantities.
Summing all pairs proves \(W_q=m+2D_q\).  Since \(q>1\),
\(\log q\neq0\), and \(\cosh x-1\geq0\) with equality only at \(x=0\),
the defect vanishes precisely when no positive \(\delta\) occurs.
\end{proof}

\begin{definition}
\label{def:atomic_measures}
Define the locally finite complex measures
\begin{align}
\mu_q
&=
\sum_{\tau>0}
m(\tau)e^{i\tau\log q}\delta_\tau,
\label{eq:muq_measure}
\\
\nu_q
&=
\sum_{\tau>0}
W_q(\tau)e^{i\tau\log q}\delta_\tau.
\label{eq:nuq_measure}
\end{align}
\end{definition}

\begin{corollary}[Atomic positive criterion]
\label{cor:atomic_positive_criterion}
The exact measure identity
\begin{equation}
\boxed{
\frac12e^{-iT\log q}(\nu_q-\mu_q)
=
\sum_{\tau>0}D_q(\tau)\delta_\tau
}
\label{eq:positive_measure_identity}
\end{equation}
holds.  For any one fixed \(q>1\),
\begin{equation}
\nu_q=\mu_q
\quad\Longleftrightarrow\quad
\RH.
\label{eq:measure_RH_equivalence}
\end{equation}
\end{corollary}

\begin{proof}
At the atom \(T=\tau\), the coefficient of \(\nu_q-\mu_q\) is
\begin{equation}
\left[W_q(\tau)-m(\tau)\right]e^{i\tau\log q}
=
2D_q(\tau)e^{i\tau\log q}.
\label{eq:atomic_difference_coefficient}
\end{equation}
Multiplication by \(\frac12e^{-iT\log q}\) cancels the phase at each
atom and gives the positive measure in
\eqref{eq:positive_measure_identity}.  If \(\nu_q=\mu_q\), this measure
is zero, so every nonnegative mass \(D_q(\tau)\) is zero.  Theorem
\ref{thm:positive_defect} then places every zero on the critical line.
Conversely, under \(\RH\), every weight \(q^{\beta-1/2}\) equals one, so
\(W_q(\tau)=m(\tau)\) and the two measures coincide.
\end{proof}

\begin{remark}
\label{rem:q_sensitivity}
No value of \(q>1\) is logically privileged.  For a small displacement
\(\delta=\beta-1/2\),
\begin{equation}
\cosh(\delta\log q)-1
=
\frac12\delta^2(\log q)^2
+
O\left(\delta^4(\log q)^4\right).
\label{eq:q_sensitivity}
\end{equation}
Larger \(q\) amplifies small horizontal displacements but also enlarges
the analytic weights.
\end{remark}

\section{The explicit-formula origin of the two prime-power branches}
\label{sec:explicit_formula_branches}

The Guinand explicit formula and its later formulations relate suitable
test-function sums over zeta zeros to the von Mangoldt function
\(\Lambda(n)\), hence to prime powers \cite{Guinand1948,Weil1952}.
Landau's formula and Gonek's uniform refinements isolate the same
prime-power resonances in exponential sums over zeros
\cite{Gonek1985,FordZaharescu2005,FordSoundararajanZaharescu2009}.

\begin{definition}
\label{def:prime_power_coefficients}
Let \(n=p^r\) be a prime power, put
\begin{equation}
\ell_n=\log n,
\label{eq:elln_definition}
\end{equation}
and suppose \(\ell_n\neq\omega\).  Define
\begin{align}
c_n^-
&=
\frac{i\Lambda(n)}
{2\pi\sqrt n(\omega-\ell_n)},
\label{eq:cminus_definition}
\\
c_n^+
&=
\frac{i\Lambda(n)}
{2\pi\sqrt n(\omega+\ell_n)}.
\label{eq:cplus_definition}
\end{align}
The associated oscillators are
\begin{align}
P_n^-(T)
&=
c_n^-e^{i(\omega-\ell_n)T},
\label{eq:Pminus_definition}
\\
P_n^+(T)
&=
c_n^+e^{i(\omega+\ell_n)T}.
\label{eq:Pplus_definition}
\end{align}
\end{definition}

\begin{heuristic}[Sharp-window prime-spectral ansatz]
\label{heur:sharp_window_ansatz}
For a sharp counting window, formal integration of the explicit formula
suggests a decomposition
\begin{equation}
S_q(T)
=
C_q+
A_{\mathrm{smooth},q}(T)
+
\frac12e^{i\omega T}
+
\sum_{n=p^r} \left(P_n^-(T)+P_n^+(T)\right),
\label{eq:formal_sharp_decomposition}
\end{equation}
where the infinite prime-power sum requires a specified truncation or
regularization.  This is a heuristic argument when used with the
discontinuous indicator \(\1_{\gamma\leq T}\).
\end{heuristic}

\begin{definition}[Smooth term used in the computation]
\label{def:Asmooth}
The MATLAB program uses
\begin{equation}
A_{\mathrm{smooth}}(T)
=
\frac{
e^{i\omega T}\log(T/2\pi)-\Ei(i\omega T)
}{
2\pi i\omega
},
\label{eq:Asmooth_definition}
\end{equation}
up to a branch-dependent additive constant.  The constant is treated by
a joint constant--carrier fit.
\end{definition}

\begin{warning}
\label{warn:Asmooth_status}
Equation \eqref{eq:Asmooth_definition} is the antiderivative of the
smooth Riemann--von Mangoldt density multiplied by \(e^{i\omega T}\).
It is not, by itself, a complete sharp-window explicit formula.
Gamma-factor, pole, endpoint, branch, and regularization conventions must
be retained when an exact global identity is claimed.
\end{warning}

\section{Exact algebra of each prime-power mode}
\label{sec:mode_algebra}

\begin{proposition}[Sum, difference, and ratio identities]
\label{prop:coefficient_identities}
For every nonresonant prime power \(n\),
\begin{align}
c_n^-+c_n^+
&=
\frac{i\Lambda(n)\omega}
{\pi\sqrt n(\omega^2-\ell_n^2)},
\label{eq:coefficient_sum}
\\
c_n^--c_n^+
&=
\frac{i\Lambda(n)\ell_n}
{\pi\sqrt n(\omega^2-\ell_n^2)},
\label{eq:coefficient_difference}
\\
\frac{c_n^-}{c_n^+}
&=
\frac{\omega+\ell_n}{\omega-\ell_n}.
\label{eq:coefficient_ratio}
\end{align}
\end{proposition}

\begin{proof}
Starting from the definitions, combine denominators:
\begin{align}
c_n^-+c_n^+
&=
\frac{i\Lambda(n)}{2\pi\sqrt n}
\left[
\frac1{\omega-\ell_n}+\frac1{\omega+\ell_n}
\right]
\label{eq:coefficient_sum_proof_one}
\\
&=
\frac{i\Lambda(n)}{2\pi\sqrt n}
\frac{2\omega}{\omega^2-\ell_n^2},
\label{eq:coefficient_sum_proof_two}
\end{align}
which is \eqref{eq:coefficient_sum}.  Replacing the plus by a minus in
the bracket gives \(2\ell_n/(\omega^2-\ell_n^2)\), proving
\eqref{eq:coefficient_difference}.  Finally the common factor cancels in
the quotient:
\begin{equation}
\frac{c_n^-}{c_n^+}
=
\frac{\omega+\ell_n}{\omega-\ell_n}.
\label{eq:coefficient_ratio_proof}
\end{equation}
All identities are exact whenever \(\ell_n\neq\omega\).
\end{proof}

\begin{definition}
\label{def:u_v_d}
Write
\begin{equation}
c_n^-=ia_n^-,
\qquad
c_n^+=ia_n^+,
\qquad
a_n^\pm\in\mathbb R,
\label{eq:a_pm_definition}
\end{equation}
and define
\begin{align}
u_n
&=
a_n^-+a_n^+,
\label{eq:un_definition}
\\
v_n
&=
a_n^--a_n^+,
\label{eq:vn_definition}
\\
d_n
&=
|c_n^-|^2-|c_n^+|^2.
\label{eq:dn_definition}
\end{align}
\end{definition}

\begin{proposition}[Signed mode invariant]
\label{prop:signed_mode_invariant}
The quantities in Definition~\ref{def:u_v_d} satisfy
\begin{align}
u_n
&=
-\frac{\omega\Lambda(n)}
{\pi\sqrt n(\ell_n^2-\omega^2)},
\label{eq:un_formula}
\\
v_n
&=
-\frac{\Lambda(n)\ell_n}
{\pi\sqrt n(\ell_n^2-\omega^2)},
\label{eq:vn_formula}
\\
\boxed{
d_n=u_nv_n
=
\frac{\omega\Lambda(n)^2\ell_n}
{\pi^2n(\ell_n^2-\omega^2)^2}.
}
\label{eq:dn_formula}
\end{align}
\end{proposition}

\begin{proof}
Because the denominators in \(c_n^\pm\) are real, the coefficients are
purely imaginary.  Taking their real scalar factors and using the sum and
difference identities gives
\begin{align}
u_n
&=
\frac{\Lambda(n)}{2\pi\sqrt n}
\left[
\frac1{\omega-\ell_n}+\frac1{\omega+\ell_n}
\right]
\label{eq:un_proof_one}
\\
&=
-\frac{\omega\Lambda(n)}
{\pi\sqrt n(\ell_n^2-\omega^2)},
\label{eq:un_proof_two}
\\
v_n
&=
\frac{\Lambda(n)}{2\pi\sqrt n}
\left[
\frac1{\omega-\ell_n}-\frac1{\omega+\ell_n}
\right]
\label{eq:vn_proof_one}
\\
&=
-\frac{\ell_n\Lambda(n)}
{\pi\sqrt n(\ell_n^2-\omega^2)}.
\label{eq:vn_proof_two}
\end{align}
The algebraic factorization
\begin{equation}
(a_n^-)^2-(a_n^+)^2
=
(a_n^-+a_n^+)(a_n^--a_n^+)
=u_nv_n
\label{eq:dn_difference_squares}
\end{equation}
then gives the displayed formula for \(d_n\).  Its positivity follows
from \(\omega>0\), \(\Lambda(n)^2\geq0\), \(\ell_n>0\), and the squared
denominator.
\end{proof}

\begin{proposition}[Elliptic form and anisotropy]
\label{prop:ellipse_anisotropy}
After demodulation by \(e^{-i\omega T}\), the mode is
\begin{equation}
g_n(T)
=
c_n^-e^{-i\ell_nT}
+
c_n^+e^{i\ell_nT}
=
\frac{\Lambda(n)}
{\pi\sqrt n(\omega^2-\ell_n^2)}
\left[
\ell_n\sin(\ell_nT)
+
i\omega\cos(\ell_nT)
\right].
\label{eq:gn_ellipse}
\end{equation}
Its unsigned elliptic area is
\begin{equation}
\pi d_n.
\label{eq:ellipse_area}
\end{equation}
For uniform phase, its anisotropy is
\begin{equation}
\mathcal A_n
=
\frac{\ell_n^2-\omega^2}{\ell_n^2+\omega^2}
\longrightarrow1
\qquad(n\to\infty).
\label{eq:single_mode_anisotropy}
\end{equation}
\end{proposition}

\begin{proof}
Using \(e^{ix}=\cos x+i\sin x\), expand
\(c_n^-e^{-i\ell_nT}+c_n^+e^{i\ell_nT}\) and collect its real and
imaginary parts.  This gives
\begin{equation}
g_n(T)=A_n\sin(\ell_nT)+iB_n\cos(\ell_nT),
\label{eq:ellipse_parametric_proof}
\end{equation}
with
\begin{equation}
A_n=
\frac{\Lambda(n)\ell_n}{\pi\sqrt n(\omega^2-\ell_n^2)},
\qquad
B_n=
\frac{\Lambda(n)\omega}{\pi\sqrt n(\omega^2-\ell_n^2)}.
\label{eq:ellipse_axes_proof}
\end{equation}
This is an ellipse with semiaxes \(|A_n|\) and \(|B_n|\).  Their product
is \(|d_n|=d_n\), so its unsigned area is \(\pi d_n\).  Under uniform
phase, the coordinate variances are \(A_n^2/2\) and \(B_n^2/2\).  The
normalized difference of principal variances is therefore
\begin{equation}
\frac{A_n^2-B_n^2}{A_n^2+B_n^2}
=
\frac{\ell_n^2-\omega^2}{\ell_n^2+\omega^2},
\label{eq:anisotropy_proof_ratio}
\end{equation}
which tends to one as \(\ell_n\to\infty\).
\end{proof}

\begin{lemma}[Dominance of \(n=3\) in the minus branch for \(\omega=1\)]
\label{lem:n3_dominance}
For \(\omega=1\), the largest coefficient
\(|c_n^-|\) among prime powers \(n\geq2\) occurs at \(n=3\).
\end{lemma}

\begin{proof}
For a prime power \(n=p^r\), \(\Lambda(n)=\log p\leq\log n\).  When
\(n\geq3\) and \(\omega=1\), \(\log n>1\), so
\begin{equation}
|c_n^-|
\leq
f(n)
:=
\frac{\log n}{2\pi\sqrt n(\log n-1)}.
\label{eq:n3_f_definition}
\end{equation}
Let \(L=\log x>1\).  Logarithmic differentiation gives
\begin{align}
x\frac{f'(x)}{f(x)}
&=
\frac1L-\frac12-\frac1{L-1}
\label{eq:n3_log_derivative_one}
\\
&=
-\frac1{L(L-1)}-\frac12<0.
\label{eq:n3_log_derivative_two}
\end{align}
Thus \(f\) is strictly decreasing on \((e,\infty)\).  At \(n=3\), the
number is prime and \(\Lambda(3)=\log3\), so equality holds in
\eqref{eq:n3_f_definition}; hence \(n=3\) dominates every \(n\geq4\).
The remaining comparison with \(n=2\) is direct.  Using the certified
bounds \(0.6931<\log2<0.6932\) and
\(1.0986<\log3<1.0987\) in the two explicit formulas gives
\(|c_3^-|>|c_2^-|\).  Therefore \(n=3\) is the unique maximum.
\end{proof}

\section{The convergent signed tail}
\label{sec:signed_tail}

\begin{definition}
\label{def:Delta}
Define
\begin{align}
D_\infty
&=
\sum_{n=p^r}d_n,
\label{eq:D_infinity}
\\
\Delta_\omega(Q)
&=
\sum_{\substack{n=p^r\\n>Q}}d_n
=
D_\infty-\sum_{\substack{n=p^r\\n\leq Q}}d_n.
\label{eq:Delta_definition}
\end{align}
\end{definition}

\begin{theorem}[Asymptotic signed tail]
\label{thm:Delta_asymptotic}
For fixed \(\omega>0\),
\begin{equation}
\boxed{
\Delta_\omega(Q)
\sim
\frac{\omega}{\pi^2\log Q}
\qquad(Q\to\infty).
}
\label{eq:Delta_asymptotic}
\end{equation}
\end{theorem}

\begin{proof}
Separate primes from proper prime powers.  For a prime \(p\),
\(\Lambda(p)=\ell_p=\log p\), and the exact formula becomes
\begin{equation}
d_p
=
\frac{\omega}{\pi^2p\log p}
\left(1-\frac{\omega^2}{(\log p)^2}\right)^{-2}.
\label{eq:Delta_dp_exact_proof}
\end{equation}
For \(p>e^{2\omega}\), expand \((1-z)^{-2}=1+O(z)\) with
\(z=\omega^2/(\log p)^2\).  Uniformly in that range,
\begin{equation}
d_p
=
\frac{\omega}{\pi^2p\log p}
+
O_\omega\left(\frac1{p(\log p)^3}\right).
\label{eq:Delta_dp_error_proof}
\end{equation}
The error tail is \(O_\omega((\log Q)^{-3})\) by the prime number theorem
and partial summation.

For the main term, write the prime sum as a Stieltjes integral:
\begin{equation}
\sum_{p>Q}\frac1{p\log p}
=
\int_Q^\infty\frac1{x\log x}\,\dd\pi(x).
\label{eq:Delta_stieltjes_prime}
\end{equation}
The prime number theorem \(\pi(x)=x/\log x+o(x/\log x)\), followed by
partial summation, gives
\begin{equation}
\int_Q^\infty\frac1{x\log x}\,\dd\pi(x)
=
\int_Q^\infty\frac{\dd x}{x(\log x)^2}
+
o\left(\frac1{\log Q}\right)
=
\frac1{\log Q}+o\left(\frac1{\log Q}\right).
\label{eq:Delta_prime_main_proof}
\end{equation}

It remains to bound \(n=p^r\) with \(r\geq2\).  Outside a fixed finite
set depending on \(\omega\), the denominator satisfies
\((r^2(\log p)^2-\omega^2)^2\gg_\omega r^4(\log p)^4\).  Hence
\begin{equation}
d_{p^r}
\ll_\omega
\frac1{r^3p^r\log p}.
\label{eq:Delta_proper_power_bound}
\end{equation}
Summing over \(p^r>Q\), \(r\geq2\), gives
\(O_\omega(Q^{-1/2})\), after separating the finitely many small primes.
This is \(o(1/\log Q)\).  Combining the prime and proper-power parts
proves the theorem.
\end{proof}

\begin{remark}
\label{rem:Delta_scale}
For \(\omega=1\),
\begin{align}
\pi\Delta_1(Q)
&\sim
\frac1{\pi\log Q},
\label{eq:pi_Delta_scale}
\\
\frac34\Delta_1(Q)
&\sim
\frac3{4\pi^2\log Q}.
\label{eq:three_quarter_Delta_scale}
\end{align}
These constants match two finite numerical quantities, but the matching
relations themselves remain conjectural.
\end{remark}

\section{Finite zero-sampled carrier identities}
\label{sec:exact_area_defect}

\begin{definition}
\label{def:abstract_sampled_model}
Fix a finite multiset of positive ordinates
\begin{equation}
\Gamma_H=\{\gamma_k:0<\gamma_k\leq H\},
\qquad N(H)=|\Gamma_H|,
\label{eq:abstract_GammaH}
\end{equation}
and let $Y$ be any complex-valued function on $\Gamma_H$.  For a cutoff
$Q$, define
\begin{equation}
R_Q(\gamma)
=
Y(\gamma)-
\sum_{\substack{n=p^r\\n\leq Q}}
\left(P_n^-(\gamma)+P_n^+(\gamma)\right).
\label{eq:abstract_RQ}
\end{equation}
The raw carrier coefficient and zero kernel are
\begin{align}
\widehat c_{Q,H}
&=
\frac1{N(H)}\sum_{\gamma\in\Gamma_H}
e^{-i\omega\gamma}R_Q(\gamma),
\label{eq:abstract_raw_carrier}
\\
M_H(x)
&=
\frac1{N(H)}\sum_{\gamma\in\Gamma_H}x^{i\gamma},
\qquad x>0.
\label{eq:abstract_MH}
\end{align}
Also put
\begin{equation}
\widehat c_{Y,H}
=
\frac1{N(H)}\sum_{\gamma\in\Gamma_H}
e^{-i\omega\gamma}Y(\gamma).
\label{eq:abstract_cYH}
\end{equation}
\end{definition}

\begin{theorem}[Exact finite-sample carrier formula]
\label{thm:abstract_carrier_formula}
For the objects in Definition~\ref{def:abstract_sampled_model},
\begin{equation}
\boxed{
\widehat c_{Q,H}
=
\widehat c_{Y,H}
-
\sum_{\substack{n=p^r\\n\leq Q}}
\left(c_n^-\overline{M_H(n)}+c_n^+M_H(n)\right).
}
\label{eq:abstract_carrier_formula}
\end{equation}
\end{theorem}

\begin{proof}
For each sampled ordinate, demodulation gives
\begin{align}
e^{-i\omega\gamma_k}P_n^-(\gamma_k)
&=
c_n^-e^{-i\ell_n\gamma_k}
=
c_n^-n^{-i\gamma_k},
\label{eq:abstract_carrier_minus_proof}
\\
e^{-i\omega\gamma_k}P_n^+(\gamma_k)
&=
c_n^+e^{i\ell_n\gamma_k}
=
c_n^+n^{i\gamma_k}.
\label{eq:abstract_carrier_plus_proof}
\end{align}
Insert the definition of the residual into the sample average defining
\(\widehat c_{Q,H}\).  Finite sums may be interchanged without any
convergence issue.  The average of \(n^{i\gamma_k}\) is \(M_H(n)\), and
the average of \(n^{-i\gamma_k}\) is its complex conjugate because
\(n>0\).  Collecting the terms yields the formula exactly.
\end{proof}


\begin{theorem}[Exact imaginary-part formula]
\label{thm:imaginary_part_identity}
Write
\begin{equation}
M_H(n)=x_n+iy_n.
\label{eq:abstract_MH_xy}
\end{equation}
Then
\begin{equation}
\Im\left(c_n^-\overline{M_H(n)}+c_n^+M_H(n)\right)
=u_nx_n,
\label{eq:abstract_mode_imaginary_projection}
\end{equation}
and consequently
\begin{equation}
\boxed{
\Im\widehat c_{Q,H}
=
\Im\widehat c_{Y,H}
-
\sum_{\substack{n=p^r\\n\leq Q}}
u_n\Re M_H(n).
}
\label{eq:abstract_imaginary_part_identity}
\end{equation}
\end{theorem}

\begin{proof}
Write \(M_H(n)=x_n+iy_n\) and \(c_n^\pm=ia_n^\pm\) with real
\(a_n^\pm\).  Then
\begin{align}
c_n^-\overline{M_H(n)}+c_n^+M_H(n)
&=
ia_n^-(x_n-iy_n)+ia_n^+(x_n+iy_n)
\label{eq:imaginary_identity_expand_one}
\\
&=
(a_n^--a_n^+)y_n
+i(a_n^-+a_n^+)x_n
\label{eq:imaginary_identity_expand_two}
\\
&=
v_ny_n+iu_nx_n.
\label{eq:imaginary_identity_expand_three}
\end{align}
Its imaginary part is therefore \(u_nx_n=u_n\Re M_H(n)\).  Taking
imaginary parts in the exact carrier formula and summing proves the result.
\end{proof}

\begin{definition}
\label{def:CH}
Define the \(Q\)-independent quantity
\begin{equation}
C_H
=
\Im\widehat c_{Y,H}-\pi D_\infty.
\label{eq:CH_definition}
\end{equation}
\end{definition}

\begin{theorem}[Exact global defect identity]
\label{thm:global_defect_identity}
For the finite-sample raw coefficient,
\begin{equation}
\boxed{
\Im\widehat c_{Q,H}-\pi\Delta_\omega(Q)
=
C_H
-
\sum_{\substack{n=p^r\\n\leq Q}}
u_n\left[
\Re M_H(n)-\pi v_n
\right].
}
\label{eq:global_defect_identity}
\end{equation}
\end{theorem}

\begin{proof}
The exact carrier identity gives
\begin{equation}
\Im\widehat c_{Q,H}
=
\Im\widehat c_{Y,H}
-
\sum_{n\leq Q}u_n\Re M_H(n).
\label{eq:global_defect_proof_carrier}
\end{equation}
On the other hand, \(d_n=u_nv_n\) and
\(\Delta_\omega(Q)=D_\infty-\sum_{n\leq Q}d_n\), so
\begin{equation}
\pi\Delta_\omega(Q)
=
\pi D_\infty
-
\pi\sum_{n\leq Q}u_nv_n.
\label{eq:global_defect_proof_Delta}
\end{equation}
Subtract the second equation from the first and group the cutoff-independent
terms:
\begin{align}
\Im\widehat c_{Q,H}-\pi\Delta_\omega(Q)
={}&
\left(\Im\widehat c_{Y,H}-\pi D_\infty\right)
\label{eq:global_defect_proof_group_one}
\\
&-
\sum_{n\leq Q}u_n
\left[\Re M_H(n)-\pi v_n\right].
\label{eq:global_defect_proof_group_two}
\end{align}
The parenthesis is exactly \(C_H\).  Every manipulation is finite except
for \(D_\infty\), whose convergence follows from
Theorem~\ref{thm:Delta_asymptotic}.
\end{proof}

\begin{corollary}[Equivalent transport target]
\label{cor:transport_target}
The conjectural relation
\begin{equation}
\Im\widehat c_{Q,H}
=
\pi\Delta_\omega(Q)
+
o\left(\frac1{\log Q}\right)
\label{eq:transport_relation}
\end{equation}
is exactly equivalent to
\begin{equation}
\boxed{
C_H
-
\sum_{n\leq Q}
u_n\left[
\Re M_H(n)-\pi v_n
\right]
=
o\left(\frac1{\log Q}\right).
}
\label{eq:equivalent_transport_target}
\end{equation}
\end{corollary}

\begin{proof}
Subtract \(\pi\Delta_\omega(Q)\) from
\(\Im\widehat c_{Q,H}\) and apply
Theorem~\ref{thm:global_defect_identity}.  Since
\(\Delta_\omega(Q)\asymp1/\log Q\) by
Theorem~\ref{thm:Delta_asymptotic}, the stated asymptotic is equivalent
to the right-hand side of \eqref{eq:global_defect_identity} being
\(o(1/\log Q)\).
\end{proof}

\begin{theorem}[Exact increment identity]
\label{thm:increment_identity}
When \(Q\) crosses a prime power \(n\),
\begin{align}
\Delta\Im\widehat c_{Q,H}
&=
-u_n\Re M_H(n),
\label{eq:increment_c}
\\
\Delta\left(\pi\Delta_\omega(Q)\right)
&=
-\pi u_nv_n.
\label{eq:increment_Delta}
\end{align}
If \(u_nv_n\neq0\), then
\begin{equation}
\boxed{
\frac{
\Delta\Im\widehat c_{Q,H}
}{
\Delta(\pi\Delta_\omega(Q))
}
=
\frac{\Re M_H(n)}{\pi v_n}.
}
\label{eq:increment_ratio}
\end{equation}
\end{theorem}

\begin{proof}
The global defect identity is exact for every finite cutoff:
\begin{equation}
\Im\widehat c_{Q,H}-\pi\Delta_\omega(Q)
=
C_H-
\sum_{n\leq Q}u_n\bigl(\Re M_H(n)-\pi v_n\bigr).
\label{eq:transport_equivalence_repeated}
\end{equation}
Therefore the left side is \(o(1/\log Q)\) if and only if the right side is
\(o(1/\log Q)\).  Since
\(\Delta_\omega(Q)\sim\omega/(\pi^2\log Q)\), the same equivalence holds
with \(o(\Delta_\omega(Q))\) in place of \(o(1/\log Q)\).  No estimate or
limit interchange is hidden in this reformulation.
\end{proof}

\section{Landau scaling and failure of termwise area transport}
\label{sec:landau_scaling}

Landau's formula concerns sums of \(x^\rho\), not sums of \(x^{i\gamma}\)
without horizontal weights.  Passing from the former to the latter requires
\(\RH\), or a separate control of the real parts.  Uniformity when \(x\)
grows with \(H\) requires a uniform Landau--Gonek formula
\cite{Gonek1985,FordZaharescu2005}.

\begin{heuristic}[Conditional Landau estimate for the ordinate kernel]
\label{heur:landau_kernel}
Assume \(\RH\) and a uniform Landau--Gonek estimate in the range under
consideration.  With
\begin{equation}
L_H=\log(H/2\pi),
\label{eq:LH_definition}
\end{equation}
one expects at a prime power \(n\)
\begin{equation}
\Re M_H(n)
\approx
-\frac{\Lambda(n)}{\sqrt n\,L_H}.
\label{eq:MH_landau_heuristic}
\end{equation}
\end{heuristic}

\begin{proposition}[Required termwise target]
\label{prop:termwise_target}
Individual transport of the signed area of mode \(n\) would require
\begin{equation}
\Re M_H(n)\sim\pi v_n.
\label{eq:termwise_target}
\end{equation}
For a prime \(p\),
\begin{equation}
\pi v_p
=
-\frac{(\log p)^2}
{\sqrt p((\log p)^2-\omega^2)}
\sim
-\frac1{\sqrt p}.
\label{eq:pi_vp_asymptotic}
\end{equation}
\end{proposition}

\begin{proof}
The increment identity \eqref{eq:increment_ratio} shows that exact
mode-by-mode transport is equivalent to
\(\Re M_H(n)/(\pi v_n)\to1\), which is
\eqref{eq:termwise_target}.  For a prime \(p\), substitute
\(\Lambda(p)=\log p\) and \(\ell_p=\log p\) into
\eqref{eq:vn_formula}; division of numerator and denominator by
\((\log p)^2\) gives \eqref{eq:pi_vp_asymptotic}.
\end{proof}

\begin{heuristic}[Termwise ratio]
\label{heur:termwise_ratio}
Combining \eqref{eq:MH_landau_heuristic} and
\eqref{eq:pi_vp_asymptotic} gives
\begin{equation}
\frac{\Re M_H(p)}{\pi v_p}
\approx
\frac{\log p}{L_H}.
\label{eq:termwise_ratio}
\end{equation}
If \(p\asymp H^\theta\), this predicts
\begin{equation}
\frac{\Re M_H(p)}{\pi v_p}
\approx\theta.
\label{eq:theta_ratio}
\end{equation}
Thus termwise transport cannot be asymptotically exact at a fixed
\(\theta<1\); it becomes plausible only for
\(p=H^{1-o(1)}\).
\end{heuristic}

\begin{warning}
\label{warn:uniform_landau}
The fixed-\(p\) Landau formula alone does not prove
\eqref{eq:theta_ratio}.  The latter is a falsifiable scaling prediction
that requires uniformity in the joint range \(p=p(H)\).
\end{warning}

\begin{conjecture}[Local transport near the height scale]
\label{conj:local_transport}
Let
\begin{equation}
Q_1,Q_2=H^{1-o(1)}.
\label{eq:local_Q_range}
\end{equation}
In a suitable uniform regime,
\begin{equation}
\Im\widehat c_{Q_2,H}-\Im\widehat c_{Q_1,H}
\sim
\pi\left[
\Delta_\omega(Q_2)-\Delta_\omega(Q_1)
\right].
\label{eq:local_transport_conjecture}
\end{equation}
\end{conjecture}

\begin{proposition}[Exact local defect]
\label{prop:local_defect}
The defect in Conjecture~\ref{conj:local_transport} is exactly
\begin{align}
&
\left[
\Im\widehat c_{Q_2,H}-\Im\widehat c_{Q_1,H}
\right]
-
\pi\left[
\Delta_\omega(Q_2)-\Delta_\omega(Q_1)
\right]
\notag
\\
&\qquad
=
-\sum_{Q_1<n\leq Q_2}
u_n\left[
\Re M_H(n)-\pi v_n
\right].
\label{eq:local_defect_identity}
\end{align}
\end{proposition}

\begin{proof}
Subtract \eqref{eq:global_defect_identity} at \(Q_1\) and \(Q_2\).
The constant \(C_H\) cancels.
\end{proof}

\section{Mixed carrier area versus prime self-area}
\label{sec:mixed_vs_self_area}

\begin{proposition}[Mixed form measured by the centroid coefficient]
\label{prop:mixed_area}
The imaginary part of the raw carrier coefficient is
\begin{align}
\Im\widehat c_{Q,H}
&=
\frac1{N(H)}
\sum_{\gamma_k\leq H}
\Im\left[
e^{-i\omega\gamma_k}R_Q(\gamma_k)
\right]
\label{eq:mixed_area_first}
\\
&=
-\frac1{N(H)}
\sum_{\gamma_k\leq H}
\Im\left[
\overline{R_Q(\gamma_k)}
e^{i\omega\gamma_k}
\right].
\label{eq:mixed_area_second}
\end{align}
It is a mixed bilinear pairing between the residual position and the
actual phase increment.
\end{proposition}

\begin{proof}
Take the imaginary part of the defining average
\eqref{eq:abstract_raw_carrier}.  For the second equality, use
\begin{equation}
\Im(e^{-i\omega\gamma}R)
=
-\Im(\overline R\,e^{i\omega\gamma}),
\label{eq:mixed_area_conjugation}
\end{equation}
which follows by complex conjugation.
\end{proof}

\begin{definition}
\label{def:prime_self_area}
For a finite prime-power model \(G\), define the formal self-area
\begin{equation}
\mathfrak A_H(G)
=
-\frac{\pi}{H}
\int
\Im\left(
\overline{G(t)}\,\cJ G(t)
\right)\dd t,
\label{eq:prime_self_area}
\end{equation}
where \(\cJ\) is a specified frequency-sign or Hilbert-type operator.
\end{definition}

\begin{warning}
\label{warn:mixed_not_self}
The mixed pairing \eqref{eq:mixed_area_second} is not identical to the
self-pairing \eqref{eq:prime_self_area}.  Such an identification would
require a justified replacement
\begin{equation}
e^{i\omega\gamma_k}
\longleftrightarrow
\cJ G_Q(\gamma_k).
\label{eq:unproved_operator_replacement}
\end{equation}
The exact termwise defect of this replacement is encoded by
\(\Re M_H(n)-\pi v_n\).
\end{warning}

\section{Exact quadratic form for the sampled residual variance}
\label{sec:exact_variance_form}

The numerical variance is the squared norm of the \emph{full residual}
after jointly removing the constant and carrier.  It is therefore not the
projected norm of the prime-power model alone.  The background vector, the
mixed background--prime term, and the projected prime Gram form all enter
exactly.

\begin{definition}[Sampled inner product]
\label{def:sample_inner_product}
For sampled vectors on $\Gamma_H$, define
\begin{equation}
\langle \mathbf f,\mathbf g\rangle_H
=
\frac1{N(H)}\mathbf f^\ast\mathbf g
=
\frac1{N(H)}\sum_{\gamma_k\leq H}
\overline{f(\gamma_k)}g(\gamma_k).
\label{eq:sample_inner_product}
\end{equation}
\end{definition}

\begin{definition}[Sampled vectors and prime-mode matrix]
\label{def:sampled_variance_vectors}
Let
\begin{align}
\mathbf 1_H
&=
(1)_{\gamma_k\in\Gamma_H},
\label{eq:sampled_one_vector}
\\
\mathbf a_H
&=
(e^{i\omega\gamma_k})_{\gamma_k\in\Gamma_H},
\label{eq:sampled_carrier_vector}
\\
\mathbf y_H
&=
(Y(\gamma_k))_{\gamma_k\in\Gamma_H}.
\label{eq:sampled_background_vector}
\end{align}
For a prime power $n$ and $\sigma\in\{-1,+1\}$, define the original,
non-demodulated branch vector
\begin{equation}
\mathbf w_{n,\sigma}
=
\left(e^{i(\omega+\sigma\log n)\gamma_k}\right)_{\gamma_k\in\Gamma_H}.
\label{eq:original_branch_vector}
\end{equation}
For a cutoff $Q$, let
\begin{equation}
\mathcal I_Q
=
\{(n,\sigma):n=p^r\leq Q,\ \sigma\in\{-1,+1\}\},
\label{eq:branch_index_set}
\end{equation}
let $V_Q$ be the matrix whose columns are the vectors
$\mathbf w_{n,\sigma}$ in any fixed order, and let $\mathbf c_Q$ be the
coefficient vector with entries $c_n^\sigma$.  Then
\begin{equation}
\mathbf p_Q=V_Q\mathbf c_Q
\label{eq:sampled_prime_model_vector}
\end{equation}
is the sampled prime-power model, and
\begin{equation}
\mathbf r_Q=\mathbf y_H-V_Q\mathbf c_Q
\label{eq:sampled_full_residual_vector}
\end{equation}
is the sampled residual used by the joint fit.
\end{definition}

\begin{definition}[Joint projection and numerical variance]
\label{def:joint_projection_variance}
Put
\begin{equation}
B_H=
\begin{pmatrix}
\mathbf 1_H&\mathbf a_H
\end{pmatrix}.
\label{eq:BH_definition}
\end{equation}
Assume $|M_H(e^\omega)|<1$, equivalently that the two columns of $B_H$
are linearly independent, and define
\begin{equation}
\Pi_H
=
I-B_H(B_H^\ast B_H)^{-1}B_H^\ast.
\label{eq:PiH_exact_definition}
\end{equation}
The variance reported after the joint constant--carrier fit is
\begin{equation}
\boxed{
V(Q,H)
=
\frac1{N(H)}\,\mathbf r_Q^\ast\Pi_H\mathbf r_Q.
}
\label{eq:numerical_variance_exact_definition}
\end{equation}
\end{definition}

\begin{remark}
\label{rem:degenerate_joint_projection}
If $|M_H(e^\omega)|=1$, the same definition is made with the orthogonal
projection onto the complement of
$\operatorname{span}\{\mathbf 1_H,\mathbf a_H\}$, or equivalently with
the Moore--Penrose inverse.  The data used in the numerical chapter are in
the nondegenerate case.
\end{remark}

\begin{proposition}[Exact unprojected branch Gram kernel]
\label{prop:unprojected_gram}
For $j=(n,\sigma)$ and $k=(m,\tau)$ in $\mathcal I_Q$,
\begin{equation}
K^{(0)}_{Q,H}(j,k)
:=
\frac1{N(H)}\mathbf w_{n,\sigma}^\ast\mathbf w_{m,\tau}
=
M_H(m^\tau n^{-\sigma}).
\label{eq:unprojected_gram}
\end{equation}
Thus the four branch interactions between $p^r$ and $q^s$ involve
\begin{equation}
M_H(q^s/p^r),
\quad
M_H(q^sp^r),
\quad
M_H((q^sp^r)^{-1}),
\quad
M_H(p^r/q^s).
\label{eq:four_gram_arguments}
\end{equation}
\end{proposition}

\begin{proof}
The common carrier cancels in the inner product.  More explicitly,
\begin{align}
\frac1{N(H)}\mathbf w_{n,\sigma}^\ast\mathbf w_{m,\tau}
&=
\frac1{N(H)}
\sum_{\gamma_k\leq H}
 e^{-i(\omega+\sigma\log n)\gamma_k}
 e^{ i(\omega+\tau\log m)\gamma_k}
\label{eq:branch_gram_expansion}
\\
&=
\frac1{N(H)}
\sum_{\gamma_k\leq H}
(m^\tau n^{-\sigma})^{i\gamma_k}
\label{eq:branch_gram_kernel_step}
\\
&=
M_H(m^\tau n^{-\sigma}).
\label{eq:branch_gram_kernel_result}
\end{align}
The four sign choices give \eqref{eq:four_gram_arguments}.
\end{proof}

\begin{proposition}[Exact projected branch Gram kernel]
\label{prop:schur_kernel}
Put $q=e^\omega$ and
\begin{equation}
G_H
=
\frac1{N(H)}B_H^\ast B_H
=
\begin{pmatrix}
1&M_H(q)\\
\overline{M_H(q)}&1
\end{pmatrix}.
\label{eq:joint_gram}
\end{equation}
For $j=(n,\sigma)$ define
\begin{equation}
\mathbf b_j
=
\frac1{N(H)}B_H^\ast\mathbf w_{n,\sigma}
=
\begin{pmatrix}
M_H(qn^\sigma)\\
M_H(n^\sigma)
\end{pmatrix}.
\label{eq:branch_projection_vector}
\end{equation}
Then
\begin{equation}
\boxed{
\widetilde K_{Q,H}(j,k)
:=
\frac1{N(H)}\mathbf w_j^\ast\Pi_H\mathbf w_k
=
K^{(0)}_{Q,H}(j,k)
-
\mathbf b_j^\ast G_H^{-1}\mathbf b_k.
}
\label{eq:schur_kernel_formula}
\end{equation}
\end{proposition}

\begin{proof}
From \eqref{eq:PiH_exact_definition},
\begin{align}
\frac1{N(H)}V_Q^\ast\Pi_HV_Q
&=
\frac1{N(H)}V_Q^\ast V_Q
-
\frac1{N(H)}V_Q^\ast B_H
(B_H^\ast B_H)^{-1}B_H^\ast V_Q.
\label{eq:projected_gram_matrix_expansion}
\end{align}
Since $B_H^\ast B_H=N(H)G_H$, the $(j,k)$ entry of the second term is
$\mathbf b_j^\ast G_H^{-1}\mathbf b_k$.  The two components of
$\mathbf b_j$ follow from
\begin{align}
\frac1{N(H)}\mathbf 1_H^\ast\mathbf w_{n,\sigma}
&=
M_H(qn^\sigma),
\label{eq:branch_projection_constant_component}
\\
\frac1{N(H)}\mathbf a_H^\ast\mathbf w_{n,\sigma}
&=
M_H(n^\sigma).
\label{eq:branch_projection_carrier_component}
\end{align}
\end{proof}

\begin{definition}[Projected background and mixed vector]
\label{def:background_mixed_vector}
Define
\begin{align}
V_Y(H)
&=
\frac1{N(H)}\mathbf y_H^\ast\Pi_H\mathbf y_H,
\label{eq:projected_background_energy}
\\
\mathbf d_{Q,H}
&=
\frac1{N(H)}V_Q^\ast\Pi_H\mathbf y_H,
\label{eq:mixed_background_prime_vector}
\\
\widetilde K_{Q,H}
&=
\frac1{N(H)}V_Q^\ast\Pi_HV_Q.
\label{eq:projected_prime_gram_matrix}
\end{align}
\end{definition}

\begin{proposition}[Entrywise formula for the mixed vector]
\label{prop:mixed_vector_entrywise}
Let
\begin{equation}
\mathbf s_{Y,H}
=
\frac1{N(H)}B_H^\ast\mathbf y_H.
\label{eq:background_removed_coordinates}
\end{equation}
For $j=(n,\sigma)$,
\begin{equation}
\boxed{
 d_{Q,H}(j)
=
\frac1{N(H)}\mathbf w_j^\ast\mathbf y_H
-
\mathbf b_j^\ast G_H^{-1}\mathbf s_{Y,H}.
}
\label{eq:mixed_vector_entrywise}
\end{equation}
\end{proposition}

\begin{proof}
Insert \eqref{eq:PiH_exact_definition} in
\eqref{eq:mixed_background_prime_vector}.  The normalization
$B_H^\ast B_H=N(H)G_H$ gives the second term in
\eqref{eq:mixed_vector_entrywise}.
\end{proof}

\begin{theorem}[Exact decomposition of the numerical residual variance]
\label{thm:exact_system_variance}
For every finite cutoff $Q$,
\begin{equation}
\boxed{
V(Q,H)
=
V_Y(H)
-
2\Re\left(\mathbf c_Q^\ast\mathbf d_{Q,H}\right)
+
\mathbf c_Q^\ast\widetilde K_{Q,H}\mathbf c_Q.
}
\label{eq:exact_residual_variance}
\end{equation}
Grouping the coefficient vector, mixed vector, and Gram matrix by prime
base gives the exact identity
\begin{align}
V(Q,H)
={}&
V_Y(H)
-
2\Re\sum_p\mathbf c_p^\ast\mathbf d_{p,H}
\label{eq:tower_variance_background_mixed}
\\
&+
\sum_p
\mathbf c_p^\ast\widetilde K_{p,p;H}\mathbf c_p
+
\sum_{p\neq r}
\mathbf c_p^\ast\widetilde K_{p,r;H}\mathbf c_r.
\label{eq:tower_variance_decomposition}
\end{align}
\end{theorem}

\begin{proof}
By \eqref{eq:sampled_full_residual_vector} and
\eqref{eq:numerical_variance_exact_definition},
\begin{align}
V(Q,H)
&=
\frac1{N(H)}
(\mathbf y_H-V_Q\mathbf c_Q)^\ast
\Pi_H
(\mathbf y_H-V_Q\mathbf c_Q)
\label{eq:variance_expansion_start}
\\
&=
\frac1{N(H)}\mathbf y_H^\ast\Pi_H\mathbf y_H
-
\frac1{N(H)}\mathbf y_H^\ast\Pi_HV_Q\mathbf c_Q
\notag
\\
&\qquad
-
\frac1{N(H)}\mathbf c_Q^\ast V_Q^\ast\Pi_H\mathbf y_H
+
\frac1{N(H)}\mathbf c_Q^\ast V_Q^\ast\Pi_HV_Q\mathbf c_Q.
\label{eq:variance_expansion_four_terms}
\end{align}
The projection is Hermitian, so the two middle terms are complex
conjugates.  Definitions~\ref{def:background_mixed_vector} therefore turn
\eqref{eq:variance_expansion_four_terms} into
\eqref{eq:exact_residual_variance}.  Partitioning the finite index set by
prime base yields \eqref{eq:tower_variance_background_mixed}--
\eqref{eq:tower_variance_decomposition}.
\end{proof}

\begin{corollary}[Prime model alone]
\label{cor:prime_model_variance_only}
The identity
\begin{equation}
\frac1{N(H)}
(V_Q\mathbf c_Q)^\ast\Pi_H(V_Q\mathbf c_Q)
=
\mathbf c_Q^\ast\widetilde K_{Q,H}\mathbf c_Q
\label{eq:prime_model_projected_norm}
\end{equation}
holds for the projected prime model alone.  It equals the numerical
residual variance only when
\begin{equation}
V_Y(H)-2\Re(\mathbf c_Q^\ast\mathbf d_{Q,H})=0.
\label{eq:prime_only_equality_condition}
\end{equation}
\end{corollary}

\begin{proof}
Equation \eqref{eq:prime_model_projected_norm} follows directly from
\eqref{eq:projected_prime_gram_matrix}.  Comparison with
\eqref{eq:exact_residual_variance} gives the necessary and sufficient
condition \eqref{eq:prime_only_equality_condition}.
\end{proof}

\begin{warning}[What the diagonal question actually contains]
\label{warn:diagonal_issue}
The scalar Parseval sum
\begin{equation}
\sum_{p,r}
\left(|c_{p^r}^-|^2+|c_{p^r}^+|^2\right)
\label{eq:ordinary_parseval_diagonal}
\end{equation}
is neither the full numerical variance nor the diagonal of the complete
residual problem.  Before any asymptotic assertion about a ``diagonal'' is
made, one must control all four exact contributions in
\eqref{eq:tower_variance_background_mixed}--
\eqref{eq:tower_variance_decomposition}: the projected background energy,
the background--prime mixed term, the within-tower prime blocks, and the
cross-tower prime blocks.  The fact that the arguments inside a fixed
prime tower remain powers of that prime is exact; a Landau--Gonek
interpretation of the ordinate-only kernel is an additional analytic step.
\end{warning}

\begin{question}[Exact variance components]
\label{question:tower_block_energy}
Determine, uniformly for a growing cutoff, the separate sizes of
\begin{equation}
V_Y(H),
\qquad
-2\Re\sum_p\mathbf c_p^\ast\mathbf d_{p,H},
\qquad
\sum_p\mathbf c_p^\ast\widetilde K_{p,p;H}\mathbf c_p,
\qquad
\sum_{p\neq r}\mathbf c_p^\ast\widetilde K_{p,r;H}\mathbf c_r.
\label{eq:four_variance_components}
\end{equation}
Only their complete sum is the variance reported by the program.
\end{question}

\section{Finite logarithmic shells in Bohr mean}
\label{sec:shell_calculation}

\begin{definition}
\label{def:short_shell}
For fixed \(a>0\), define the finite multiplicative shell
\begin{equation}
\mathcal S_a(Q)
=
\{n=p^r:Q<n\leq e^aQ\}
\label{eq:short_shell}
\end{equation}
and its demodulated sum
\begin{equation}
G_{Q,a}(T)
=
\sum_{n\in\mathcal S_a(Q)}g_n(T).
\label{eq:shell_sum}
\end{equation}
\end{definition}

\begin{proposition}[Bohr energy of a fixed logarithmic shell]
\label{prop:shell_energy}
For fixed \(a>0\),
\begin{equation}
\lim_{X\to\infty}
\frac1X\int_0^X|G_{Q,a}(T)|^2\dd T
=
\frac{a}{2\pi^2\log Q}
+
o\left(\frac1{\log Q}\right)
\label{eq:shell_energy}
\end{equation}
as \(Q\to\infty\).
\end{proposition}

\begin{proof}
Distinct frequencies \(\pm\log n\) are orthogonal in Bohr mean.  For
primes,
\begin{equation}
|c_p^-|^2+|c_p^+|^2
=
\frac1{2\pi^2p}
+
O_\omega\left(\frac1{p(\log p)^2}\right).
\label{eq:ordinary_prime_energy}
\end{equation}
The prime number theorem and partial summation give \cite{Davenport2000,IwaniecKowalski2004}
\begin{equation}
\sum_{Q<p\leq e^aQ}\frac1p
=
\frac{a}{\log Q}
+
o\left(\frac1{\log Q}\right).
\label{eq:prime_shell_reciprocals}
\end{equation}
Proper prime powers contribute a lower-order term.
\end{proof}

\begin{corollary}
\label{cor:three_half_shell}
The choice \(a=3/2\) gives
\begin{equation}
\frac{3}{4\pi^2\log Q}
+
o\left(\frac1{\log Q}\right).
\label{eq:three_half_shell}
\end{equation}
\end{corollary}

\begin{proof}
Substitute \(a=3/2\) in Proposition~\ref{prop:shell_energy}.
\end{proof}

\begin{warning}
\label{warn:no_shell_localization}
Corollary~\ref{cor:three_half_shell} proves only the Bohr energy of a
deliberately selected hard shell.  It does not prove that the actual
residual is supported in \((Q,e^{3/2}Q]\).  Many transfer profiles have
the same quadratic mass, and the projected zero-sampled variance is not
the Bohr norm in Proposition~\ref{prop:shell_energy}.
\end{warning}

\section{Regularization and the remaining explicit-formula bridge}
\label{sec:regularization_bridge}

\begin{definition}
\label{def:regularized_window}
A regularized counting window is a family
\(H_{T,\eta}(u)\) tending to the midpoint indicator of \([0,T]\) as
\(\eta\to\infty\).  A formal Gaussian example is
\begin{equation}
H_{T,\eta}(u)
=
\int_0^T
\frac{\eta}{\sqrt\pi}
e^{-\eta^2(u-v)^2}\dd v.
\label{eq:gaussian_window}
\end{equation}
\end{definition}

\begin{heuristic}[Fourier displacement]
\label{heur:gaussian_fourier}
For
\begin{equation}
h_{T,\eta,\omega}(u)
=
e^{i\omega u}H_{T,\eta}(u),
\label{eq:gaussian_test}
\end{equation}
formal Fourier transformation gives
\begin{equation}
\widehat h_{T,\eta,\omega}(v)
=
e^{-(v-\omega)^2/(4\eta^2)}
\frac{1-e^{-i(v-\omega)T}}{i(v-\omega)}.
\label{eq:gaussian_fourier}
\end{equation}
Insertion into the explicit formula generates the two branches in
Definition~\ref{def:prime_power_coefficients}, multiplied by explicit
regularizing weights.
\end{heuristic}

\begin{warning}
\label{warn:regularization_requirements}
To turn Section~\ref{sec:regularization_bridge} into a theorem one must
specify a test class for which the explicit formula is valid, retain every
archimedean and boundary term, and control jointly:
\begin{equation}
H,\qquad Q,\qquad\eta.
\label{eq:joint_parameters}
\end{equation}
A formal Gaussian manipulation is not a proof of a sharp-window formula.
\end{warning}

\begin{conjecture}[Uniform explicit formula for the ordinate walk]
\label{conj:uniform_explicit_formula}
There exist choices \(Q=Q(H)\to\infty\) and
\(\eta=\eta(H)\to\infty\) for which the regularized prime-power formula,
including its complete archimedean and endpoint terms, converges locally
to the midpoint ordinate walk, with an error tending to zero in a
topology strong enough to identify the atomic measure
\eqref{eq:positive_measure_identity}.
\end{conjecture}

\begin{remark}
\label{rem:uniform_bridge_is_RH_interface}
The ordinary explicit formula reconstructs complete-zero weights, which
retain \(q^{\beta-1/2}\).  Replacing that object by the unweighted
ordinate walk is precisely where the positive defect
\(D_q(\tau)\) enters.  Truncation control alone does not remove the
horizontal zero information.
\end{remark}

\section{Energy exhaustion and distributional convergence}
\label{sec:two_main_lemmas}

\begin{conjecture}[Energy exhaustion]
\label{conj:energy_exhaustion}
For a suitable joint regime \(Q=Q(H)\to\infty\),
\begin{equation}
\Var_{\gamma\leq H}(F_{Q,H})
\leq
\varepsilon(Q,H),
\qquad
\varepsilon(Q(H),H)\to0.
\label{eq:energy_exhaustion}
\end{equation}
The variance is the full projected residual form
\eqref{eq:exact_residual_variance}, not the prime Gram term alone and not
an unprojected Lebesgue \(L^2\)-norm.
\end{conjecture}

\begin{theorem}[Abstract local distributional implication]
\label{thm:distributional_implication}
Suppose a sequence \(F_j\) satisfies
\begin{equation}
F_j-\frac12\longrightarrow0
\quad\text{in }L^2_{\mathrm{loc}}(\mathbb R),
\label{eq:local_L2_convergence}
\end{equation}
and
\begin{equation}
\frac{\dd}{\dd T}
\left[
e^{i\omega T}
\left(F_j(T)-\frac12\right)
\right]
\longrightarrow
\nu_q-\mu_q
\quad\text{in }\mathcal D'(\mathbb R).
\label{eq:derivative_distributional_identification}
\end{equation}
Then \(\RH\) holds.
\end{theorem}

\begin{proof}
Define
\begin{equation}
G_j(T)
=
e^{i\omega T}\left(F_j(T)-\frac12\right).
\label{eq:distributional_Gj_detailed}
\end{equation}
Multiplication by the bounded smooth phase preserves local \(L^2\)
convergence, so \(G_j\to0\) in \(L^2_{\mathrm{loc}}\).  On every compact
set, Cauchy--Schwarz then gives \(G_j\to0\) in \(L^1\).  For
\(\varphi\in C_c^\infty(\mathbb R)\),
\begin{equation}
\langle G_j',\varphi\rangle
=
-\int_{\mathbb R}G_j(T)\varphi'(T)\,\dd T
\longrightarrow0.
\label{eq:distributional_derivative_limit_detailed}
\end{equation}
The assumed identification of the derivatives therefore implies
\(\nu_q-\mu_q=0\) in \(\mathcal D'\).  A locally finite measure that is
zero as a distribution is the zero measure.  Multiplying by
\(\frac12e^{-iT\log q}\) and applying
Corollary~\ref{cor:atomic_positive_criterion} gives
\(D_q(\tau)=0\) at every ordinate, hence \(\RH\).
\end{proof}

\begin{warning}
\label{warn:global_mean_not_local}
A normalized mean-square limit on intervals \([0,H]\) does not imply
Theorem~\ref{thm:distributional_implication}.  A fixed compact defect, or
even a sparse family of atomic defects, is diluted by division by \(H\).
Local or distributional convergence is indispensable.
\end{warning}


\section{Further analytic consequences of the finite Möbius interpolation}
\label{sec:further_mobius_consequences}

\subsection{Meromorphic structure and residues}
\label{subsec:meromorphic_structure}

\begin{proposition}[Sources of poles]
\label{prop:sources_of_poles}
The meromorphic function $\Phi_N$ can have poles only at zeros of $\zeta$
and at zeros of the exponential polynomial $P_N$.  At a simple zero $\rho$
of $\zeta$ satisfying $P_N(\rho)\neq0$, the residue is
\begin{equation}
\operatorname*{Res}_{s=\rho}\Phi_N(s)
=
\frac{h_N}{\zeta'(\rho)P_N(\rho)}.
\label{eq:Phi_residue_zeta_zero}
\end{equation}
At a simple zero $s_0$ of $P_N$ that is not a zero of $\zeta$ and satisfies
$\zeta(s_0)^{-1}\neq A_N(s_0)$, the residue is
\begin{equation}
\operatorname*{Res}_{s=s_0}\Phi_N(s)
=
h_N\frac{\zeta(s_0)^{-1}-A_N(s_0)}{P_N'(s_0)}.
\label{eq:Phi_residue_P_zero}
\end{equation}
\end{proposition}

\begin{proof}
Both statements follow directly from the quotient in
\eqref{eq:PhiN_definition}.  At a simple zeta zero,
$\zeta(s)^{-1}$ has principal part
$1/(\zeta'(\rho)(s-\rho))$.  At a simple zero of $P_N$, the denominator has
principal part $P_N'(s_0)(s-s_0)$.
\end{proof}

\begin{proposition}[Value at the pole of zeta]
\label{prop:Phi_at_one}
The reciprocal zeta function has a simple zero at $s=1$, and therefore
\begin{equation}
\Phi_N(1)
=
B_N-h_N\frac{A_N(1)}{P_N(1)}.
\label{eq:Phi_at_one}
\end{equation}
\end{proposition}

\begin{proof}
Use $\zeta(s)^{-1}\to0$ as $s\to1$ in
\eqref{eq:PhiN_definition}.  The denominator $P_N(1)$ is positive.
\end{proof}

\subsection{Eventually exact sign and quantitative rounding}
\label{subsec:eventual_sign_rounding}

\begin{corollary}[Eventually exact sign]
\label{cor:eventual_sign}
For sufficiently large real $\sigma$,
\begin{equation}
\operatorname{sgn}\bigl(\Phi_N(\sigma)-M(N)\bigr)
=
\mu(a_N).
\label{eq:eventual_sign}
\end{equation}
\end{corollary}

\begin{proof}
The leading asymptotic has the form
\begin{equation}
\Phi_N(\sigma)-M(N)
=
A_N(\sigma)(1+o(1)),
\qquad
A_N(\sigma)
=
h_N\mu(a_N)\left(\frac{q_N}{a_N}\right)^\sigma.
\label{eq:eventual_sign_leading_term}
\end{equation}
The magnitude multiplying \(\mu(a_N)\) is strictly positive.  Choose
\(\sigma_0\) so that the relative error has modulus less than one half for
\(\sigma\geq\sigma_0\).  Then the full expression cannot change the sign
of its leading term, and its sign is \(\mu(a_N)\).
\end{proof}

\begin{proposition}[Existence of a finite rounding threshold]
\label{prop:finite_rounding_threshold}
For every $N\geq2$ there exists a finite $\sigma_N>1$ such that
\begin{equation}
\operatorname{nint}\Phi_N(\sigma)=M(N)
\label{eq:eventual_rounding}
\end{equation}
for all real $\sigma\geq\sigma_N$, where $\operatorname{nint}$ denotes the
nearest integer.
\end{proposition}

\begin{proof}
Let
\begin{equation}
B_N(\sigma)
=
\frac{h_NN}{\sigma-1}
\left(\frac{q_N}{N}\right)^\sigma.
\label{eq:rounding_threshold_bound_function}
\end{equation}
The base \(q_N/N\) lies strictly between zero and one.  Hence its
exponential power tends to zero, while \(h_NN/(\sigma-1)\) grows at most
polynomially and in fact decreases after the constant factor is removed.
Thus \(B_N(\sigma)\to0\).  Choose \(\sigma_N\) so that
\(B_N(\sigma)<1/2\) for every \(\sigma\geq\sigma_N\).  The rounding
criterion then gives the conclusion uniformly for all such \(\sigma\).
\end{proof}

\begin{proposition}[A sufficient explicit threshold]
\label{prop:explicit_rounding_threshold}
Let
\begin{equation}
\lambda_N=\log\frac{N}{q_N}>0.
\label{eq:lambdaN_definition}
\end{equation}
Every $\sigma>1$ satisfying
\begin{equation}
\sigma\lambda_N
>
\log(2h_NN)+\log\frac1{\sigma-1}
\label{eq:explicit_threshold_implicit}
\end{equation}
is a valid rounding point.  In particular, any $\sigma\geq2$ satisfying
\begin{equation}
\sigma
>
\frac{\log(2h_NN)}{\lambda_N}
\label{eq:explicit_threshold_simple}
\end{equation}
is sufficient.
\end{proposition}

\begin{proof}
Rewrite the rounding condition as
\begin{equation}
\frac{h_NN}{\sigma-1}e^{-\sigma\lambda_N}<\frac12.
\label{eq:explicit_threshold_exponential_form}
\end{equation}
All factors are positive.  Taking logarithms gives
\begin{equation}
\log(h_NN)-\log(\sigma-1)-\sigma\lambda_N<-
\log2,
\label{eq:explicit_threshold_log_form}
\end{equation}
which rearranges to the implicit condition in the statement.  When
\(\sigma\geq2\), one has \(\log(1/(\sigma-1))\leq0\).  The stronger but
simpler inequality
\(\sigma\lambda_N>\log(2h_NN)\) therefore implies the implicit condition
and is sufficient.
\end{proof}

\subsection{Recursive recovery of the nonzero tail}
\label{subsec:recursive_tail_recovery}

\begin{definition}
\label{def:tail_support_sequence}
Let
\begin{equation}
a_{N,1}<a_{N,2}<a_{N,3}<\cdots
\label{eq:tail_support_sequence}
\end{equation}
be the increasing sequence of integers exceeding $N$ for which
$\mu(a_{N,j})\neq0$.
\end{definition}

\begin{theorem}[Recursive exponential extraction]
\label{thm:recursive_extraction}
The real-axis function $E_N(\sigma)$ determines the entire signed sequence
\begin{equation}
\bigl(a_{N,j},\mu(a_{N,j})\bigr)_{j\geq1}
\label{eq:entire_tail_sequence}
\end{equation}
recursively.  More precisely, once the first $J-1$ terms are known, define
\begin{equation}
E_N^{(J)}(\sigma)
=
E_N(\sigma)-
\frac{\sum_{j<J}\mu(a_{N,j})a_{N,j}^{-\sigma}}{P_N(\sigma)}.
\label{eq:residual_tail_J}
\end{equation}
Then
\begin{align}
a_{N,J}
&=
q_N\lim_{\sigma\to\infty}
|E_N^{(J)}(\sigma)|^{-1/\sigma},
\label{eq:recursive_position}
\\
\mu(a_{N,J})
&=
\lim_{\sigma\to\infty}
\operatorname{sgn}E_N^{(J)}(\sigma).
\label{eq:recursive_sign}
\end{align}
\end{theorem}

\begin{proof}
By definition,
\begin{equation}
E_N(\sigma)
=
\frac{1}{P_N(\sigma)}
\sum_{j\geq1}\mu(a_{N,j})a_{N,j}^{-\sigma}.
\label{eq:recursive_full_expansion}
\end{equation}
Subtracting the first \(J-1\) known terms leaves
\begin{equation}
E_N^{(J)}(\sigma)
=
\frac{1}{P_N(\sigma)}
\sum_{j\geq J}\mu(a_{N,j})a_{N,j}^{-\sigma}.
\label{eq:recursive_residual_exact}
\end{equation}
Factor the first remaining term and the least denominator term exactly as
in Theorem~\ref{thm:leading_exponential}.  Dominated convergence gives
\begin{equation}
E_N^{(J)}(\sigma)
=
\mu(a_{N,J})
\left(\frac{q_N}{a_{N,J}}\right)^\sigma
(1+o(1)).
\label{eq:recursive_leading_asymptotic}
\end{equation}
Taking the negative \(\sigma\)-th root recovers the ratio
\(a_{N,J}/q_N\), and eventual sign recovers \(\mu(a_{N,J})\).  Induction
on \(J\) determines the entire signed support sequence uniquely.
\end{proof}

\begin{remark}
\label{rem:recursive_not_algorithmic_gain}
Theorem~\ref{thm:recursive_extraction} is an exact uniqueness statement.
It does not by itself yield a computational advantage, because evaluating
$E_N$ already requires the reciprocal zeta function with sufficient
precision and control of exponentially small differences.
\end{remark}

\section{Additional structure of the positive horizontal defect}
\label{sec:additional_positive_defect}

\subsection{Even-moment expansion}
\label{subsec:even_moment_expansion}

\begin{proposition}[Power-series expansion in $\log q$]
\label{prop:Dq_power_series}
At a fixed ordinate $\tau$, put
\begin{equation}
\delta_\rho=\beta-\frac12
\label{eq:delta_rho_definition}
\end{equation}
for zeros with $\beta>1/2$.  Then
\begin{equation}
D_q(\tau)
=
\sum_{m=1}^{\infty}
\frac{(\log q)^{2m}}{(2m)!}
\sum_{\substack{\rho=\beta+i\tau\\\beta>1/2}}
\delta_\rho^{2m}.
\label{eq:Dq_even_moments}
\end{equation}
\end{proposition}

\begin{proof}
At a fixed ordinate there are finitely many zeros, counted with finite
multiplicity.  Put \(y=\log q\) and
\(\delta_\rho=\beta-1/2>0\).  The entire expansion
\begin{equation}
\cosh(\delta_\rho y)-1
=
\sum_{m=1}^{\infty}
\frac{\delta_\rho^{2m}y^{2m}}{(2m)!}
\label{eq:Dq_cosh_series_proof}
\end{equation}
converges absolutely for every complex \(y\).  Because the outer zero sum
is finite at the fixed ordinate, it may be interchanged with the power
series.  Summing the coefficients gives the asserted even-moment
expansion, with no uniform-convergence subtlety.
\end{proof}

\begin{corollary}[Quadratic lower bound]
\label{cor:Dq_quadratic_lower}
For every $q>1$,
\begin{equation}
D_q(\tau)
\geq
\frac{(\log q)^2}{2}
\sum_{\substack{\rho=\beta+i\tau\\\beta>1/2}}
\left(\beta-\frac12\right)^2.
\label{eq:Dq_quadratic_lower}
\end{equation}
\end{corollary}

\begin{proof}
The Taylor series
\begin{equation}
\cosh x-1
=
\frac{x^2}{2!}+\frac{x^4}{4!}+\cdots
\label{eq:cosh_nonnegative_series}
\end{equation}
has nonnegative terms for real \(x\).  Dropping every term except the first
gives \(\cosh x-1\geq x^2/2\).  Apply this with
\(x=(\beta-1/2)\log q\) to every off-line pair at the fixed ordinate.
Since all summands remain nonnegative, summing preserves the inequality and
produces the stated quadratic moment.
\end{proof}

\begin{proposition}[Monotonicity and convexity]
\label{prop:Dq_monotonicity}
For fixed $\tau$, the function
\begin{equation}
t\longmapsto D_{e^t}(\tau),
\qquad t\geq0,
\label{eq:D_exp_t}
\end{equation}
is nonnegative, increasing, and convex.  It is identically zero if and only
if every zero at ordinate $\tau$ lies on the critical line.
\end{proposition}

\begin{proof}
Set \(t=\log q\geq0\).  For a horizontal displacement \(\delta>0\),
define \(f_\delta(t)=\cosh(\delta t)-1\).  Then
\begin{equation}
f_\delta'(t)=\delta\sinh(\delta t)\geq0,
\qquad
f_\delta''(t)=\delta^2\cosh(\delta t)>0.
\label{eq:Dq_monotonicity_derivatives_detailed}
\end{equation}
Thus every summand is increasing and strictly convex.  Their finite sum has
the same properties whenever at least one summand is present.  If no
off-line pair is present, the sum is identically zero.  Conversely, an
identically zero sum of nonnegative summands cannot contain a positive
summand, so no off-line pair exists.
\end{proof}

\begin{corollary}[One parameter and many parameters]
\label{cor:one_many_q}
One fixed $q>1$ suffices for the logical criterion
\eqref{eq:measure_RH_equivalence}.  A continuum of values of $q$ additionally
encodes all even power sums of the horizontal displacements at each
ordinate through \eqref{eq:Dq_even_moments}.
\end{corollary}

\begin{proof}
The fixed-parameter assertion is Corollary~\ref{cor:atomic_positive_criterion}.
For the family assertion, Proposition~\ref{prop:Dq_power_series} writes the
analytic function of \(y=\log q\) as
\begin{equation}
D_{e^y}(\tau)
=
\sum_{m\geq1}\frac{y^{2m}}{(2m)!}P_{2m}(\tau),
\label{eq:one_many_q_generating_function}
\end{equation}
where \(P_{2m}(\tau)\) is the \(2m\)-th power sum of the positive
horizontal displacements.  Differentiating \(2m\) times at zero gives
\(P_{2m}(\tau)\).  Thus a neighborhood of one in the \(q\)-parameter
encodes all even moments, whereas positivity makes a single nontrivial
parameter enough merely to test whether they all vanish.
\end{proof}

\subsection{Jump formulation}
\label{subsec:jump_formulation}

\begin{definition}
\label{def:cumulative_complete_walks}
Define the right-continuous cumulative functions
\begin{align}
U_q(T)
&=
\sum_{0<\tau\leq T}m(\tau)e^{i\tau\log q},
\label{eq:Uq_cumulative}
\\
V_q(T)
&=
\sum_{0<\tau\leq T}W_q(\tau)e^{i\tau\log q}.
\label{eq:Vq_cumulative}
\end{align}
\end{definition}

\begin{proposition}[Exact jump defect]
\label{prop:jump_defect}
At every ordinate $\tau$,
\begin{equation}
\boxed{
\frac12e^{-i\tau\log q}
\Delta(V_q-U_q)(\tau)
=
D_q(\tau).
}
\label{eq:jump_defect}
\end{equation}
\end{proposition}

\begin{proof}
The jump of the complete-zero walk at \(\tau\) is
\begin{equation}
\Delta V_q(\tau)
=
W_q(\tau)e^{i\tau\log q},
\label{eq:jump_complete_walk}
\end{equation}
whereas the jump of the ordinate-only walk is
\begin{equation}
\Delta U_q(\tau)
=
m(\tau)e^{i\tau\log q}.
\label{eq:jump_ordinate_walk}
\end{equation}
Subtracting and using \(W_q-m=2D_q\) gives
\begin{equation}
\Delta(V_q-U_q)(\tau)
=
2D_q(\tau)e^{i\tau\log q}.
\label{eq:jump_difference_detailed}
\end{equation}
Demodulation by \(\frac12e^{-i\tau\log q}\) leaves the nonnegative scalar
\(D_q(\tau)\).
\end{proof}

\section{Regularized explicit formula with explicit prime-power weights}
\label{sec:regularized_explicit_formula_detailed}

\subsection{The explicit formula in the Fourier normalization used here}
\label{subsec:guinand_weil_normalization}

\begin{theorem}[Guinand--Weil explicit formula]
\label{thm:guinand_weil_normalized}
Let \(h\) be holomorphic in a strip
\(|\Im z|\leq1/2+\epsilon\), real-valued on \(\mathbb R\), and suppose
that, uniformly in the strip,
\(h(z)\ll(1+|\Re z|)^{-1-\delta}\) for some \(\delta>0\).  Define
\begin{equation}
\widehat h(v)
=
\int_{\mathbb R}h(t)e^{-ivt}\,\dd t.
\label{eq:guinand_fourier_convention}
\end{equation}
Then, with zeros counted with multiplicity,
\begin{align}
\sum_\rho
h\!\left(\frac{\rho-1/2}{i}\right)
={}&
h\!\left(\frac1{2i}\right)
+
h\!\left(-\frac1{2i}\right)
-
\frac{\widehat h(0)}{2\pi}\log\pi
\label{eq:guinand_archimedean_first}
\\
&+
\frac1{2\pi}
\int_{-\infty}^{\infty}
 h(u)
 \Re\frac{\Gamma'}{\Gamma}
 \!\left(\frac14+\frac{iu}{2}\right)
\,\dd u
\label{eq:guinand_archimedean_second}
\\
&-
\frac1{2\pi}
\sum_{n=2}^{\infty}
\frac{\Lambda(n)}{\sqrt n}
\left[
\widehat h(\log n)+\widehat h(-\log n)
\right].
\label{eq:guinand_prime_side}
\end{align}
The formula is understood by the usual limiting procedure whenever the
chosen test lies at the boundary of the stated class
\cite{Guinand1948,Weil1952}.
\end{theorem}

\begin{proof}
Consider the completed function \(\xi(s)\) from
\eqref{eq:rvm_xi_definition} and the contour integral
\begin{equation}
\frac1{2\pi i}
\int_{\partial\mathcal R}
\frac{\xi'}{\xi}(s)
 h\!\left(\frac{s-1/2}{i}\right)\,\dd s
\label{eq:guinand_contour_integral}
\end{equation}
over a rectangle whose vertical sides are \(\Re s=1/2\pm a\), with
\(a>1/2\), and whose horizontal sides tend to infinite height through
values avoiding zeros.  The decay assumptions on \(h\) make the horizontal
integrals vanish.  By the residue theorem, the enclosed zeros of \(\xi\)
contribute the spectral sum on the left, while the factors \(s(s-1)\)
produce the two evaluations at \(\pm1/(2i)\).

On the right vertical side, decompose the logarithmic derivative as
\begin{equation}
\frac{\xi'}{\xi}(s)
=
\frac1s+\frac1{s-1}
-
\frac12\log\pi
+
\frac12\frac{\Gamma'}{\Gamma}\!\left(\frac s2\right)
+
\frac{\zeta'}{\zeta}(s).
\label{eq:guinand_log_derivative_decomposition}
\end{equation}
The functional equation \(\xi(s)=\xi(1-s)\) transforms the left vertical
side into the conjugate contribution of the right side.  After the two
sides are combined and parametrized by \(s=1/2+iu\), the constant
\(-\frac12\log\pi\) gives
\(-\widehat h(0)\log\pi/(2\pi)\), and the gamma logarithmic derivative
gives the integral in \eqref{eq:guinand_archimedean_second}.

For \(\Re s>1\), the absolutely convergent Dirichlet series
\begin{equation}
-\frac{\zeta'}{\zeta}(s)
=
\sum_{n=2}^{\infty}\frac{\Lambda(n)}{n^s}
\label{eq:guinand_log_zeta_series}
\end{equation}
permits termwise integration.  The contribution of \(n^{-s}\) on the two
paired vertical sides is
\begin{equation}
\frac{\Lambda(n)}{\sqrt n}
\left[
\widehat h(\log n)+\widehat h(-\log n)
\right],
\label{eq:guinand_fourier_prime_pair}
\end{equation}
with the prefactor \(-1/(2\pi)\) arising from the contour normalization and
the parametrization \(\dd s=i\,\dd u\).  Summing the contributions gives
\eqref{eq:guinand_prime_side}.  Absolute convergence in the initial test
class justifies all interchanges; the standard limiting argument extends
the identity to the stated boundary cases.
\end{proof}

\begin{corollary}[Complex-linear extension]
\label{cor:guinand_complex_extension}
The formula in Theorem~\ref{thm:guinand_weil_normalized} remains valid
for complex-valued \(h\) satisfying the same analytic and decay
conditions.
\end{corollary}

\begin{proof}
Define
\begin{align}
h_1(z)
&=
\frac12\left(h(z)+\overline{h(\overline z)}\right),
\label{eq:complex_extension_h1}
\\
h_2(z)
&=
\frac1{2i}\left(h(z)-\overline{h(\overline z)}\right).
\label{eq:complex_extension_h2}
\end{align}
Both functions are holomorphic in the strip, real-valued on the real
axis, satisfy the same decay, and obey \(h=h_1+ih_2\).  Apply the real
formula separately to \(h_1\) and \(h_2\), then combine by complex
linearity.
\end{proof}

\begin{corollary}[Why only prime powers occur]
\label{cor:explicit_formula_prime_power_support}
The arithmetic side of the formula is supported exactly on prime powers,
because
\begin{equation}
\Lambda(n)\neq0
\quad\Longleftrightarrow\quad
n=p^r
\label{eq:von_mangoldt_support}
\end{equation}
for a prime \(p\) and integer \(r\geq1\).
\end{corollary}

\begin{proof}
This is the defining support property of the von Mangoldt function:
\(\Lambda(p^r)=\log p\) and \(\Lambda(n)=0\) otherwise.  Applying it to
the prime sum \eqref{eq:guinand_prime_side} removes every non-prime-power
integer.
\end{proof}

\subsection{Gaussian convolution of a finite interval}
\label{subsec:gaussian_convolution}

\begin{definition}
\label{def:gaussian_kernel_detailed}
For $\eta>0$, define
\begin{equation}
g_\eta(t)
=
\frac{\eta}{\sqrt\pi}e^{-\eta^2t^2},
\label{eq:gaussian_kernel_detailed}
\end{equation}
and
\begin{equation}
H_{T,\eta}(z)
=
\int_0^Tg_\eta(z-u)\dd u.
\label{eq:H_T_eta_detailed}
\end{equation}
The oscillatory test is
\begin{equation}
h_{T,\eta,\omega}(z)
=
e^{i\omega z}H_{T,\eta}(z).
\label{eq:h_T_eta_omega_detailed}
\end{equation}
\end{definition}

\begin{proposition}[Fourier transform]
\label{prop:gaussian_fourier_transform}
With the convention
\begin{equation}
\widehat f(v)=\int_{\mathbb R}f(t)e^{-ivt}\dd t,
\label{eq:fourier_convention}
\end{equation}
one has
\begin{equation}
\boxed{
\widehat h_{T,\eta,\omega}(v)
=
e^{-(v-\omega)^2/(4\eta^2)}
\frac{1-e^{-i(v-\omega)T}}{i(v-\omega)},
}
\label{eq:exact_gaussian_fourier_transform}
\end{equation}
where the quotient is interpreted continuously as $T$ at $v=\omega$.
\end{proposition}

\begin{proof}
With the Fourier convention
\(\widehat f(v)=\int_{\mathbb R}f(t)e^{-ivt}\,\dd t\), the normalized
Gaussian has transform
\begin{equation}
\widehat g_\eta(v)=e^{-v^2/(4\eta^2)}.
\label{eq:gaussian_transform_proof}
\end{equation}
The regularized interval is the convolution
\(H_{T,\eta}=\1_{[0,T]}*g_\eta\), so the convolution theorem gives
\begin{equation}
\widehat H_{T,\eta}(v)
=
e^{-v^2/(4\eta^2)}
\frac{1-e^{-ivT}}{iv},
\label{eq:interval_gaussian_transform_proof}
\end{equation}
with the quotient interpreted continuously at \(v=0\).  Multiplication by
\(e^{i\omega t}\) translates the Fourier transform:
\(\widehat{e^{i\omega t}H}(v)=\widehat H(v-\omega)\).  Substitution proves
the formula.
\end{proof}

\begin{proposition}[Regularized two-branch prime contribution]
\label{prop:regularized_two_branches}
In the prime-power side of the Guinand--Weil explicit formula, the test
\eqref{eq:h_T_eta_omega_detailed} produces, apart from terms independent of
$T$, the oscillatory sum
\begin{align}
\cP_{\eta,\omega}^{\mathrm{osc}}(T)
={}&
\sum_{n=p^r}
c_n^-e^{-(\ell_n-\omega)^2/(4\eta^2)}
e^{i(\omega-\ell_n)T}
\label{eq:regularized_minus_branch}
\\
&+
\sum_{n=p^r}
c_n^+e^{-(\ell_n+\omega)^2/(4\eta^2)}
e^{i(\omega+\ell_n)T}.
\label{eq:regularized_plus_branch}
\end{align}
\end{proposition}

\begin{proof}
By Corollary~\ref{cor:guinand_complex_extension}, the explicit formula
applies to the complex oscillatory test.  Insert \(v=\ell_n\) and
\(v=-\ell_n\) into the exact Fourier transform.
For the first argument,
\begin{equation}
\widehat h(\ell_n)
=
w_{n,-}
\frac{1-e^{-i(\ell_n-\omega)T}}
{i(\ell_n-\omega)},
\label{eq:branch_minus_fourier_insert}
\end{equation}
where \(w_{n,-}=e^{-(\ell_n-\omega)^2/(4\eta^2)}\).  Multiplication by the
prime-side factor \(-\Lambda(n)/(2\pi\sqrt n)\) separates a constant and
the oscillatory term
\begin{equation}
w_{n,-}c_n^-e^{i(\omega-\ell_n)T}.
\label{eq:branch_minus_oscillatory_proof}
\end{equation}
The argument \(-\ell_n\) similarly gives
\begin{equation}
w_{n,+}c_n^+e^{i(\omega+\ell_n)T},
\qquad
w_{n,+}=e^{-(\ell_n+\omega)^2/(4\eta^2)}.
\label{eq:branch_plus_oscillatory_proof}
\end{equation}
The remaining terms are independent of \(T\) and belong to the additive
constant that must be retained or fitted before demodulation.
\end{proof}

\subsection{A uniform bound for the regularized prime tail}
\label{subsec:regularized_tail_bound}

\begin{theorem}[Uniform Gaussian tail]
\label{thm:uniform_gaussian_tail}
Assume $Q>e^{2\omega}$ and put $L=\log Q$.  If
$L>\omega+\eta^2$, then
\begin{align}
&\sup_{T\in\mathbb R}
\left|
\sum_{\substack{n=p^r\\n>Q}}
\left[
 c_n^-e^{-(\ell_n-\omega)^2/(4\eta^2)}e^{i(\omega-\ell_n)T}
+c_n^+e^{-(\ell_n+\omega)^2/(4\eta^2)}e^{i(\omega+\ell_n)T}
\right]
\right|
\notag
\\
&\qquad\ll_\omega
\frac{\eta^2}{L-\omega-\eta^2}
\exp\left(
\frac{\eta^2}{4}
-
\frac{(L-\omega-\eta^2)^2}{4\eta^2}
\right).
\label{eq:uniform_gaussian_tail}
\end{align}
\end{theorem}

\begin{proof}
For \(n>Q>e^{2\omega}\),
\begin{equation}
|\ell_n-\omega|\geq\frac12\ell_n,
\qquad
\ell_n+\omega\geq\ell_n.
\label{eq:tail_denominator_bounds_detailed}
\end{equation}
Using \(|1-e^{ix}|\leq2\) and \(\Lambda(n)\leq\log n\), both branches of
the omitted prime side are bounded, up to a constant depending on
\(\omega\), by
\begin{equation}
\sum_{n>Q}n^{-1/2}
\exp\left[-\frac{(\log n-\omega)^2}{4\eta^2}\right].
\label{eq:tail_integer_sum_detailed}
\end{equation}
When \(L=\log Q>\omega+\eta^2\), the corresponding continuous integrand is
decreasing beyond \(Q\), so the sum is bounded by a constant multiple of
\begin{equation}
I
=
\int_Q^\infty x^{-1/2}
\exp\left[-\frac{(\log x-\omega)^2}{4\eta^2}\right]\,\dd x.
\label{eq:tail_integral_detailed}
\end{equation}
Put \(u=\log x\).  Completing the square gives
\begin{align}
I
&=
\int_L^\infty
\exp\left[
\frac u2-
\frac{(u-\omega)^2}{4\eta^2}
\right]\,\dd u
\label{eq:tail_square_one}
\\
&=
\exp\left(\frac\omega2+\frac{\eta^2}{4}\right)
\int_L^\infty
\exp\left[-
\frac{(u-\omega-\eta^2)^2}{4\eta^2}
\right]\,\dd u.
\label{eq:tail_square_two}
\end{align}
For \(r=L-\omega-\eta^2>0\), the elementary Gaussian-tail estimate
\begin{equation}
\int_r^\infty e^{-v^2/(4\eta^2)}\,\dd v
\leq
\frac{2\eta^2}{r}e^{-r^2/(4\eta^2)}
\label{eq:gaussian_tail_elementary}
\end{equation}
follows by comparing \(1\leq v/r\) and integrating the derivative of the
Gaussian.  Substitution yields the asserted uniform bound.  The estimate is
independent of \(T\) because only \(|1-e^{ix}|\leq2\) was used.
\end{proof}

\begin{corollary}[One explicit joint choice]
\label{cor:joint_Q_eta_choice}
If
\begin{equation}
\log Q=\omega+3\eta^2,
\label{eq:joint_Q_eta_choice}
\end{equation}
then the right-hand side of \eqref{eq:uniform_gaussian_tail} is
\begin{equation}
O_\omega(e^{-3\eta^2/4}).
\label{eq:joint_Q_eta_decay}
\end{equation}
\end{corollary}

\begin{proof}
Under \eqref{eq:joint_Q_eta_choice},
\begin{equation}
L-\omega-\eta^2=2\eta^2.
\label{eq:joint_choice_r}
\end{equation}
Substitution in \eqref{eq:uniform_gaussian_tail} gives a bounded
prefactor and exponent \(\eta^2/4-\eta^2=-3\eta^2/4\).
\end{proof}

\begin{warning}
\label{warn:regularized_tail_not_ordinate_formula}
Theorem~\ref{thm:uniform_gaussian_tail} controls the prime side of a
regularized explicit formula.  It does not identify the complete-zero
spectral side with the unweighted ordinate walk.  That identification is
exactly where the positive measure \eqref{eq:positive_measure_identity}
appears.
\end{warning}

\section{Mean-square statements and their exact logical strength}
\label{sec:mean_square_logic}

\subsection{Finite Dirichlet-polynomial tails in Lebesgue mean}
\label{subsec:finite_dirichlet_mean}

\begin{proposition}[Finite-tail mean-value estimate]
\label{prop:finite_tail_mean_value}
Let
\begin{equation}
E_{Q,R}(T)
=
\sum_{\substack{Q<n\leq R\\n=p^r}}
\left(c_n^-n^{-iT}+c_n^+n^{iT}\right).
\label{eq:finite_tail_ER}
\end{equation}
Then
\begin{align}
\frac1X\int_0^X|E_{Q,R}(T)|^2\dd T
\ll_\omega
\left(1+\frac RX\right)
\left[
\log\frac{\log R}{\log Q}
+
\frac1{\log Q}
+
Q^{-1/2}
\right]
\label{eq:finite_tail_mean_value}
\end{align}
for $e^{2\omega}<Q<R$.
\end{proposition}

\begin{proof}
Apply the Montgomery--Vaughan mean-value theorem for Dirichlet polynomials
\cite[Chapter~7]{MontgomeryVaughan2007} separately to
\begin{equation}
\sum_{Q<n\leq R}c_n^-n^{-iT}
\quad\text{and}\quad
\sum_{Q<n\leq R}c_n^+n^{iT}.
\label{eq:MV_two_polynomials}
\end{equation}
It gives an upper bound of the form
\begin{equation}
\int_0^X\left|\sum_{n\leq R}a_nn^{-iT}\right|^2\,\dd T
\ll
(X+R)\sum_{n\leq R}|a_n|^2.
\label{eq:MV_mean_value_used}
\end{equation}
The inequality \(|u+v|^2\leq2|u|^2+2|v|^2\) combines the branches.  For
\(n=p^r>e^{2\omega}\),
\begin{equation}
|c_n^-|^2+|c_n^+|^2
\ll_\omega
\frac{\Lambda(n)^2}{n(\log n)^2}
=
\frac1{r^2p^r}.
\label{eq:finite_tail_coeff_square}
\end{equation}
The prime terms sum to
\(\log(\log R/\log Q)+O(1/\log Q)\) by Mertens' theorem for primes, while
proper powers contribute \(O(Q^{-1/2})\).  Division by \(X\) yields the
stated estimate.
\end{proof}

\begin{corollary}[A finite logarithmic tail can be exhausted]
\label{cor:finite_log_tail_exhaustion}
Let
\begin{equation}
R(X)=X^{1/2},
\qquad
Q(X)=R(X)^{1-1/\sqrt{\log\log X}}.
\label{eq:finite_tail_choice}
\end{equation}
Then
\begin{equation}
\frac1X\int_0^X|E_{Q(X),R(X)}(T)|^2\dd T
\longrightarrow0.
\label{eq:finite_tail_exhaustion}
\end{equation}
\end{corollary}

\begin{proof}
For the choice \eqref{eq:finite_tail_choice},
\begin{equation}
\frac{R(X)}{X}=X^{-1/2}\longrightarrow0.
\label{eq:finite_tail_RX}
\end{equation}
If \(\delta_X=1/\sqrt{\log\log X}\), then
\(\log Q=(1-\delta_X)\log R\), so
\begin{equation}
\log\frac{\log R}{\log Q}
=
-\log(1-\delta_X)
=
O(\delta_X)\longrightarrow0.
\label{eq:finite_tail_log_ratio}
\end{equation}
Also \(1/\log Q\to0\) and \(Q^{-1/2}\to0\).  Apply
Proposition~\ref{prop:finite_tail_mean_value}.
\end{proof}

\begin{warning}
\label{warn:finite_tail_not_full_residual}
Corollary~\ref{cor:finite_log_tail_exhaustion} concerns a finite tail in
Lebesgue measure.  The numerical variance is sampled on zeta ordinates and
contains the projected prime kernel \eqref{eq:projected_prime_gram_matrix}; the two norms are
not interchangeable.
\end{warning}

\subsection{Normalized global convergence does not detect fixed atoms}
\label{subsec:global_not_local}

\begin{proposition}[Counterexample to global normalized sufficiency]
\label{prop:global_normalized_counterexample}
There exist functions $F_X$ such that
\begin{equation}
\frac1X\int_0^X
\left|F_X(T)-\frac12\right|^2\dd T
\longrightarrow0,
\label{eq:global_normalized_convergence}
\end{equation}
but whose distributional derivatives retain a nonzero fixed atomic part.
\end{proposition}

\begin{proof}
Take \(F_X(T)=1/2+\1_{[0,1]}(T)\).  Then for every \(X\geq1\),
\begin{equation}
\frac1X\int_0^X\left|F_X(T)-\frac12\right|^2\,\dd T
=
\frac1X\longrightarrow0.
\label{eq:global_counterexample_mean}
\end{equation}
However, on the fixed compact interval \([0,1]\), the defect remains
identically one.  In the sense of distributions,
\begin{equation}
\frac{\dd}{\dd T}\1_{[0,1]}
=
\delta_0-\delta_1,
\label{eq:global_counterexample_derivative}
\end{equation}
which does not tend to zero.  Thus normalization by the length of an
expanding interval can erase a fixed local or atomic obstruction.
\end{proof}

\begin{corollary}
\label{cor:local_topology_required}
Any route from energy exhaustion to the vanishing of the positive atomic
measure \eqref{eq:positive_measure_identity} must use a local topology,
atomwise convergence, or distributional convergence against compactly
supported test functions.
\end{corollary}

\begin{proof}
Proposition~\ref{prop:global_normalized_counterexample} shows that global
normalized convergence alone permits a persistent compact atomic defect.
Therefore any sufficient argument must contain additional local
information of one of the stated forms.
\end{proof}

\section{Weighted prime-power profiles}
\label{sec:weighted_profiles}

\begin{definition}
\label{def:weighted_profile}
Let \(W\in C_c([0,\infty))\).  Define
\begin{equation}
G_{Q,W}(T)
=
\sum_{n=p^r}
W\left(\log\frac nQ\right)g_n(T).
\label{eq:weighted_profile_sum}
\end{equation}
\end{definition}

\begin{theorem}[Weighted Bohr energy]
\label{thm:weighted_bohr_energy}
As $Q\to\infty$,
\begin{equation}
\lim_{X\to\infty}\frac1X\int_0^X|G_{Q,W}(T)|^2\dd T
=
\frac1{2\pi^2\log Q}
\int_0^\infty|W(u)|^2\dd u
+
o\left(\frac1{\log Q}\right).
\label{eq:weighted_bohr_energy}
\end{equation}
\end{theorem}

\begin{proof}
Expand the squared modulus of the finite weighted sum.  Every diagonal term
contributes
\begin{equation}
|W(\log(n/Q))|^2
\left(|c_n^-|^2+|c_n^+|^2\right).
\label{eq:weighted_bohr_diagonal}
\end{equation}
Every unequal pair has a nonzero frequency difference and therefore zero
Bohr mean.  For primes, use
\begin{equation}
|c_p^-|^2+|c_p^+|^2
=
\frac1{2\pi^2p}
+
O_\omega\left(\frac1{p(\log p)^2}\right).
\label{eq:weighted_bohr_prime_coeff}
\end{equation}
Write \(p=Qe^u\) and apply the prime number theorem by partial summation.
For compactly supported continuous \(W\),
\begin{equation}
\sum_p\frac1p
\left|W\!\left(\log\frac pQ\right)\right|^2
=
\frac1{\log Q}
\int_0^\infty|W(u)|^2\,\dd u
+
o\left(\frac1{\log Q}\right).
\label{eq:weighted_bohr_pnt}
\end{equation}
The coefficient error and all proper prime powers are lower order.  This
proves the formula.
\end{proof}

\begin{corollary}[Interpretation of the constant $3/4$]
\label{cor:three_quarter_profile_mass}
A weighted Bohr model has leading energy
$3/(4\pi^2\log Q)$ precisely when
\begin{equation}
\int_0^\infty|W(u)|^2\dd u
=
\frac32.
\label{eq:profile_mass_three_half}
\end{equation}
The constant determines a quadratic profile mass, not a unique hard support.
\end{corollary}

\begin{proof}
Compare the leading term in
Theorem~\ref{thm:weighted_bohr_energy} with
\(3/(4\pi^2\log Q)\).  Equality of the coefficients is equivalent to
\(\int_0^\infty|W(u)|^2\,\dd u=3/2\).
\end{proof}

\section{Detailed projected Gram algebra}
\label{sec:detailed_gram_algebra}

\subsection{Four branch expansion}
\label{subsec:four_branch_expansion}

\begin{proposition}[Full unprojected quadratic expansion]
\label{prop:full_four_branch_expansion}
For a finite set $\mathcal N$ of prime powers, put
\begin{equation}
G_{\mathcal N}(\gamma)
=
\sum_{n\in\mathcal N}
\left(c_n^-n^{-i\gamma}+c_n^+n^{i\gamma}\right).
\label{eq:G_N_set}
\end{equation}
Then
\begin{align}
\|G_{\mathcal N}\|_H^2
={}&
\sum_{m,n\in\mathcal N}
\overline{c_m^-}c_n^-M_H(m/n)
\label{eq:four_branch_mm}
\\
&+
\sum_{m,n\in\mathcal N}
\overline{c_m^-}c_n^+M_H(mn)
\label{eq:four_branch_mp}
\\
&+
\sum_{m,n\in\mathcal N}
\overline{c_m^+}c_n^-\overline{M_H(mn)}
\label{eq:four_branch_pm}
\\
&+
\sum_{m,n\in\mathcal N}
\overline{c_m^+}c_n^+M_H(n/m).
\label{eq:four_branch_pp}
\end{align}
\end{proposition}

\begin{proof}
Write the finite model as
\begin{equation}
G(\gamma)
=
\sum_n\left(c_n^-n^{-i\gamma}+c_n^+n^{i\gamma}\right).
\label{eq:four_branch_G_definition}
\end{equation}
Multiplying by its conjugate produces four sign choices.  For example,
\begin{equation}
\overline{c_m^-m^{-i\gamma}}\,c_n^-n^{-i\gamma}
=
\overline{c_m^-}c_n^-
\left(\frac mn\right)^{i\gamma},
\label{eq:four_branch_minus_minus}
\end{equation}
whose sample average is
\(\overline{c_m^-}c_n^-M_H(m/n)\).  The other three products give the
arguments \(mn\), \((mn)^{-1}\), and \(n/m\).  Summing all pairs yields
the displayed expansion exactly.
\end{proof}

\subsection{Prime-tower blocks}
\label{subsec:prime_tower_blocks}

\begin{proposition}[Same-tower multiplicative closure]
\label{prop:same_tower_resonance}
For $m=p^r$ and $n=p^s$, every argument of $M_H$ in
\eqref{eq:four_branch_mm}--\eqref{eq:four_branch_pp} is one of
\begin{equation}
p^{s-r},
\qquad
p^{r-s},
\qquad
p^{r+s},
\qquad
p^{-(r+s)}.
\label{eq:same_tower_arguments}
\end{equation}
Thus the exact Gram arguments of a two-branch prime tower remain inside
the multiplicative set generated by \(p\); the block is not reduced to the
scalar diagonal \(r=s\).
\end{proposition}

\begin{proof}
Substitute \(m=p^r\) and \(n=p^s\) into the four multiplicative arguments.
They become
\begin{equation}
p^{r-s},
\qquad
p^{r+s},
\qquad
p^{-(r+s)},
\qquad
p^{s-r}.
\label{eq:same_tower_arguments_proof}
\end{equation}
Each is one, a prime power, or the reciprocal of a prime power.  This is
an exact algebraic closure statement.  Interpreting these entries by the
main term of a Landau--Gonek formula for the ordinate-only kernel
\(M_H\) additionally requires the Riemann Hypothesis, or an independent
replacement of \(x^{i\gamma}\) by \(x^{\rho-1/2}\).
\end{proof}

\begin{proposition}[Projection correction]
\label{prop:projection_correction_detailed}
Let $V$ be the matrix whose columns are the sampled prime-power modes, and
let $B$ be the two-column matrix with columns $1$ and $e^{i\omega\gamma}$.
Then the residual Gram matrix is
\begin{equation}
\widetilde K_H
=
\frac1{N(H)}V^\ast V
-
\frac1{N(H)^2}
V^\ast B
\left(\frac1{N(H)}B^\ast B\right)^{-1}
B^\ast V.
\label{eq:matrix_schur_projection}
\end{equation}
In particular, even the norm of a single frequency is altered by the
constant--carrier projection.
\end{proposition}

\begin{proof}
Let \(B\) be the matrix whose columns are the sampled constant and carrier
vectors.  Since these columns are linearly independent when
\(|M_H(e^\omega)|<1\), the orthogonal projection onto their span is
\begin{equation}
P_B=B(B^\ast B)^{-1}B^\ast.
\label{eq:projection_PB_detailed}
\end{equation}
The complementary projection is \(I-P_B\).  If \(V\) contains the sampled
prime frequency vectors, the projected Gram matrix is
\begin{align}
\frac1{N(H)}V^\ast(I-P_B)V
={}&
\frac1{N(H)}V^\ast V
\label{eq:projection_gram_one}
\\
&-
\frac1{N(H)}V^\ast B(B^\ast B)^{-1}B^\ast V.
\label{eq:projection_gram_two}
\end{align}
Writing \(B^\ast B=N(H)G_H\) gives the normalized Schur-complement form in
the proposition.  In particular, projection changes both diagonal and
off-diagonal entries.
\end{proof}

\section{Analytic targets isolated by the exact identities}
\label{sec:analytic_targets}

\begin{conjecture}[Collective carrier law]
\label{conj:collective_carrier_law}
In a specified joint range $Q=Q(H)\to\infty$,
\begin{equation}
\Im\widehat c_{Q,H}
=
\pi\Delta_\omega(Q)
+
o\left(\Delta_\omega(Q)\right).
\label{eq:collective_carrier_law}
\end{equation}
By Corollary~\ref{cor:transport_target}, this is equivalent to a weighted
collective cancellation in the exact kernel $M_H(n)$.
\end{conjecture}

\begin{conjecture}[Projected variance law]
\label{conj:projected_variance_law}
For a suitable joint range,
\begin{equation}
V(Q,H)
=
\frac34\Delta_\omega(Q)
+
o\left(\Delta_\omega(Q)\right).
\label{eq:projected_variance_law}
\end{equation}
The left-hand side is the full projected residual expression
\eqref{eq:exact_residual_variance}.  Its prime Gram term is only one of
three cutoff-dependent contributions in addition to the projected
background energy.
\end{conjecture}

\begin{question}[Background, mixed, tower, and cross-tower separation]
\label{question:tower_cross_tower_separation}
Can one prove, uniformly in a growing cutoff, asymptotic formulas for
\begin{align}
-2\Re\sum_p\mathbf c_p^\ast\mathbf d_{p,H},
\label{eq:mixed_sum_target}
\\
\sum_p\mathbf c_p^\ast\widetilde K_{p,p;H}\mathbf c_p,
\label{eq:tower_sum_target}
\\
\sum_{p\neq r}\mathbf c_p^\ast\widetilde K_{p,r;H}\mathbf c_r
\label{eq:cross_tower_sum_target}
\end{align}
relative to the projected background energy $V_Y(H)$?  These are the exact
components that must be separated before any assertion about a
``diagonal'' can be made.
\end{question}




\section{Exact routes from the exponential zero walk to the Riemann Hypothesis}
\label{sec:routes_to_RH}

This section records the routes that arise directly from the phase-walk
construction.  Each route is an exact implication or equivalence.  What is
missing in every incomplete route is stated separately; no numerical
observation is used as a premise.

\subsection{Route I: equality of the complete and ordinate-only atomic measures}
\label{subsec:route_atomic_equality}

\begin{theorem}[Fixed-\(q\) atomic criterion]
\label{thm:route_fixed_q_atomic}
For any fixed \(q>1\),
\begin{equation}
\nu_q=\mu_q
\quad\Longleftrightarrow\quad
\RH.
\label{eq:route_fixed_q_atomic}
\end{equation}
Equivalently, it is enough to prove equality of the jumps of the two walks
at every positive ordinate.
\end{theorem}

\begin{proof}
The measure difference is given exactly by
\begin{equation}
\frac12e^{-iT\log q}(\nu_q-\mu_q)
=
\sum_{\tau>0}D_q(\tau)\delta_\tau.
\label{eq:route_atomic_positive_measure}
\end{equation}
Every coefficient on the right is nonnegative.  Equality of the measures
forces every coefficient to vanish, and Theorem~\ref{thm:positive_defect}
then places every zero on the critical line.  Conversely, under \(\RH\),
all complete-zero weights equal one and the measures are identical.  Since
a locally finite atomic measure is determined by its jumps, equality of
jumps is equivalent to equality of measures.
\end{proof}

\begin{corollary}[Midpoint residual criterion]
\label{cor:midpoint_residual_criterion}
Suppose an explicit-formula reconstruction is proved to converge locally to
the complete-zero midpoint walk, and suppose its residual at every ordinate
\(\tau\) is shown to equal
\begin{equation}
\frac{m(\tau)}2e^{i\tau\log q}.
\label{eq:midpoint_residual_target}
\end{equation}
Then \(\RH\) follows.
\end{corollary}

\begin{proof}
The complete midpoint jump is
\(\frac12W_q(\tau)e^{i\tau\log q}\).  Equality with
\eqref{eq:midpoint_residual_target} gives
\(W_q(\tau)=m(\tau)\) for every \(\tau\).  Apply
Theorem~\ref{thm:route_fixed_q_atomic}.
\end{proof}

\subsection{Route II: local distributional convergence}
\label{subsec:route_distributional}

\begin{theorem}[Test-function criterion]
\label{thm:test_function_criterion}
Let \(\eta_j\) be a sequence of locally finite complex measures such that
\begin{equation}
\eta_j\longrightarrow\nu_q-\mu_q
\quad\text{in }\mathcal D'(\mathbb R).
\label{eq:test_measure_identification}
\end{equation}
If, for every \(\varphi\in C_c^\infty(\mathbb R)\),
\begin{equation}
\langle\eta_j,\varphi\rangle\longrightarrow0,
\label{eq:test_measure_zero}
\end{equation}
then \(\RH\) holds.
\end{theorem}

\begin{proof}
For every test function \(\varphi\), convergence
\eqref{eq:test_measure_identification} gives
\begin{equation}
\langle\eta_j,\varphi\rangle
\longrightarrow
\langle\nu_q-\mu_q,\varphi\rangle.
\label{eq:test_measure_limit_first}
\end{equation}
The second hypothesis says that the same left-hand side tends to zero.
Uniqueness of limits in \(\mathbb C\) yields
\begin{equation}
\langle\nu_q-\mu_q,\varphi\rangle=0
\label{eq:test_measure_limit_zero}
\end{equation}
for every \(\varphi\in C_c^\infty\).  Therefore the measure difference is
the zero distribution, hence the zero locally finite measure.  The
fixed-\(q\) atomic criterion then implies \(\RH\).
\end{proof}

\begin{remark}
\label{rem:route_distributional_missing}
For the prime-power construction, the missing theorem is a joint
\((H,Q,\eta)\) limit that identifies the derivative of the regularized
prime side with \(\nu_q-\mu_q\), while the corresponding residual tends
locally to zero.  A normalized average over \([0,H]\) is not sufficient,
because it can dilute fixed atoms.
\end{remark}

\subsection{Route III: small-\(q\) curvature at every ordinate}
\label{subsec:route_small_q}

\begin{theorem}[Quadratic curvature criterion]
\label{thm:small_q_curvature}
Put \(q=e^y\).  At each ordinate \(\tau\),
\begin{equation}
D_{e^y}(\tau)
=
\frac{y^2}{2}
\sum_{\substack{\rho=\beta+i\tau\\\beta>1/2}}
\left(\beta-\frac12\right)^2
+
O_\tau(y^4)
\qquad(y\to0).
\label{eq:small_q_curvature_expansion}
\end{equation}
Consequently,
\begin{equation}
\RH
\quad\Longleftrightarrow\quad
D_{e^y}(\tau)=o(y^2)
\text{ for every }\tau.
\label{eq:small_q_curvature_criterion}
\end{equation}
\end{theorem}

\begin{proof}
At a fixed ordinate, enumerate the positive horizontal displacements as
\(\delta_1,\ldots,\delta_m\), repeated with multiplicity.  Taylor's theorem
at zero gives, uniformly over this finite set,
\begin{equation}
\cosh(\delta_jy)-1
=
\frac{\delta_j^2y^2}{2}
+
\frac{\delta_j^4y^4}{24}
+
O_\tau(y^6).
\label{eq:small_q_taylor_detailed}
\end{equation}
Summing proves the expansion.  If the total is \(o(y^2)\), division by
\(y^2\) and passage to the limit gives
\begin{equation}
\frac12\sum_{j=1}^m\delta_j^2=0.
\label{eq:small_q_sum_squares_zero}
\end{equation}
Every term is nonnegative, so every displacement is zero.  This holds at
every ordinate and yields \(\RH\).  Under \(\RH\), the defect is
identically zero, proving the converse.
\end{proof}

\begin{corollary}[Second-derivative formulation]
\label{cor:small_q_second_derivative}
The horizontal second moment at ordinate \(\tau\) is recovered by
\begin{equation}
\left.\frac{\dd^2}{\dd y^2}D_{e^y}(\tau)\right|_{y=0}
=
\sum_{\substack{\rho=\beta+i\tau\\\beta>1/2}}
\left(\beta-\frac12\right)^2.
\label{eq:small_q_second_derivative}
\end{equation}
\end{corollary}

\begin{proof}
Differentiate the finite sum
\begin{equation}
D_{e^y}(\tau)
=
\sum_j\left[\cosh(\delta_jy)-1\right]
\label{eq:small_q_finite_sum_for_derivative}
\end{equation}
twice.  The second derivative of the \(j\)-th term is
\(\delta_j^2\cosh(\delta_jy)\), which at \(y=0\) equals
\(\delta_j^2\).  Finite summation commutes with differentiation, proving
the formula.
\end{proof}

\subsection{Route IV: large-\(q\) amplification}
\label{subsec:route_large_q}

\begin{definition}
\label{def:max_horizontal_displacement}
At a positive ordinate \(\tau\), define
\begin{equation}
\delta_{\max}(\tau)
=
\max_{\rho=\beta+i\tau}
\left|\beta-\frac12\right|,
\label{eq:max_horizontal_displacement}
\end{equation}
with value zero when all zeros at the ordinate are on the critical line.
\end{definition}

\begin{theorem}[Recovery of the maximal horizontal displacement]
\label{thm:large_q_growth_exponent}
For every positive ordinate \(\tau\),
\begin{equation}
\boxed{
\lim_{q\to\infty}
\frac{\log\left(1+D_q(\tau)\right)}{\log q}
=
\delta_{\max}(\tau).
}
\label{eq:large_q_growth_exponent}
\end{equation}
\end{theorem}

\begin{proof}
If \(\delta_{\max}=0\), then \(D_q(\tau)=0\) identically and the limit is
zero.  Otherwise, let \(m_{\max}\geq1\) denote the number of paired zeros
with displacement \(\delta_{\max}\).  Since
\begin{equation}
\cosh(\delta\log q)-1
=
\frac12q^\delta+\frac12q^{-\delta}-1,
\label{eq:large_q_each_term}
\end{equation}
we have the lower bound
\begin{equation}
D_q(\tau)
\geq
m_{\max}\left(\frac12q^{\delta_{\max}}-1\right).
\label{eq:large_q_lower}
\end{equation}
If \(M(\tau)\) is the finite number of off-line pairs at the ordinate,
every summand is at most \(q^{\delta_{\max}}\) for sufficiently large
\(q\), so
\begin{equation}
D_q(\tau)
\leq
M(\tau)q^{\delta_{\max}}.
\label{eq:large_q_upper}
\end{equation}
Taking logarithms, dividing by \(\log q\), and applying the squeeze theorem
gives the result.
\end{proof}

\begin{corollary}[Subpolynomial-growth route]
\label{cor:subpolynomial_growth_route}
If one can prove from the prime side that, for every ordinate \(\tau\),
\begin{equation}
D_q(\tau)=q^{o(1)}
\qquad(q\to\infty),
\label{eq:subpolynomial_defect}
\end{equation}
then \(\RH\) follows.
\end{corollary}

\begin{proof}
The growth exponent in Theorem~\ref{thm:large_q_growth_exponent} must then
be zero at every ordinate.
\end{proof}

\subsection{Route V: reconstruction from a family of frequencies}
\label{subsec:route_family_q}

\begin{theorem}[Finite exponential-polynomial reconstruction]
\label{thm:q_family_reconstruction}
At a fixed ordinate \(\tau\), the function
\begin{equation}
\mathcal W_\tau(y)
=
W_{e^y}(\tau)
=
\sum_{\rho:\Im\rho=\tau}
e^{(\beta-1/2)y}
\label{eq:q_family_exponential_polynomial}
\end{equation}
is an entire exponential polynomial.  Knowledge of
\(\mathcal W_\tau(y)\) on any real set having an accumulation point
uniquely determines the multiset of horizontal displacements
\(\{\beta-1/2\}\).  In particular,
\begin{equation}
\mathcal W_\tau(y)\equiv m(\tau)
\quad\Longleftrightarrow\quad
\text{all zeros at }\tau\text{ lie on }\Re s=\frac12.
\label{eq:q_family_constant_criterion}
\end{equation}
\end{theorem}

\begin{proof}
The sum is finite, so \(\mathcal W_\tau\) is entire.  If two finite
exponential polynomials agree on a set with an accumulation point, the
identity theorem makes them equal everywhere.  Their derivatives at zero
are the power sums
\begin{equation}
\mathcal W_\tau^{(k)}(0)
=
\sum_{\rho:\Im\rho=\tau}
\left(\beta-\frac12\right)^k.
\label{eq:q_family_power_sums}
\end{equation}
The finite multiset of exponents is determined by these power sums through
Newton identities, after combining repeated values with multiplicity.  If
the function is the constant \(m(\tau)\), its second derivative is zero;
by symmetry and nonnegativity this forces every squared displacement to be
zero.  The converse is immediate.
\end{proof}

\subsection{Route VI: projected energy exhaustion plus local atomic control}
\label{subsec:route_energy_local}

\begin{theorem}[A sufficient two-part energy route]
\label{thm:energy_local_route}
Let \(F_j\) be regularized residuals and suppose:
\begin{enumerate}[label=\textup{(\roman*)}]
\item for every compact interval \(I\),
\begin{equation}
\int_I\left|F_j(T)-\frac12\right|^2\,\dd T
\longrightarrow0;
\label{eq:energy_route_local_L2}
\end{equation}
\item the distributional derivatives satisfy
\begin{equation}
\frac{\dd}{\dd T}
\left[e^{i\omega T}\left(F_j(T)-\frac12\right)\right]
\longrightarrow
\nu_q-\mu_q.
\label{eq:energy_route_derivative}
\end{equation}
\end{enumerate}
Then \(\RH\) holds.
\end{theorem}

\begin{proof}
Condition (i) is local \(L^2\)-convergence of \(F_j-1/2\) to zero.
Multiplication by the smooth unit-modulus factor \(e^{i\omega T}\) preserves
that convergence.  On compact sets, local \(L^2\) implies local \(L^1\) by
Cauchy--Schwarz.  Consequently, the distributional derivative of the
product tends to zero when tested against a compactly supported smooth
function.  Condition (ii) identifies that same derivative limit with
\(\nu_q-\mu_q\).  Hence the measure difference is zero, and the positive
atomic criterion gives \(\RH\).
\end{proof}

\begin{remark}
\label{rem:global_variance_insufficient_RH_route}
The measured quantity \(N(H)^{-1}\sum_{\gamma\leq H}|F(\gamma)|^2\) is a
global normalized energy.  Even if it tends to zero, it does not imply
\eqref{eq:energy_route_local_L2}.  The route requires either local windows,
compactly supported test functions, or an independent argument excluding a
finite or sparse atomic defect.
\end{remark}

\subsection{Route VII: a one-sided positive test family}
\label{subsec:route_positive_tests}

\begin{theorem}[Exhausting positive tests]
\label{thm:exhausting_positive_tests}
Let
\begin{equation}
\mathcal D_q
=
\sum_{\tau>0}D_q(\tau)\delta_\tau.
\label{eq:positive_defect_measure_route}
\end{equation}
Suppose there is a sequence of nonnegative test functions
\(\varphi_j\in C_c^\infty(0,\infty)\) such that for every \(T>0\),
\begin{equation}
\liminf_{j\to\infty}
\inf_{0<t\leq T}\varphi_j(t)>0,
\label{eq:positive_tests_cover}
\end{equation}
and
\begin{equation}
\langle\mathcal D_q,\varphi_j\rangle\longrightarrow0.
\label{eq:positive_tests_vanish}
\end{equation}
Then \(\RH\) holds.
\end{theorem}

\begin{proof}
Fix an ordinate \(\tau\).  For all sufficiently large \(j\),
\(\varphi_j(\tau)>0\).  Positivity of all masses gives
\begin{equation}
0\leq
D_q(\tau)\varphi_j(\tau)
\leq
\langle\mathcal D_q,\varphi_j\rangle.
\label{eq:positive_tests_atom_bound}
\end{equation}
By \eqref{eq:positive_tests_cover}, the values
\(\varphi_j(\tau)\) are bounded below by a positive constant for all
sufficiently large \(j\).  Since the right-hand side tends to zero, the
fixed coefficient \(D_q(\tau)\) must vanish.  Since \(\tau\) was arbitrary, the positive defect
measure is zero and \(\RH\) follows.
\end{proof}

\begin{remark}
\label{rem:routes_summary_missing_estimates}
The exact algebra does not choose among these routes.  The analytic work
still required is one of the following: local reconstruction of the
ordinate-only measure, a sufficiently strong uniform bound in \(q\), or a
local projected-energy theorem accompanied by derivative identification.
The finite numerical decay of a global variance is evidence for none of
these implications by itself.
\end{remark}


\section{Open problems and falsifiable next computations}
\label{sec:open_problems}

\begin{question}[Increment audit]
\label{question:increment_audit}
For every prime power \(n\) in a growing range, compute the exact ratio
\begin{equation}
\mathcal I_H(n)
=
\frac{\Re M_H(n)}{\pi v_n}
\label{eq:increment_audit_ratio}
\end{equation}
and compare it with
\begin{equation}
\frac{\log n}{\log(H/2\pi)}.
\label{eq:increment_audit_prediction}
\end{equation}
This tests the local Landau mechanism without fitting a global law.
\end{question}

\begin{question}[Height--cutoff matrix]
\label{question:height_cutoff_matrix}
Repeat the complete computation for a two-dimensional grid of
\((H,Q)\), including regimes
\begin{equation}
Q=H^\theta
\label{eq:height_cutoff_regime}
\end{equation}
with several fixed \(\theta\) and a near-height regime
\(Q=H^{1-o(1)}\).  Record raw and joint coefficients separately.
\end{question}

\begin{question}[Numerical audit of the exact variance components]
\label{question:projected_tower_blocks}
For every cutoff and height, compute separately
\begin{equation}
V_Y(H),
\quad
-2\Re\sum_p\mathbf c_p^\ast\mathbf d_{p,H},
\quad
\sum_p\mathbf c_p^\ast\widetilde K_{p,p;H}\mathbf c_p,
\quad
\sum_{p\neq r}\mathbf c_p^\ast\widetilde K_{p,r;H}\mathbf c_r,
\label{eq:numerical_variance_components}
\end{equation}
and verify numerically that their sum equals the directly fitted residual
variance in \eqref{eq:numerical_variance_exact_definition}.
\end{question}

\begin{question}[Local test-function convergence]
\label{question:local_test_functions}
For compactly supported smooth \(\varphi\), determine whether the
regularized residual satisfies
\begin{equation}
\int
\left(F_{Q,H}(T)-\frac12\right)\varphi(T)\dd T
\longrightarrow0
\label{eq:test_function_target}
\end{equation}
under a joint limiting regime, while the prime-side derivative is
identified with the complete spectral measure.
\end{question}


\part{Numerical support}

\section{Computational object and exact program identities}
\label{sec:numerical_program}

\subsection{Input data and sampled height}
\label{subsec:numerical_input_data}

\begin{experiment}[Zero data]
\label{exp:numerical_zero_data}
The full-data run used $2{,}001{,}052$ positive ordinates from
\texttt{zeross.bin}, with
\begin{equation}
\gamma_1=14.134725142,
\qquad
\gamma_{2{,}001{,}052}=1{,}132{,}490.6587144111.
\label{eq:numerical_zero_endpoints_detailed}
\end{equation}
The phase parameter was
\begin{equation}
\omega=1,
\qquad q=e.
\label{eq:numerical_omega}
\end{equation}
All ordinates were used; the stride was one.
\end{experiment}

\subsection{Exact definition of the program}
\label{subsec:exact_program_definition}

\begin{definition}[Program variables]
\label{def:numerical_program_variables}
At each sampled ordinate, the code forms
\begin{equation}
Y(\gamma_k)
=
S_q(\gamma_k)-A_{\mathrm{smooth}}(\gamma_k),
\label{eq:numerical_Y_program}
\end{equation}
and
\begin{equation}
R_Q(\gamma_k)
=
Y(\gamma_k)-P_Q^-(\gamma_k)-P_Q^+(\gamma_k),
\label{eq:numerical_RQ_program}
\end{equation}
where
\begin{align}
P_Q^-(\gamma)
&=
\sum_{\substack{n=p^r\\n\leq Q}}
 c_n^-e^{i(\omega-\ell_n)\gamma},
\label{eq:numerical_PQ_minus}
\\
P_Q^+(\gamma)
&=
\sum_{\substack{n=p^r\\n\leq Q}}
 c_n^+e^{i(\omega+\ell_n)\gamma}.
\label{eq:numerical_PQ_plus}
\end{align}
The raw carrier statistic is
\begin{equation}
\widehat c_{Q,H}
=
\frac1{N(H)}
\sum_{\gamma_k\leq H}
e^{-i\omega\gamma_k}R_Q(\gamma_k).
\label{eq:numerical_raw_carrier}
\end{equation}
\end{definition}

\begin{proposition}[Identity evaluated by the raw program]
\label{prop:numerical_raw_identity}
The program variables satisfy exactly
\begin{equation}
\widehat c_{Q,H}
=
\widehat c_{Y,H}
-
\sum_{\substack{n=p^r\\n\leq Q}}
\left(c_n^-\overline{M_H(n)}+c_n^+M_H(n)\right).
\label{eq:numerical_raw_identity}
\end{equation}
Its imaginary part is
\begin{equation}
\Im\widehat c_{Q,H}
=
\Im\widehat c_{Y,H}
-
\sum_{\substack{n=p^r\\n\leq Q}}
u_n\Re M_H(n).
\label{eq:numerical_imag_identity}
\end{equation}
\end{proposition}

\begin{proof}
The arrays used by the program are finite, so the abstract identity in
Theorem~\ref{thm:abstract_carrier_formula} applies without any limiting
operation.  The code generates exactly the prime powers \(n\leq Q\), uses
the coefficients in Definition~\ref{def:prime_power_coefficients}, and
forms the sample kernel by averaging \(n^{i\gamma_k}\).  Substitution of
these arrays in the theorem gives the program identity entry by entry.
\end{proof}

\subsection{Joint constant--carrier fit used for the reported values}
\label{subsec:numerical_joint_fit}

\begin{definition}
\label{def:numerical_joint_fit}
For every cutoff, the code fits
\begin{equation}
R_Q(\gamma_k)
=
\alpha_{Q,H}+c_{Q,H}e^{i\omega\gamma_k}+E_{Q,H}(\gamma_k)
\label{eq:numerical_joint_model}
\end{equation}
by least squares in the sampled inner product
\eqref{eq:sample_inner_product}.  The reported variance is
\begin{equation}
V(Q,H)
=
\frac1{N(H)}
\sum_{\gamma_k\leq H}|E_{Q,H}(\gamma_k)|^2.
\label{eq:numerical_reported_variance}
\end{equation}
\end{definition}

\begin{proposition}[Closed form of the reported carrier]
\label{prop:numerical_joint_carrier}
Let
\begin{equation}
\mu_H=M_H(e^\omega).
\label{eq:numerical_muH}
\end{equation}
Then
\begin{equation}
\boxed{
c_{Q,H}
=
\frac{
\langle e^{i\omega\gamma},R_Q\rangle_H
-
\overline{\mu_H}\langle1,R_Q\rangle_H
}{1-|\mu_H|^2}.
}
\label{eq:numerical_joint_carrier}
\end{equation}
The fitted constant is
\begin{equation}
\alpha_{Q,H}
=
\frac{
\langle1,R_Q\rangle_H
-
\mu_H\langle e^{i\omega\gamma},R_Q\rangle_H
}{1-|\mu_H|^2}.
\label{eq:numerical_joint_constant}
\end{equation}
\end{proposition}

\begin{proof}
The fit minimizes
\begin{equation}
\sum_{k=1}^{N(H)}|R_k-a-cb_k|^2
\label{eq:numerical_joint_objective_proof}
\end{equation}
over \(a,c\in\mathbb C\).  Differentiating with respect to the conjugate
parameters, or equivalently imposing orthogonality of the residual to the
basis vectors \(1\) and \(b\), gives the two normal equations.  Their Gram
matrix is
\(\bigl(\begin{smallmatrix}1&\mu_H\\\overline\mu_H&1\end{smallmatrix}\bigr)\).
Its inverse is explicit, and multiplication by the right-hand side gives
the displayed formulas for the constant and carrier coefficients.
\end{proof}

\subsection{Prime-power generation and exact coefficients}
\label{subsec:numerical_prime_powers}

\begin{experiment}[Prime powers]
\label{exp:numerical_prime_power_counts}
The cutoffs and counts were
\begin{center}
\begin{tabular}{rr}
\toprule
$Q$ & number of prime powers $n=p^r\leq Q$\\
\midrule
0 & 0\\
3 & 2\\
10 & 7\\
100 & 35\\
1000 & 193\\
10000 & 1280\\
\bottomrule
\end{tabular}
\end{center}
Every coefficient was generated from
\eqref{eq:cminus_definition}--\eqref{eq:cplus_definition}; no frequency or
amplitude was fitted individually.
\end{experiment}

\section{Radial geometry and low prime-power resonances}
\label{sec:numerical_radial_geometry}

\subsection{Direct radial audit}
\label{subsec:direct_radial_audit}

\begin{experiment}[Original walk and theoretical subtractions]
\label{exp:direct_radial_audit_detailed}
After subtraction of the empirical complex center, the radial statistics
were
\begin{center}
\begin{tabular}{lccc}
\toprule
case & mean radius & radial standard deviation & smoothed radial modes\\
\midrule
original & $1.9203104476$ & $0.7265477726$ & $1.35731059,\ 2.49789956$\\
after $n=3$ & $1.8079247921$ & $0.3083120307$ & $1.68017926$\\
after $n=2,3$ & $1.7884459085$ & $0.2214048337$ & $1.81300495$\\
\bottomrule
\end{tabular}
\end{center}
The theoretical $n=3$ coefficients were
\begin{equation}
c_3^-=-1.02370050729i,
\qquad
c_3^+=0.0481029537851i.
\label{eq:numerical_c3_coefficients_detailed}
\end{equation}
Subtracting $n=3$ reduced the radial variance by approximately
$81.99\%$; subtracting $n=2,3$ reduced it by approximately $90.71\%$.
The finite conclusion is that the observed global double radial mode is
dominated by the theoretical $n=3$ oscillator.
\end{experiment}

\subsection{Geometric interpretation}
\label{subsec:numerical_geometric_interpretation}

\begin{experiment}[Anisotropy]
\label{exp:numerical_anisotropy}
At $Q=10000$, the centered residual anisotropy was
\begin{equation}
0.981212.
\label{eq:numerical_anisotropy_value}
\end{equation}
This is consistent with the exact single-mode limit
\eqref{eq:single_mode_anisotropy}: after low frequencies are removed,
large-$n$ modes are strongly elongated in one direction.
\end{experiment}

\section{Cutoff sweep and carrier stabilization}
\label{sec:numerical_cutoff_sweep}

\subsection{Global cutoff table}
\label{subsec:numerical_global_cutoff_table}

\begin{experiment}[Full-data cutoff sweep]
\label{exp:full_data_sweep}
The jointly fitted carrier and residual energy were
\begin{center}
\begin{tabular}{rrrrrr}
\toprule
$Q$ & terms & $\Re c_{Q,H}$ & $\Im c_{Q,H}$ & $V(Q,H)$ & $\sqrt{V(Q,H)}$\\
\midrule
0 & 0 & 0.4999989123 & 0.1381673175 & 1.2420277147 & 1.1144630\\
3 & 2 & 0.4999998679 & 0.0956865241 & 0.1339324657 & 0.3659678\\
10 & 7 & 0.4999999862 & 0.0755224108 & 0.0596099186 & 0.2441514\\
100 & 35 & 0.5000000416 & 0.0545919879 & 0.0220396104 & 0.1484574\\
1000 & 193 & 0.5000000755 & 0.0426449212 & 0.0112822119 & 0.1062178\\
10000 & 1280 & 0.4999998710 & 0.0343125459 & 0.0082559398 & 0.0908622\\
\bottomrule
\end{tabular}
\end{center}
At the largest cutoff the finite variance reduction relative to $Q=0$ was
\begin{equation}
1-\frac{0.0082559398}{1.2420277147}
=0.99335\ldots.
\label{eq:numerical_variance_reduction_detailed}
\end{equation}
\end{experiment}

\subsection{Normalized $1/\log Q$ comparisons}
\label{subsec:numerical_normalized_scaling}

\begin{experiment}[Two normalized quantities]
\label{exp:numerical_normalized_quantities}
For $Q\geq3$, the finite normalized values were
\begin{center}
\begin{tabular}{rcc}
\toprule
$Q$ & $\pi\log Q\,\Im c_{Q,H}$ &
$\dfrac{4\pi^2}{3}\log Q\,V(Q,H)$\\
\midrule
3 & 0.33025 & 1.93628\\
10 & 0.54631 & 1.80623\\
100 & 0.78981 & 1.33564\\
1000 & 0.92545 & 1.02558\\
10000 & 0.99284 & 1.00065\\
\bottomrule
\end{tabular}
\end{center}
At $Q=10000$,
\begin{align}
\frac1{\pi\log Q}
&=0.0345601\ldots,
\label{eq:numerical_imag_prediction}
\\
\frac3{4\pi^2\log Q}
&=0.00825061\ldots.
\label{eq:numerical_variance_prediction}
\end{align}
The corresponding measured values were
$0.0343125459$ and $0.0082559398$.  These are finite coincidences to be
interpreted through the exact defect identities, not proofs of asymptotic
laws.
\end{experiment}

\section{Height dependence}
\label{sec:numerical_height_dependence}

\subsection{Twenty-four block audit}
\label{subsec:numerical_block_audit}

\begin{experiment}[Block ranges]
\label{exp:numerical_block_ranges}
For $Q=10000$ and $24$ consecutive height blocks,
\begin{align}
\Re\operatorname{mean}_{\mathrm{block}}F_Q
&\in[0.499994991361,0.500003656384],
\label{eq:numerical_block_real_range_detailed}
\\
\Im\operatorname{mean}_{\mathrm{block}}F_Q
&\in[0.033825249952,0.035105891286],
\label{eq:numerical_block_imag_range_detailed}
\\
V_{\mathrm{block}}
&\in[0.0054744557,0.0089851833].
\label{eq:numerical_block_var_range_detailed}
\end{align}
The standard deviation of the block means was
\begin{equation}
0.0003812589.
\label{eq:numerical_block_mean_std}
\end{equation}
The carrier is stable across height at this scale, while the block variance
has a visible drift.  Consequently the global factor $3/4$ cannot yet be
regarded as height-uniform.
\end{experiment}

\subsection{Termwise increment scale}
\label{subsec:numerical_increment_scale}

\begin{experiment}[Present range of the Landau ratio]
\label{exp:numerical_landau_ratio}
For
\begin{equation}
H=1.13249\times10^6,
\qquad Q=10^4,
\label{eq:numerical_H_Q_pair}
\end{equation}
one has approximately
\begin{equation}
\frac{\log Q}{\log(H/2\pi)}\approx0.76.
\label{eq:numerical_log_ratio}
\end{equation}
The global ratio
\begin{equation}
\frac{\Im c_{Q,H}}{\pi\Delta_1(Q)}\approx0.986
\label{eq:numerical_global_ratio_again}
\end{equation}
is therefore not explained by independent termwise transport at the
leading Landau scale.  The exact increment statistic
\eqref{eq:increment_ratio} is the direct falsifiable test.
\end{experiment}

\section{Branch decomposition and omitted frequencies}
\label{sec:numerical_branches_omitted}

\subsection{Branch energies}
\label{subsec:numerical_branch_energies}

\begin{experiment}[Minus and plus branches]
\label{exp:numerical_branch_energy_detailed}
At $Q=10000$,
\begin{align}
\|P_Q^-\|^2
&=1.2433344613,
\label{eq:numerical_minus_energy_detailed}
\\
\|P_Q^+\|^2
&=0.0309962015,
\label{eq:numerical_plus_energy_detailed}
\\
2\Re\langle P_Q^-,P_Q^+\rangle
&=-0.0005963448.
\label{eq:numerical_cross_energy_detailed}
\end{align}
The minus branch dominates, and direct plus--minus energy cancellation is
small.
\end{experiment}

\subsection{Omitted prime-power shell}
\label{subsec:numerical_omitted_shell}

\begin{experiment}[Exact-frequency omitted fit]
\label{exp:numerical_omitted_fit_detailed}
For the $56$ omitted prime powers immediately above $Q=10000$, the median
amplitude-to-null ratios were
\begin{equation}
3.35349
\quad\text{and}\quad
2.67559
\label{eq:numerical_omitted_ratios_detailed}
\end{equation}
for the minus and plus branches.  The absolute correlation between the
jointly fitted coefficient vector and the theoretical coefficient vector was
\begin{equation}
0.916474.
\label{eq:numerical_omitted_correlation_detailed}
\end{equation}
The narrow shell explained
\begin{equation}
0.0144446
\label{eq:numerical_omitted_energy_fraction_detailed}
\end{equation}
of the remaining residual energy.  Thus the next prime-power frequencies
are present, but this narrow shell does not account for the full residual.
\end{experiment}

\section{Signed invariant and finite-area comparisons}
\label{sec:numerical_signed_invariant}

\subsection{Direct evaluation of $\Delta_1(Q)$}
\label{subsec:numerical_Delta_evaluation}

\begin{experiment}[Common invariant]
\label{exp:numerical_common_invariant}
At $Q=10000$,
\begin{equation}
\Delta_1(10000)\approx0.01108.
\label{eq:numerical_Delta_detailed}
\end{equation}
Hence
\begin{equation}
\pi\Delta_1(10000)\approx0.03481,
\qquad
\frac34\Delta_1(10000)\approx0.00831.
\label{eq:numerical_Delta_multiples_detailed}
\end{equation}
The same exact signed invariant therefore reproduces both observed scales
to finite accuracy.  The analytic claims that the centroid and projected
variance equal these multiples remain Conjectures~\ref{conj:collective_carrier_law}
and \ref{conj:projected_variance_law}.
\end{experiment}

\subsection{Coefficient prefix sums}
\label{subsec:numerical_coefficient_prefixes}

\begin{experiment}[Prefix sums]
\label{exp:numerical_prefix_sums}
The coefficient combinations
\begin{equation}
A_Q=\sum_{n\leq Q}(c_n^-+c_n^+),
\qquad
B_Q=\sum_{n\leq Q}(c_n^--c_n^+)
\label{eq:numerical_AQ_BQ}
\end{equation}
were, at $Q=10000$,
\begin{equation}
A_Q=-2.5552137617i,
\qquad
B_Q=-10.6932626383i.
\label{eq:numerical_AQ_BQ_values_detailed}
\end{equation}
They do not account for the real carrier value $1/2$.
\end{experiment}

\section{Finite-part computations for $\Phi_N$}
\label{sec:numerical_finite_parts}

\subsection{Direct algebraic verification}
\label{subsec:numerical_finite_part_verification}

\begin{experiment}[Finite-part identity]
\label{exp:numerical_finite_part_identity}
A numerical comparison between the direct tail expression and the value
reconstructed from Proposition~\ref{prop:finite_part_Phi} gave
\begin{equation}
|R_{\mathrm{from\,FP}}-R_{\mathrm{direct}}|
=4.043\times10^{-10}.
\label{eq:numerical_finite_part_error_detailed}
\end{equation}
This verifies the implemented identity at the working precision.
\end{experiment}

\subsection{Scale scan}
\label{subsec:numerical_finite_part_scale_scan}

\begin{experiment}[Normalized finite-part differences]
\label{exp:numerical_finite_part_scale_scan_detailed}
For
\begin{equation}
N=2000,4000,8000,16000,32000,64000,
\label{eq:numerical_finite_part_N_values}
\end{equation}
the successive magnitudes
\begin{equation}
|\mathcal F_{2N}(\rho)-\mathcal F_N(\rho)|
\label{eq:numerical_finite_part_difference_definition}
\end{equation}
were approximately
\begin{equation}
1.12,
\quad0.75,
\quad0.80,
\quad0.75,
\quad0.44.
\label{eq:numerical_finite_part_difference_values}
\end{equation}
No monotone cancellation law is established by these six scales.
\end{experiment}

\section{Numerical tests dictated by the analytic identities}
\label{sec:numerical_next_tests}

\subsection{Prime-power increment audit}
\label{subsec:numerical_increment_audit}

\begin{experiment}[Required output]
\label{exp:required_increment_output}
For every prime power $n$ in the selected range, compute and store
\begin{equation}
\mathcal I_H(n)
=
\frac{\Re M_H(n)}{\pi v_n},
\label{eq:numerical_increment_output}
\end{equation}
together with
\begin{equation}
\theta_H(n)
=
\frac{\log n}{\log(H/2\pi)}.
\label{eq:numerical_theta_output}
\end{equation}
The comparison is pointwise and requires no global fitting.
\end{experiment}

\subsection{Two-parameter height--cutoff scan}
\label{subsec:numerical_two_parameter_scan}

\begin{experiment}[Required grid]
\label{exp:required_height_cutoff_grid}
For several heights $H_j$ and exponents $\theta$, set
\begin{equation}
Q_{j,\theta}=H_j^\theta
\label{eq:numerical_Q_theta_grid}
\end{equation}
and record separately:
\begin{equation}
\widehat c_{Q,H},
\quad
c_{Q,H},
\quad
V(Q,H),
\quad
\Delta_\omega(Q),
\quad
\sum_p a_p^\ast\widetilde K_{p,p;H}a_p,
\quad
\sum_{p\neq q}a_p^\ast\widetilde K_{p,q;H}a_q.
\label{eq:numerical_required_grid_outputs}
\end{equation}
This directly separates the raw carrier, the joint fitted carrier, the
projected background energy, the mixed background--prime term, the tower
blocks, and the cross-tower interactions.
\end{experiment}



\section{Complete MATLAB implementation}
\label{sec:numerical_complete_matlab}

\subsection{Self-contained implementation}
\label{subsec:numerical_self_contained_implementation}

\begin{remark}
\label{rem:complete_matlab_scope}
The following program is the complete implementation used for the
prime-power audit.  It reads only the original binary ordinate file,
generates all prime powers internally, retains all figures, performs the
joint constant--carrier fit, and exports every table used in the preceding
numerical sections.
\end{remark}

\lstinputlisting[
  style=matlabclean,
  language=Matlab,
  caption={Complete MATLAB implementation of the prime-power phase audit.},
  label={lst:complete_matlab}
]{ZETA_PRIME_POWER_ANALYTIC_AUDIT_RIGOROUS.m}

\section{Complete console output of the full-data run}
\label{sec:numerical_complete_output}

\subsection{Coefficient identities, radial audit, branches, towers, blocks,
and omitted frequencies}
\label{subsec:numerical_complete_output_tables}

\begin{remark}
\label{rem:complete_output_scope}
The console transcript below is included as numerical support.  The analytic
claims of the manuscript are established in the theorem--proof sections and
do not depend on numerical rounding in this transcript.
\end{remark}

\lstinputlisting[
  style=consoleclean,
  caption={Complete console output for the corrected full-data run.},
  label={lst:complete_output}
]{PRIME_POWER_FULL_DATA_RUN.txt}

\section{Conclusion}
\label{sec:conclusion}

The finite Möbius interpolation is completely oriented by
\begin{equation}
\Phi_N(s)
=
M(N)+h_N\frac{R_N(s)}{P_N(s)}.
\label{eq:conclusion_Phi}
\end{equation}
Its value at zero, its integer limit, its exponential rate, its rounding
criterion, and its recovery of \(a_N\) are mutually consistent.  The
absolutely convergent half-plane does not select \(M(N)\) among feasible
terminal values.

For the zeta phase walk, the following components are exact:

\begin{enumerate}[label=\textup{(\roman*)}]
\item the endpoint conversion and triangular half identity;
\item the positive atomic horizontal defect \(D_q(\tau)\);
\item the two-branch coefficient algebra;
\item the convergent signed invariant
\(\Delta_\omega(Q)\sim\omega/(\pi^2\log Q)\);
\item the raw zero-kernel identity for the carrier coefficient;
\item the exact global and local defects
\eqref{eq:global_defect_identity} and
\eqref{eq:local_defect_identity};
\item the joint constant--carrier projection formula;
\item the exact four-term decomposition of the projected sampled
residual variance in \eqref{eq:exact_residual_variance} and
\eqref{eq:tower_variance_background_mixed}--
\eqref{eq:tower_variance_decomposition}.
\end{enumerate}

The finite computation shows strong prime-power structure,
including the disappearance of the radial double mode after subtraction of
\(n=3\), a \(99.335\%\) finite variance reduction through \(Q=10000\),
stable real carrier coefficient near \(1/2\), and detectable omitted
prime-power frequencies.  None of these finite observations proves
uniform energy exhaustion.

The possible routes to the Riemann Hypothesis are now explicit: equality of
atomic jumps for one fixed \(q>1\), local distributional convergence,
small-\(q\) quadratic curvature, large-\(q\) subpolynomial growth, or local
projected-energy exhaustion accompanied by derivative identification.  The
main unresolved statements are not ordinary Parseval estimates.  They are
uniform estimates for the projected zero-sampled kernel, together with a
local explicit-formula limit strong enough to identify the positive atomic
measure
\begin{equation}
\sum_\tau D_q(\tau)\delta_\tau.
\label{eq:conclusion_positive_measure}
\end{equation}
The termwise area mechanism fails away from the near-height regime.
Any stable global relation between the centroid and
\(\Delta_\omega(Q)\) must therefore be collective.

\begin{thebibliography}{99}

\bibitem{Corless1996}
R.~M. Corless, G.~H. Gonnet, D.~E.~G. Hare, D.~J. Jeffrey, and D.~E. Knuth,
\emph{On the Lambert \(W\) function},
Adv. Comput. Math. \textbf{5} (1996), 329--359.
doi:10.1007/BF02124750.



\bibitem{Davenport2000}
H.~Davenport,
\emph{Multiplicative Number Theory},
3rd ed., revised by H.~L. Montgomery,
Graduate Texts in Mathematics 74, Springer, New York, 2000.
doi:10.1007/978-1-4757-5927-3.

\bibitem{IwaniecKowalski2004}
H.~Iwaniec and E.~Kowalski,
\emph{Analytic Number Theory},
American Mathematical Society Colloquium Publications 53,
Providence, 2004.

\bibitem{MontgomeryVaughan2007}
H.~L. Montgomery and R.~C. Vaughan,
\emph{Multiplicative Number Theory I: Classical Theory},
Cambridge Studies in Advanced Mathematics 97,
Cambridge University Press, Cambridge, 2007.
doi:10.1017/CBO9780511618314.


\bibitem{BombieriLagarias1999}
E.~Bombieri and J.~C. Lagarias,
\emph{Complements to Li's criterion for the Riemann hypothesis},
J. Number Theory \textbf{77} (1999), 274--287.
doi:10.1006/jnth.1999.2392.

\bibitem{Edwards1974}
H.~M. Edwards,
\emph{Riemann's Zeta Function},
Academic Press, New York, 1974.

\bibitem{FordSoundararajanZaharescu2009}
K.~Ford, K.~Soundararajan, and A.~Zaharescu,
\emph{On the distribution of imaginary parts of zeros of the Riemann
zeta function, II},
Math. Ann. \textbf{343} (2009), 487--505.
doi:10.1007/s00208-008-0280-x.
arXiv:0805.2745.

\bibitem{FordZaharescu2005}
K.~Ford and A.~Zaharescu,
\emph{On the distribution of imaginary parts of zeros of the Riemann
zeta function},
J. Reine Angew. Math. \textbf{579} (2005), 145--158.
doi:10.1515/crll.2005.2005.579.145.
arXiv:math/0405459.

\bibitem{Gonek1985}
S.~M. Gonek,
\emph{A formula of Landau and mean values of \(\zeta(s)\)},
in S.~W. Graham and J.~D. Vaaler (eds.),
\emph{Topics in Analytic Number Theory},
University of Texas Press, Austin, 1985, pp.~92--97.

\bibitem{Guinand1948}
A.~P. Guinand,
\emph{A summation formula in the theory of prime numbers},
Proc. London Math. Soc. (2) \textbf{50} (1948), 107--119.
doi:10.1112/plms/s2-50.2.107.

\bibitem{Montgomery1973}
H.~L. Montgomery,
\emph{The pair correlation of zeros of the zeta function},
in \emph{Analytic Number Theory},
Proc. Sympos. Pure Math. \textbf{24},
American Mathematical Society, Providence, 1973, pp.~181--193.

\bibitem{Titchmarsh1986}
E.~C. Titchmarsh,
\emph{The Theory of the Riemann Zeta-Function},
2nd ed., revised by D.~R. Heath-Brown,
Oxford University Press, Oxford, 1986.

\bibitem{Weil1952}
A.~Weil,
\emph{Sur les ``formules explicites'' de la théorie des nombres premiers},
Comm. Sém. Math. Univ. Lund, Tome Supplémentaire (1952), 252--265.

\end{thebibliography}

\end{document}
